% Les acides $\alpha$-aminés et la synthèse des peptides — Chimie Tle D — Cours signé OnBuch+
% Programme officiel MINESEC, Terminale C, D, E. Compiler avec : tectonic L04-acides-alpha-amines-et-peptides.tex
\newcommand{\DOCMATIERE}{Chimie}
\newcommand{\DOCNIVEAU}{Tle D}
\newcommand{\DOCMODULE}{Module 1 · Chimie organique}
\newcommand{\DOCLECON}{04}
\newcommand{\DOCTITRE}{Les acides $\alpha$-aminés et la synthèse des peptides}
% OnBuch+ — préambule « cours signés » ENRICHI (compilé avec tectonic / XeLaTeX)
% Système visuel « Pop » d'OnBuch+ : fond crème (jamais blanc), bordures encre
% épaisses, ombres « dures » décalées, titres Archivo Black en capitales.
% Version enrichie pour les cours de chimie : chimie (mhchem, chemfig),
% graphes (pgfplots), boîtes pédagogiques supplémentaires (définition,
% propriété, méthode, attention, le savais-tu, exercice résolu, exercice,
% TP, bilan), page de garde refaite et signature OnBuch+ ancrée en bas.
%
% Macros attendues dans main.tex AVANT \input{preamble} :
%   \DOCMATIERE  \DOCNIVEAU  \DOCMODULE  \DOCLECON (numéro)  \DOCTITRE
\documentclass[11pt]{article}

\usepackage{fontspec}
\usepackage[french]{babel}
\usepackage[a4paper,margin=2cm,top=3.4cm,bottom=2.8cm,headheight=1.4cm,footskip=1.3cm]{geometry}
\usepackage{amsmath,amssymb,amsfonts}
\usepackage[table]{xcolor} % [table] : fournit aussi \rowcolors (alternance de lignes)
\usepackage{pagecolor}
\usepackage{tikz}
\usetikzlibrary{arrows.meta,positioning,calc,shapes.geometric,shapes.misc,shapes.arrows,decorations.pathmorphing,decorations.markings,patterns,shadows,mindmap,trees,fit,backgrounds,matrix,intersections}
\usepackage{pgfplots}
\pgfplotsset{compat=1.18}
\usepackage[version=4]{mhchem}
\usepackage{chemfig}
\usepackage{enumitem}
\usepackage{fancyhdr}
% [most] charge toutes les bibliothèques tcolorbox (listings, minted, raster,
% documentation...) et gonfle inutilement la mémoire TeX sur un document long
% (~300 boîtes) : on ne charge que ce qui est réellement utilisé.
\usepackage[skins,breakable]{tcolorbox}
\usepackage{graphicx}
\usepackage{colortbl}
\usepackage{booktabs}
\usepackage{tabularx}
\usepackage{array}
\usepackage{multicol}
\usepackage{siunitx}
% parse-numbers=false : accepte aussi \SI{2,5 \times 10^{-3}}{...}
\sisetup{locale=FR,output-decimal-marker={,},group-minimum-digits=4,parse-numbers=false}
\usepackage{qrcode}
% \degree : commande fréquemment utilisée par les modèles pour un angle ou une
% température en dehors de \SI{}{} (non fournie par les paquets chargés).
\providecommand{\degree}{^{\circ}}
\usepackage{lastpage}
\usepackage{hyperref}
\hypersetup{hidelinks}

% ============================================================
% Palette « Pop » OnBuch+ — mêmes hex que la classe P (plus_ui.dart)
% ============================================================
\definecolor{poporange}{HTML}{F59321}
\definecolor{popdark}{HTML}{E07A0C}
\definecolor{popcream}{HTML}{FFF6E8}
\definecolor{popink}{HTML}{1C1714}
\definecolor{poppurple}{HTML}{7A5AE0}
\definecolor{popgreen}{HTML}{2E9E6A}
\definecolor{poppink}{HTML}{E0457B}
\definecolor{popblue}{HTML}{2F7DE1}
\definecolor{popgold}{HTML}{C98A12}
\definecolor{popmuted}{HTML}{8A7F76}
\definecolor{popline}{HTML}{EADFD2}
\definecolor{popgray}{HTML}{8A7F76}
\colorlet{popgrey}{popgray}
% Alias tolérants (le modèle nomme parfois les couleurs différemment)
\colorlet{popwhite}{white}
\colorlet{popblack}{popink}
\colorlet{popred}{poppink}
\colorlet{popyellow}{popgold}
% Teintes claires pour fonds de tableaux / zones
\colorlet{popblueL}{popblue!10!white}
\colorlet{poporangeL}{poporange!14!white}
\colorlet{popgreenL}{popgreen!12!white}
\colorlet{poppurpleL}{poppurple!12!white}
\colorlet{poppinkL}{poppink!12!white}
\colorlet{popgoldL}{popgold!14!white}

\pagecolor{popcream}

% ============================================================
% Typographie
% ============================================================
\defaultfontfeatures{Ligatures=TeX}
\setmainfont{PlusJakartaSans-Medium.ttf}[Path=../../../fonts/,BoldFont=PlusJakartaSans-Bold.ttf]
% Maths en sans-serif (Fira Math) avec chiffres et lettres droites de Plus
% Jakarta, pour que les formules se fondent dans le texte.
\usepackage[math-style=ISO,bold-style=ISO]{unicode-math}
\setmathfont{FiraMath-Regular.otf}[Path=../../../fonts/]
\setmathfont{PlusJakartaSans-Medium.ttf}[Path=../../../fonts/,range={up/num,up/{Latin,latin},\mathup}]
\usepackage{icomma} % virgule décimale sans espace en mode math
% Caractères absents de la police de texte (sinon carré vide) : rendus par
% leur équivalent mathématique, en texte comme en formule.
\usepackage{newunicodechar}
\newunicodechar{α}{\ensuremath{\alpha}}
\newunicodechar{Α}{\ensuremath{\alpha}}
\newunicodechar{β}{\ensuremath{\beta}}
\newunicodechar{δ}{\ensuremath{\delta}}
\newunicodechar{✓}{\popcheck{popgreen}}
\newfontfamily\popdisplay{ArchivoBlack-Regular.ttf}[Path=../../../fonts/]
\newfontfamily\poplogo{SpaceGrotesk-Bold.ttf}[Path=../../../fonts/]
\newfontfamily\popsemi{PlusJakartaSans-SemiBold.ttf}[Path=../../../fonts/]
\newfontfamily\popxbold{PlusJakartaSans-ExtraBold.ttf}[Path=../../../fonts/]
\newfontfamily\popxboldls{PlusJakartaSans-ExtraBold.ttf}[Path=../../../fonts/,LetterSpace=3.0]

\setlist[enumerate]{leftmargin=1.7em,itemsep=5pt,topsep=4pt}
\setlist[itemize]{leftmargin=1.7em,itemsep=4pt,topsep=4pt,label={\color{poporange}\textbullet}}
\setlength{\parindent}{0pt}
\setlength{\parskip}{5pt}
\renewcommand{\arraystretch}{1.25}

% chemfig : liaisons un peu plus épaisses, encre Pop
\setchemfig{atom sep=2em,bond style={line width=0.9pt,popink},cram width=3pt}

% pgfplots : style Pop par défaut
\pgfplotsset{
  every axis/.append style={axis line style={popink,line width=1pt},tick style={popink},
    grid style={popline,line width=0.5pt},label style={font=\small},tick label style={font=\footnotesize}},
  popcurve/.style={poporange,line width=1.8pt,no markers,smooth},
}

% ============================================================
% Pill / squiggle
% ============================================================
\newcommand{\poppill}[3]{%
  \tikz[baseline=-0.55ex]\node[draw=popink,line width=1.2pt,rounded corners=2.6pt,%
    fill=#2,inner xsep=6.5pt,inner ysep=3pt]{%
    \popxboldls\fontsize{7}{7}\selectfont\color{#3}\MakeUppercase{#1}};%
}
\newcommand{\popsquiggle}[1]{%
  \tikz[baseline=-0.4ex]{
    \draw[#1,line width=2.6pt,line cap=round]
      plot [smooth] coordinates {(0,0.09) (0.34,0.01) (0.68,0.21) (1.02,0.01) (1.36,0.10)};
  }%
}
% Mot-clé surligné (marqueur orange sous le texte)
\newcommand{\cle}[1]{{\popxbold\color{popdark}#1}}

% ============================================================
% En-tête / pied
% ============================================================
\pagestyle{fancy}
\fancyhf{}
\renewcommand{\headrulewidth}{0pt}
\renewcommand{\footrulewidth}{0pt}
\lhead{\raisebox{-0.15ex}{\poplogo\fontsize{13}{13}\selectfont\color{popink}OnBuch\color{poporange}+}}
\rhead{\poppill{\DOCMATIERE\ \textbullet\ \DOCNIVEAU\ \textbullet\ Leçon \DOCLECON}{white}{popink}}
\lfoot{\popxbold\fontsize{7.6}{7.6}\selectfont\color{popmuted}\MakeUppercase{Cours signé OnBuch+}\ \textbullet\ {\popsemi\DOCTITRE}}
\rfoot{\tikz[baseline=-0.6ex]\node[rounded corners=8.5pt,fill=popink,minimum height=17pt,minimum width=17pt,inner xsep=5pt,inner sep=0]{\popxbold\fontsize{8}{8}\selectfont\color{popcream}\thepage\,/\,\pageref*{LastPage}};}

% ============================================================
% Titres
% ============================================================
\newcommand{\chaptitre}[2]{%
  \begin{tcolorbox}[enhanced,colback=poporange,colframe=popink,boxrule=2.6pt,arc=11pt,
      drop shadow=popink,left=15pt,right=15pt,top=12pt,bottom=11pt,before skip=3pt,after skip=18pt]
    \poppill{#2}{popink}{popcream}\par\vspace{9pt}
    {\raggedright\hyphenpenalty=10000\exhyphenpenalty=10000\popdisplay\fontsize{22}{24}\selectfont\color{popcream}\MakeUppercase{#1}\par}
  \end{tcolorbox}
}
% Section numérotée : pastille orange avec le numéro + titre Archivo
\newcounter{popsec}
\newcommand{\coursec}[1]{%
  \stepcounter{popsec}\setcounter{popsub}{0}%
  \par\vspace{16pt}\noindent
  \begin{minipage}{\linewidth}
  \tikz[baseline=-0.7ex]\node[draw=popink,line width=1.6pt,fill=poporange,rounded corners=4pt,minimum size=22pt,inner sep=0]{\popdisplay\fontsize{12}{12}\selectfont\color{popcream}\thepopsec};%
  \hspace{8pt}{\popdisplay\fontsize{15}{17}\selectfont\color{popink}\MakeUppercase{#1}}\par
  \vspace{4pt}\noindent{\color{poporange}\rule{36pt}{3.6pt}}
  \end{minipage}\par\nopagebreak\vspace{8pt}
}
\newcounter{popsub}
\newcommand{\courssub}[1]{%
  \stepcounter{popsub}%
  \par\vspace{9pt}\noindent{\popxbold\fontsize{11.5}{13}\selectfont\color{popink}{\color{poporange}\thepopsec.\thepopsub}\hspace{6pt}#1}\par\nopagebreak\vspace{4pt}
}

% ============================================================
% Boîtes pédagogiques — même recette : fond blanc, bordure épaisse colorée,
% ombre dure encre, badge plein. Titre optionnel [#1].
% ============================================================
\tcbset{popbase/.style={breakable,enhanced,colback=white,boxrule=2.3pt,arc=9pt,
  shadow={3pt}{-3pt}{0pt}{popink},left=14pt,right=14pt,top=11pt,bottom=11pt,before skip=11pt,after skip=15pt,
  fontupper=\popsemi\fontsize{10.6}{14.5}\selectfont\color{popink},
  fontlower=\popsemi\fontsize{10.6}{14.5}\selectfont\color{popink}}}
% Coche dessinée (le glyphe ✓ n'existe pas dans les polices du document)
\newcommand{\popcheck}[1]{\tikz[baseline=-0.5ex]\draw[#1,line width=1.6pt,line cap=round,line join=round](-0.13,0.0)--(-0.04,-0.09)--(0.13,0.1);}
\newcommand{\popboxhead}[3]{\poppill{#1}{#2}{white}\ifx\relax#3\relax\else\hspace{6pt}{\popxbold\fontsize{10.6}{12}\selectfont\color{popink}#3}\fi\par\vspace{7pt}}

\newtcolorbox{aretenir}[1][]{popbase,colframe=popgreen,before upper={\popboxhead{À retenir}{popgreen}{#1}}}
\newtcolorbox{exemplebox}[1][]{popbase,colframe=poporange,before upper={\popboxhead{Exemple}{poporange}{#1}}}
\newtcolorbox{definition}[1][]{popbase,colframe=popblue,before upper={\popboxhead{Définition}{popblue}{#1}}}
\newtcolorbox{propriete}[1][]{popbase,colframe=poppurple,before upper={\popboxhead{Propriété}{poppurple}{#1}}}
\newtcolorbox{methode}[1][]{popbase,colframe=popgold,before upper={\popboxhead{Méthode}{popgold}{#1}}}
\newtcolorbox{attention}[1][]{popbase,colframe=poppink,before upper={\popboxhead{Attention}{poppink}{#1}}}
\newtcolorbox{experience}[1][]{popbase,colframe=popblue,colback=popblueL,before upper={\popboxhead{Expérience}{popblue}{#1}}}
% « Le savais-tu ? » : carte violette pleine (contexte, culture, Cameroun)
\newtcolorbox{savaistu}[1][]{popbase,colback=poppurple,colframe=popink,
  fontupper=\popsemi\fontsize{10.4}{14.2}\selectfont\color{white},
  before upper={\poppill{Le savais-tu\,?}{white}{poppurple}\ifx\relax#1\relax\else\hspace{6pt}{\popxbold\color{white}#1}\fi\par\vspace{7pt}}}
% Exercice résolu : énoncé en haut, \tcblower puis la solution sur fond crème
\newcounter{popexo}\newcounter{popexr}
\newtcolorbox{exoresolu}[1][]{popbase,colframe=poporange,colbacklower=popcream,
  segmentation style={popink,line width=1.2pt,dashed},
  before upper={\stepcounter{popexr}\popboxhead{Exercice résolu \thepopexr}{poporange}{#1}},
  before lower={\poppill{Solution}{popink}{popcream}\par\vspace{6pt}}}
% Exercice (non résolu en place) — la correction est dans la partie Corrigés
\newtcolorbox{exercice}[1][]{popbase,colframe=popink,
  before upper={\stepcounter{popexo}\popboxhead{Exercice \thepopexo}{popink}{#1}}}
% Corrigé
\newtcolorbox{corrige}[1][]{popbase,colframe=popgreen,colback=popgreenL,
  before upper={\popboxhead{Corrigé}{popgreen}{#1}}}
% Objectifs / prérequis / bilan
\newtcolorbox{objectifs}{popbase,colframe=popink,colback=poporangeL,
  before upper={\popboxhead{Objectifs de la leçon}{popink}{}}}
\newtcolorbox{prerequis}{popbase,colframe=popink,colback=popblueL,
  before upper={\popboxhead{Prérequis}{popblue}{}}}
\newtcolorbox{bilan}{popbase,colframe=popink,colback=white,boxrule=2.8pt,
  before upper={\popboxhead{Fiche bilan}{poporange}{L'essentiel à connaître pour l'examen}}}
% Figure encadrée avec légende
\newtcolorbox{popfigure}[1][]{enhanced,breakable=false,colback=white,colframe=popline,boxrule=1.4pt,arc=8pt,
  left=10pt,right=10pt,top=10pt,bottom=8pt,before skip=10pt,after skip=12pt,halign=center,
  lower separated=false,halign lower=center,
  fontlower=\popsemi\fontsize{9}{11.5}\selectfont\color{popmuted},
  #1}
\newcommand{\legende}[1]{\par\vspace{6pt}{\popsemi\fontsize{9}{11.5}\selectfont\color{popmuted}#1}}

% ============================================================
% Signature OnBuch+
% ============================================================
\newcommand{\poplogoinline}[1]{{\poplogo\fontsize{#1}{#1}\selectfont\color{popink}OnBuch\color{poporange}+}}
\newcommand{\signatureonbuch}{%
  \begin{tcolorbox}[enhanced,colback=popink,colframe=popink,boxrule=2.2pt,arc=10pt,
      drop shadow=poporange,left=14pt,right=14pt,top=10pt,bottom=10pt,before skip=0pt,after skip=0pt]
    \tikz[baseline=-0.7ex]\node[circle,fill=popgreen,draw=popcream,line width=1.3pt,minimum size=22pt,inner sep=0]{\popcheck{white}};%
    \hspace{9pt}\begin{minipage}[c]{0.62\linewidth}
      {\popxbold\fontsize{12}{13}\selectfont\color{popcream}Signé\ \textbullet\ {\poplogo OnBuch\color{poporange}+}}\par\vspace{2pt}
      {\raggedright\popsemi\fontsize{8}{10.5}\selectfont\color{popline}\MakeUppercase{Équipe pédagogique OnBuch+}\\\MakeUppercase{Conforme au programme officiel MINESEC}\par}
    \end{minipage}\hfill
    {\poplogo\fontsize{22}{22}\selectfont\color{popcream}OnBuch\color{poporange}+}
  \end{tcolorbox}%
}

% ============================================================
% Page de garde
% #1 titre · #2 matière · #3 niveau · #4 module · #5 n° leçon · #6 durée ·
% #7 description / objectifs (texte ou itemize)
% ============================================================
\newcommand{\pagegarde}[7]{%
  \thispagestyle{empty}
  \newgeometry{a4paper,margin=1.7cm,top=1.6cm,bottom=1.4cm}%
  \begin{tikzpicture}[remember picture,overlay]
    \fill[poporange!18] ([xshift=-3.2cm,yshift=-3.6cm]current page.north east) circle (4.6cm);
    \fill[poppurple!14] ([xshift=2.4cm,yshift=7.6cm]current page.south west) circle (3.8cm);
  \end{tikzpicture}%
  {\poplogo\fontsize{30}{30}\selectfont\color{popink}OnBuch\color{poporange}+}\hfill
  \raisebox{0.6ex}{\poppill{Cours signé \textbullet\ Premium}{popink}{popcream}}\par
  \vspace{1pt}\popsquiggle{poppurple}\par
  \vspace{8pt}
  \begin{tcolorbox}[enhanced,colback=poporange,colframe=popink,boxrule=2.8pt,arc=13pt,
      drop shadow=popink,left=18pt,right=18pt,top=12pt,bottom=13pt,after skip=10pt]
    \poppill{#2 \textbullet\ #3}{popink}{popcream}\hspace{5pt}\poppill{#4}{white}{popink}\par\vspace{8pt}
    {\popdisplay\fontsize{14}{14}\selectfont\color{popink}LEÇON #5}\par\vspace{4pt}
    {\raggedright\hyphenpenalty=10000\exhyphenpenalty=10000\popdisplay\fontsize{25}{27}\selectfont\color{popcream}\MakeUppercase{#1}\par}
    \vspace{8pt}
    {\popxbold\fontsize{9.5}{9.5}\selectfont\color{popink}\MakeUppercase{Durée conseillée : #6}}
  \end{tcolorbox}
  \begin{tcolorbox}[enhanced,colback=white,colframe=popink,boxrule=2.2pt,arc=10pt,
      drop shadow=popink,left=16pt,right=16pt,top=10pt,bottom=10pt,after skip=8pt]
    \poppill{Dans cette leçon}{poppurple}{white}\par\vspace{5pt}
    \popsemi\fontsize{9.3}{12.6}\selectfont\color{popink}#7
  \end{tcolorbox}
  \noindent\begin{minipage}[t]{0.487\linewidth}
    \begin{tcolorbox}[enhanced,colback=popgreen,colframe=popink,boxrule=2.2pt,arc=10pt,
        drop shadow=popink,left=11pt,right=11pt,top=10pt,bottom=10pt,height=3.1cm,valign=center]
      \tcbox[colback=white,colframe=popink,boxrule=1.3pt,arc=5pt,boxsep=0pt,nobeforeafter,
        left=3pt,right=3pt,top=3pt,bottom=3pt,valign=center]{\qrcode[height=1.9cm]{https://play.google.com/store/apps/details?id=com.luvvix.onbuch&hl=fr}}%
      \hspace{8pt}\begin{minipage}[c]{0.5\linewidth}
        {\popdisplay\fontsize{11}{12.5}\selectfont\color{white}\MakeUppercase{Télécharge OnBuch}}\par\vspace{4pt}
        {\raggedright\popsemi\fontsize{8.4}{11}\selectfont\color{white}Gratuit \textbullet\ Google Play\\Tuteur IA, annales, concours, résultats\par}
      \end{minipage}
    \end{tcolorbox}
  \end{minipage}\hfill
  \begin{minipage}[t]{0.487\linewidth}
    \begin{tcolorbox}[enhanced,colback=poppurple,colframe=popink,boxrule=2.2pt,arc=10pt,
        drop shadow=popink,left=11pt,right=11pt,top=10pt,bottom=10pt,height=3.1cm,valign=center]
      \tikz[baseline=-0.7ex]\node[circle,fill=white,draw=popink,line width=1.3pt,minimum size=1.6cm,inner sep=0]{\popdisplay\fontsize{24}{24}\selectfont\color{poppurple}+};%
      \hspace{8pt}\begin{minipage}[c]{0.6\linewidth}
        {\popdisplay\fontsize{11}{12.5}\selectfont\color{white}\MakeUppercase{Passe à OnBuch+}}\par\vspace{4pt}
        {\raggedright\popsemi\fontsize{8.4}{11}\selectfont\color{white}Tous les cours signés, Tuteur IA illimité, fiches PDF — onglet OnBuch+ de l'app\par}
      \end{minipage}
    \end{tcolorbox}
  \end{minipage}
  \vspace{8pt}
  \popcontenu
  \vfill
  \signatureonbuch
  \restoregeometry
  \newpage
}

% Rangée « Ce cours contient » (page de garde)
\newcommand{\poptile}[3]{% #1 couleur #2 gros texte #3 légende
  \begin{tcolorbox}[enhanced,colback=#1,colframe=popink,boxrule=2pt,arc=9pt,shadow={2.5pt}{-2.5pt}{0pt}{popink},
      left=6pt,right=6pt,top=9pt,bottom=9pt,halign=center,valign=center,height=1.75cm,width=\linewidth,nobeforeafter]
    {\popdisplay\fontsize{12.5}{13}\selectfont\color{white}\MakeUppercase{#2}}\par\vspace{3pt}
    {\centering\popsemi\fontsize{7.8}{9.5}\selectfont\color{white}#3\par}
  \end{tcolorbox}}
\newcommand{\popcontenu}{%
  \par\noindent{\popxboldls\fontsize{7.5}{7.5}\selectfont\color{popmuted}CE COURS SIGNÉ CONTIENT}\par\vspace{7pt}
  \noindent\begin{minipage}[t]{0.235\linewidth}\poptile{popblue}{Cours}{complet, illustré, pas à pas}\end{minipage}\hfill
  \begin{minipage}[t]{0.235\linewidth}\poptile{poporange}{Méthodes}{et exercices résolus}\end{minipage}\hfill
  \begin{minipage}[t]{0.235\linewidth}\poptile{poppink}{Exercices}{type examen + corrigés}\end{minipage}\hfill
  \begin{minipage}[t]{0.235\linewidth}\poptile{popgreen}{Fiche bilan}{l'essentiel à retenir}\end{minipage}\par}

% Sommaire
\newtcolorbox{sommaire}{enhanced,colback=white,colframe=popink,boxrule=2.2pt,arc=10pt,
  drop shadow=popink,left=18pt,right=18pt,top=13pt,bottom=13pt,before skip=6pt,after skip=20pt,
  before upper={\poppill{Sommaire}{poppurple}{white}\par\vspace{9pt}\popsemi\fontsize{10.6}{14.5}\selectfont\color{popink}}}

% Signature de fin de cours — poussée tout en bas de la dernière page
\newcommand{\findecours}{%
  \par\vfill\nopagebreak
  \begin{center}\popsquiggle{poporange}\end{center}\vspace{4pt}
  \signatureonbuch
}

%% FIGFIX-BEGIN v3 — correctifs globaux des figures (docs/onbuch-generation/tools/figfix)
\usepackage{adjustbox}\usepackage{etoolbox}
% toute figure TikZ/pgfplots plus large que la ligne est réduite à la largeur disponible
\BeforeBeginEnvironment{tikzpicture}{\begin{adjustbox}{max width=\linewidth}}
\AfterEndEnvironment{tikzpicture}{\end{adjustbox}}
% légendes pgfplots lisibles par-dessus les courbes
\pgfplotsset{every axis/.append style={legend style={fill=white,fill opacity=0.92,text opacity=1,draw=black!15,inner sep=3pt}}}
% étiquettes d'arêtes (syntaxe edge["..."]) avec fond blanc
\tikzset{every edge quotes/.append style={fill=white,inner sep=1.2pt}}
% bulles de cartes mentales : texte à la ligne dans la bulle, bulle un peu plus grande
\tikzset{every concept/.append style={text width=2.0cm,align=center,minimum size=2.3cm,inner sep=2pt}}
%% FIGFIX-END
\begin{document}

\pagegarde{Les acides $\alpha$-aminés et la synthèse des peptides}{Chimie}{Tle D}{Module 1 · Chimie organique}{04}{4 h}{Cette leçon présente les acides α-aminés, briques élémentaires des protéines : structure, stéréochimie (configurations D et L), propriétés acido-basiques et amphion. Elle explique ensuite la formation de la liaison peptidique, la nomenclature des peptides et la stratégie de synthèse sélective d'un dipeptide par protection et activation des fonctions.
\vspace{4pt}
\begin{multicols}{2}\small
\begin{itemize}
\item À la fin de cette leçon, je sais écrire la formule générale d'un acide α-aminé et citer les principaux acides aminés naturels avec leurs abréviations
\item À la fin de cette leçon, je sais nommer un acide α-aminé en nomenclature systématique et usuelle
\item À la fin de cette leçon, je sais identifier un carbone asymétrique dans une molécule d'acide α-aminé
\item À la fin de cette leçon, je sais construire la représentation de Fischer d'un acide α-aminé et distinguer les configurations D et L
\item À la fin de cette leçon, je sais écrire l'équilibre entre forme moléculaire et forme ionique (amphion) d'un acide α-aminé
\item À la fin de cette leçon, je sais écrire les équations-bilans montrant le caractère acide et le caractère basique de l'amphion et tracer le diagramme de prédominance
\end{itemize}\end{multicols}}

\begin{sommaire}
\begin{multicols}{2}
\begin{enumerate}
\item Structure et nomenclature des acides α-aminés
\item Stéréochimie : carbone asymétrique et configurations D/L
\item Propriétés acido-basiques : l'amphion
\item La liaison peptidique et les polypeptides
\item Synthèse sélective d'un dipeptide
\item Travaux pratiques
\item Méthodes et exercices résolus
\item Exercices
\item Corrigés des exercices
\item Fiche bilan
\end{enumerate}
\end{multicols}
\end{sommaire}

\begin{objectifs}
\begin{itemize}
\item À la fin de cette leçon, je sais écrire la formule générale d'un acide α-aminé et citer les principaux acides aminés naturels avec leurs abréviations
\item À la fin de cette leçon, je sais nommer un acide α-aminé en nomenclature systématique et usuelle
\item À la fin de cette leçon, je sais identifier un carbone asymétrique dans une molécule d'acide α-aminé
\item À la fin de cette leçon, je sais construire la représentation de Fischer d'un acide α-aminé et distinguer les configurations D et L
\item À la fin de cette leçon, je sais écrire l'équilibre entre forme moléculaire et forme ionique (amphion) d'un acide α-aminé
\item À la fin de cette leçon, je sais écrire les équations-bilans montrant le caractère acide et le caractère basique de l'amphion et tracer le diagramme de prédominance
\item À la fin de cette leçon, je sais identifier la liaison peptidique, écrire l'équation de formation d'un dipeptide et nommer un peptide
\item À la fin de cette leçon, je sais décrire les étapes de la synthèse sélective d'un dipeptide (blocage, activation, condensation, déblocage)
\end{itemize}
\end{objectifs}

\begin{prerequis}
\begin{itemize}
\item Fonctions organiques : acide carboxylique (–COOH) et amine (–NH2), couples acide/base (Première)
\item Nomenclature des composés organiques oxygénés et azotés
\item Stéréochimie : carbone asymétrique et énantiomères (début d'année de Terminale)
\item Réactions de condensation et d'hydrolyse (formation d'un ester, analogie utile)
\item Écriture et équilibrage d'équations-bilans en chimie organique
\end{itemize}
\end{prerequis}

\newpage

\coursec{Structure et nomenclature des acides \texorpdfstring{$\alpha$}{alpha}-aminés}

\textbf{Situation d'accroche.} Il est midi au marché Mokolo, à Yaoundé. Sur les étals, le poisson fumé, le ndolé bien vert, les arachides grillées et le soja attirent les acheteurs. « Ça donne de la force ! », affirme la vendeuse. Derrière cette expression populaire se cache une vérité scientifique profonde : tous ces aliments sont riches en \cle{protéines}. Or une protéine n'est rien d'autre qu'une longue chaîne de petites molécules appelées \cle{acides $\alpha$-aminés}, reliées entre elles par une liaison très particulière : la \cle{liaison peptidique}. Comment ces briques élémentaires sont-elles agencées ? Pourquoi l'organisme humain ne sait-il fabriquer que certaines d'entre elles ? Et comment le chimiste parvient-il, en laboratoire, à assembler deux acides aminés dans l'ordre voulu pour synthétiser un dipeptide précis ? Toute cette leçon répond à ces questions. Commençons par le commencement : qu'est-ce qu'un acide $\alpha$-aminé, et comment le nomme-t-on ?

\courssub{Qu'est-ce qu'un acide $\alpha$-aminé ?}

\textbf{L'intuition d'abord.} Le nom dit presque tout : un acide $\alpha$-aminé est une molécule qui possède \emph{à la fois} un caractère \emph{acide} et un groupe \emph{aminé}. Le caractère acide vient de la fonction \cle{acide carboxylique} \ce{-COOH}, que tu connais depuis la Première (souviens-toi de l'acide éthanoïque du vinaigre). Le caractère aminé vient de la fonction \cle{amine} \ce{-NH2}, dérivée de l'ammoniac \ce{NH3}. La lettre grecque $\alpha$ (alpha) précise \emph{où} se trouve l'amine : sur le carbone \emph{voisin} du groupe carboxyle, appelé \cle{carbone $\alpha$}.

\begin{definition}[Acide $\alpha$-aminé]
Un \cle{acide $\alpha$-aminé} est une molécule organique qui possède une fonction \textbf{acide carboxylique} \ce{-COOH} et une fonction \textbf{amine} \ce{-NH2} portées par le \textbf{même atome de carbone}, appelé \cle{carbone $\alpha$} (le carbone adjacent au groupe carboxyle).
\end{definition}

\textbf{Pourquoi « $\alpha$ » ?} En chimie organique, on repère les positions par rapport à une fonction principale à l'aide des lettres grecques : le carbone directement lié au \ce{-COOH} est le carbone $\alpha$, le suivant est le carbone $\beta$, puis $\gamma$, etc. Dire « acide $\alpha$-aminé » signifie donc précisément : \emph{le groupe \ce{-NH2} est sur le premier carbone après le \ce{-COOH}}.

\begin{attention}[Ne pas confondre !]
Tout acide aminé n'est pas un acide $\alpha$-aminé. Si le groupe \ce{-NH2} est porté par le carbone $\beta$ ou $\gamma$, la molécule est un acide aminé, certes, mais \textbf{pas} un acide $\alpha$-aminé. Ainsi :
\begin{itemize}
\item \ce{H2N-CH2-COOH} (glycine) : \ce{-NH2} sur le carbone $\alpha$ $\Rightarrow$ acide $\alpha$-aminé ;
\item \ce{H2N-CH2-CH2-COOH} (acide $\beta$-alanine) : \ce{-NH2} sur le carbone $\beta$ $\Rightarrow$ ce n'est \textbf{pas} un acide $\alpha$-aminé.
\end{itemize}
Au Baccalauréat, seuls les acides $\alpha$-aminés sont au programme : ce sont eux qui constituent les protéines.
\end{attention}

\courssub{Formule générale et nature du radical R}

Tous les acides $\alpha$-aminés naturels partagent le même squelette : un carbone central (le carbone $\alpha$) relié à quatre groupes : un hydrogène \ce{H}, un groupe amine \ce{-NH2}, un groupe carboxyle \ce{-COOH}, et un reste variable noté \cle{R}, appelé \textbf{radical} ou \textbf{chaîne latérale}.

\begin{aretenir}[Formule générale]
La formule générale d'un acide $\alpha$-aminé s'écrit :
\[ \ce{H2N-CHR-COOH} \]
C'est la nature du radical \ce{R} qui distingue un acide aminé d'un autre.
\end{aretenir}

\begin{popfigure}
\begin{center}
\chemfig{H_2N-[:30]C(-[:90]H)(-[:-90]R)-[:-30]C(=[:30]O)-[:-90]OH}
\end{center}
\legende{Formule générale d'un acide $\alpha$-aminé : les deux fonctions \ce{-NH2} et \ce{-COOH} sont portées par le même carbone $\alpha$ (central) ; le radical R varie d'un acide aminé à l'autre.}
\end{popfigure}

\textbf{La nature du radical R.} C'est lui qui donne à chaque acide aminé sa « personnalité » chimique :
\begin{itemize}
\item \ce{R} peut être un simple atome d'hydrogène : \ce{R = H} donne la \textbf{glycine}, le plus petit des acides aminés ;
\item \ce{R} peut être un groupe \textbf{alkyle} : \ce{-CH3} (alanine), \ce{-CH(CH3)2} (valine), \ce{-CH2-CH(CH3)2} (leucine) ;
\item \ce{R} peut porter un groupe \textbf{hydroxyle} \ce{-OH} : \ce{-CH2-OH} (sérine) ;
\item \ce{R} peut porter un groupe \textbf{thiol} \ce{-SH} : \ce{-CH2-SH} (cystéine) ;
\item \ce{R} peut contenir un \textbf{noyau aromatique} : \ce{-CH2-C6H5} (phénylalanine).
\end{itemize}

\begin{exemplebox}[Identifier les deux fonctions dans la sérine]
La sérine a pour formule semi-développée \ce{HO-CH2-CH(NH2)-COOH}. Identifions ses groupes :
\begin{itemize}
\item \ce{-COOH} à l'extrémité droite : fonction \textbf{acide carboxylique} ;
\item \ce{-NH2} porté par le carbone voisin du \ce{-COOH} : fonction \textbf{amine}, en position $\alpha$ ;
\item le radical est \ce{R = -CH2-OH} : il porte une fonction \textbf{alcool} supplémentaire.
\end{itemize}
La sérine est donc bien un acide $\alpha$-aminé, dont la chaîne latérale est hydroxylée.
\end{exemplebox}

\courssub{Les principaux acides $\alpha$-aminés naturels}

Une vingtaine d'acides $\alpha$-aminés entrent dans la composition des protéines des êtres vivants. Le programme de Terminale en retient six, que tu dois connaître avec leur \cle{abréviation à trois lettres} (code international utilisé pour écrire rapidement la composition d'un peptide).

\begin{center}
\begin{tabularx}{\linewidth}{|l|l|c|X|}
\hline
\rowcolor{popblueL} \textbf{Nom usuel} & \textbf{Formule semi-développée} & \textbf{Abrév.} & \textbf{Nature du radical R} \\
\hline
Glycine & \ce{H2N-CH2-COOH} & Gly & \ce{R = H} (hydrogène) \\
\hline
Alanine & \ce{H2N-CH(CH3)-COOH} & Ala & \ce{R = -CH3} (alkyle) \\
\hline
Valine & \ce{H2N-CH(CH(CH3)2)-COOH} & Val & \ce{R = -CH(CH3)2} (alkyle ramifié) \\
\hline
Leucine & \ce{H2N-CH(CH2CH(CH3)2)-COOH} & Leu & \ce{R = -CH2-CH(CH3)2} (alkyle ramifié) \\
\hline
Sérine & \ce{H2N-CH(CH2OH)-COOH} & Ser & \ce{R = -CH2-OH} (porteur d'un \ce{-OH}) \\
\hline
Phénylalanine & \ce{H2N-CH(CH2C6H5)-COOH} & Phe & \ce{R = -CH2-C6H5} (noyau aromatique) \\
\hline
\end{tabularx}
\end{center}

\begin{popfigure}
\begin{center}
\chemfig{H_2N-[:30]C(-[:90]CH_3)(-[:-90]H)-[:-30]C(=[:30]O)-[:-90]OH}
\hspace{2cm}
\chemfig{H_2N-[:30]C(-[:90]CH_2OH)(-[:-90]H)-[:-30]C(=[:30]O)-[:-90]OH}
\end{center}
\legende{Formules semi-développées de l'alanine (à gauche, \ce{R = -CH3}) et de la sérine (à droite, \ce{R = -CH2-OH}).}
\end{popfigure}

\begin{attention}[Piège classique au Bac]
Dans la glycine, \ce{R = H} : le carbone $\alpha$ porte alors \emph{deux} atomes d'hydrogène. Conséquence importante pour la suite de la leçon : la glycine est le \textbf{seul} acide $\alpha$-aminé naturel dont le carbone $\alpha$ n'est \textbf{pas asymétrique}. Ne l'oublie pas lors des questions de stéréochimie !
\end{attention}

\courssub{Nomenclature systématique et nomenclature usuelle}

\textbf{Nomenclature systématique.} L'Union internationale de chimie (UICPA) considère un acide $\alpha$-aminé comme un \emph{acide carboxylique} (fonction principale) portant un substituant \emph{amino} en position 2. D'où le nom général : \cle{acide 2-aminoalcanoïque}.

\begin{methode}[Nommer un acide $\alpha$-aminé en nomenclature systématique]
\begin{enumerate}
\item Identifier la \textbf{chaîne carbonée la plus longue} contenant le groupe \ce{-COOH} : elle donne le nom de l'acide alcanoïque correspondant (le carbone du \ce{-COOH} porte le numéro 1).
\item Remplacer la terminaison « -ane » de l'alcane par « \textbf{-oïque} », précédée de « acide ».
\item Repérer le groupe \ce{-NH2} : il est toujours en position \textbf{2} (carbone $\alpha$) ; ajouter le préfixe « \textbf{2-amino} ».
\item Nommer et numéroter les autres ramifications éventuelles (méthyl, éthyl\ldots) par ordre alphabétique.
\end{enumerate}
\end{methode}

\begin{exemplebox}[Application de la méthode]
\begin{itemize}
\item \textbf{Alanine} \ce{H2N-CH(CH3)-COOH} : chaîne de 3 carbones (acide propanoïque), \ce{-NH2} en position 2 $\Rightarrow$ \textbf{acide 2-aminopropanoïque}.
\item \textbf{Valine} \ce{H2N-CH(CH(CH3)2)-COOH} : chaîne la plus longue contenant \ce{-COOH} = 4 carbones (acide butanoïque), \ce{-NH2} en 2, un méthyle en 3 $\Rightarrow$ \textbf{acide 2-amino-3-méthylbutanoïque}.
\item \textbf{Glycine} \ce{H2N-CH2-COOH} : chaîne de 2 carbones $\Rightarrow$ \textbf{acide 2-aminoéthanoïque}.
\item \textbf{Sérine} : chaîne de 3 carbones, \ce{-NH2} en 2, \ce{-OH} en 3 $\Rightarrow$ \textbf{acide 2-amino-3-hydroxypropanoïque}.
\end{itemize}
\end{exemplebox}

\textbf{Nomenclature usuelle.} Bien avant la nomenclature systématique, les chimistes ont baptisé ces molécules de noms « triviaux », souvent liés à leur origine ou à leur goût. Ces noms restent les plus utilisés en biochimie et au Baccalauréat :
\begin{itemize}
\item la \textbf{glycine} est aussi appelée \cle{glycocolle} (du grec \emph{glykys}, « doux », car elle a un goût sucré) ;
\item l'\textbf{alanine}, la \textbf{valine}, la \textbf{leucine} (du grec \emph{leukos}, « blanc »), la \textbf{sérine} (isolée de la soie, \emph{sericum} en latin) et la \textbf{phénylalanine} complètent la liste à connaître.
\end{itemize}

\begin{exoresolu}[Nommer un acide aminé]
On considère la leucine de formule \ce{H2N-CH(CH2-CH(CH3)2)-COOH}. Donner son nom systématique.
\tcblower
\textbf{Étape 1.} Chaîne la plus longue contenant \ce{-COOH} : \ce{COOH}-\ce{CH}-\ce{CH2}-\ce{CH}-\ce{CH3}, soit \textbf{5 carbones} $\Rightarrow$ acide pentanoïque.

\textbf{Étape 2.} Le groupe \ce{-NH2} est sur le carbone 2 $\Rightarrow$ préfixe « 2-amino ».

\textbf{Étape 3.} Un groupe méthyle \ce{-CH3} est porté par le carbone 4 $\Rightarrow$ « 4-méthyl ».

\textbf{Conclusion.} La leucine est l'\textbf{acide 2-amino-4-méthylpentanoïque}. Vérification rapide : formule brute \ce{C6H13NO2}, masse molaire $M = 6\times 12 + 13\times 1 + 14 + 2\times 16 = \SI{131}{g.mol^{-1}}$.
\end{exoresolu}

\courssub{Des monomères aux protéines : les acides aminés essentiels}

Les acides $\alpha$-aminés sont les \cle{monomères} des protéines : comme des perles enfilées sur un collier, ils s'assemblent par des liaisons peptidiques (que nous étudierons dans la suite de la leçon) pour former des \textbf{polypeptides}, puis des \textbf{protéines} pouvant contenir des centaines de maillons. Une protéine alimentaire (poisson, viande, soja) est digérée en acides aminés, que l'organisme réassemble ensuite pour fabriquer ses propres protéines : muscles, enzymes, hémoglobine, anticorps\ldots

\begin{definition}[Acide aminé essentiel]
Un acide aminé est dit \cle{essentiel} (ou indispensable) lorsque l'organisme humain est \textbf{incapable de le synthétiser} : il doit alors obligatoirement être apporté par l'\textbf{alimentation}. C'est le cas, entre autres, de la \textbf{valine} et de la \textbf{leucine}.
\end{definition}

\begin{savaistu}[Le soja et le niébé, trésors du marché camerounais]
Au Cameroun, le soja et le niébé (haricot local) sont très consommés, et ce n'est pas un hasard : ces légumineuses sont exceptionnellement riches en protéines (jusqu'à \SI{35}{\%} en masse pour le soja), donc en acides aminés \emph{essentiels}. Associés à un tubercule comme le manioc ou le macabo — pauvre en protéines mais riche en énergie —, ils permettent un repas équilibré à moindre coût. Les nutritionnistes parlent de « complémentarité protéique » : les acides aminés absents d'un aliment sont fournis par l'autre. La chimie des acides $\alpha$-aminés est donc aussi une affaire de santé publique !
\end{savaistu}

\begin{aretenir}[L'essentiel de la section]
\begin{itemize}
\item Un acide $\alpha$-aminé possède \ce{-NH2} et \ce{-COOH} sur le \textbf{même} carbone (le carbone $\alpha$) : formule générale \ce{H2N-CHR-COOH}.
\item Le radical R caractérise chaque acide aminé : \ce{H} (glycine), alkyle (alanine, valine, leucine), \ce{-CH2-OH} (sérine), \ce{-CH2-C6H5} (phénylalanine).
\item Nom systématique : \textbf{acide 2-aminoalcanoïque} ; noms usuels et abréviations (Gly, Ala, Val, Leu, Ser, Phe) à connaître.
\item Les acides aminés essentiels (valine, leucine\ldots) doivent être apportés par l'alimentation.
\end{itemize}
\end{aretenir}

\coursec{Stéréochimie : carbone asymétrique et configurations D/L}

Dans la section précédente, nous avons établi la formule générale des acides $\alpha$-aminés, \ce{H2N-CHR-COOH}, et appris à les nommer. Mais une formule semi-développée ne dit pas tout : elle ne précise pas \emph{comment} les groupes d'atomes sont disposés dans l'espace autour du carbone $\alpha$. Or cette disposition spatiale a des conséquences biologiques immenses : c'est elle qui distingue un acide aminé « naturel », utilisable par ton organisme pour fabriquer ses protéines, de son image dans un miroir, biologiquement inactive. C'est toute la question de la \cle{stéréochimie} des acides $\alpha$-aminés, que nous abordons maintenant.

\courssub{Le carbone asymétrique}

\textbf{Une intuition pour commencer.}\par
Regarde tes deux mains. Elles sont constituées des mêmes « pièces » (une paume, quatre doigts, un pouce), et pourtant il est impossible de les superposer exactement : la main gauche est l'image de la main droite dans un miroir. On dit que les mains sont \emph{chirales} (du grec \emph{kheir}, la main). Certaines molécules possèdent exactement cette propriété : elles existent sous deux formes, images l'une de l'autre dans un miroir et \textbf{non superposables}. La cause la plus fréquente de cette chiralité en chimie organique est la présence d'un \cle{carbone asymétrique}.

\begin{definition}[Carbone asymétrique]
Un \cle{carbone asymétrique} est un atome de carbone \textbf{tétraédrique} (hybridation $sp^3$, quatre liaisons simples) lié à \textbf{quatre atomes ou groupes d'atomes tous différents}. On le repère par un astérisque : $C^*$.
\end{definition}

Appliquons cette définition au carbone $\alpha$ d'un acide $\alpha$-aminé de formule \ce{H2N-CHR-COOH}. Ce carbone porte quatre substituants :
\begin{enumerate}
\item le groupe amine \ce{-NH2} ;
\item le groupe carboxyle \ce{-COOH} ;
\item un atome d'hydrogène \ce{-H} ;
\item la chaîne latérale \ce{-R}.
\end{enumerate}
Ces quatre groupes sont différents \textbf{à condition que} \ce{R} soit différent de \ce{H}. Autrement dit :

\begin{propriete}[Chiralité des acides $\alpha$-aminés]
Tout acide $\alpha$-aminé possède un carbone $\alpha$ asymétrique, \textbf{sauf la glycine} (\ce{H2N-CH2-COOH}) pour laquelle $\ce{R} = \ce{H}$ : le carbone $\alpha$ porte alors deux atomes d'hydrogène identiques et n'est pas asymétrique.
\par\smallskip
Conséquence : tout acide $\alpha$-aminé autre que la glycine existe sous la forme de \textbf{deux énantiomères}, c'est-à-dire deux stéréoisomères images l'un de l'autre dans un miroir et non superposables.
\end{propriete}

\begin{methode}[Repérer un carbone asymétrique]
Pour savoir si une molécule possède un carbone asymétrique :
\begin{enumerate}
\item repérer tous les atomes de carbone \textbf{tétraédriques} (quatre liaisons simples ; éliminer d'office les carbones des doubles liaisons \ce{C=C} et \ce{C=O}, qui sont plans) ;
\item pour chacun, lister les quatre substituants ;
\item si les quatre substituants sont \textbf{tous différents}, le carbone est asymétrique : on le marque $C^*$ ;
\item dans un acide $\alpha$-aminé \ce{H2N-CHR-COOH}, le seul candidat est le carbone $\alpha$ : il est asymétrique sauf si $\ce{R} = \ce{H}$ (glycine).
\end{enumerate}
\end{methode}

\begin{attention}[Le piège de la glycine — classique du Bac]
La question « tous les acides $\alpha$-aminés sont-ils chiraux ? » revient très souvent au Baccalauréat. La réponse est \textbf{NON} : la \cle{glycine} (Gly), dont la chaîne latérale est un simple atome d'hydrogène, possède \emph{deux} atomes d'hydrogène sur son carbone $\alpha$. Elle est donc \textbf{achirale} (superposable à son image miroir) et optiquement inactive. C'est le seul des vingt acides aminés protéinogènes dans ce cas.
\end{attention}

\courssub{Le carbone $\alpha$ de l'alanine dans l'espace}

Pour bien visualiser, représentons le carbone $\alpha$ de l'alanine ($\ce{R} = \ce{CH3}$) dans l'espace, en respectant la géométrie qui sera codée par la projection de Fischer : les liaisons verticales (\ce{COOH} et \ce{R}) pointent vers l'arrière (hachurées), les liaisons horizontales (\ce{NH2} et \ce{H}) pointent vers l'avant (pleines). C'est la géométrie de la \textbf{L-alanine} (forme naturelle).

\begin{popfigure}
\begin{center}
\begin{tikzpicture}[scale=1.2, line join=round]
% Carbone central
\shade[ball color=popblue!60] (0,0) circle (0.35);
\node[white, font=\bfseries] at (0,0) {C};
% Liaison verticale haute : COOH (hachurée, vers l'arrière)
\foreach \i in {0,...,6} {
  \draw[popdark, thick] ({-0.15+0.05*\i}, {0.25+0.22*\i}) -- ({0.15-0.05*\i}, {0.25+0.22*\i});
}
\shade[ball color=poporange!80] (0,1.85) circle (0.30);
\node[white, font=\small\bfseries] at (0,1.85) {COOH};
% Liaison verticale basse : CH3 (hachurée, vers l'arrière)
\foreach \i in {0,...,6} {
  \draw[popdark, thick] ({-0.15+0.05*\i}, {-0.25-0.22*\i}) -- ({0.15-0.05*\i}, {-0.25-0.22*\i});
}
\shade[ball color=popgreen!80] (0,-1.85) circle (0.30);
\node[white, font=\small\bfseries] at (0,-1.85) {CH$_3$};
% Liaison horizontale gauche : NH2 (pleine, coin, vers l'avant)
\fill[poppink!80] (0,0) -- (-1.55,0.42) -- (-1.55,-0.42) -- cycle;
\shade[ball color=poppink!80] (-1.85,0) circle (0.30);
\node[white, font=\small\bfseries] at (-1.85,0) {NH$_2$};
% Liaison horizontale droite : H (pleine, coin, vers l'avant)
\fill[popmuted!70] (0,0) -- (1.55,0.42) -- (1.55,-0.42) -- cycle;
\shade[ball color=popmuted!70] (1.85,0) circle (0.28);
\node[white, font=\small\bfseries] at (1.85,0) {H};
% Annotations
\node[popdark, font=\small\itshape, align=center] at (-2.7,0.9) {coin plein :\\ vers l'avant};
\node[popdark, font=\small\itshape, align=center] at (2.7,0.9) {coin plein :\\ vers l'avant};
\node[popdark, font=\small\itshape, align=center] at (0,2.5) {hachuré :\\ vers l'arrière};
\node[popdark, font=\small\itshape, align=center] at (0,-2.5) {hachuré :\\ vers l'arrière};
\end{tikzpicture}
\end{center}
\legende{Le carbone $\alpha$ tétraédrique de la L-alanine : \ce{COOH} et \ce{CH3} (verticales) vers l'arrière (hachurés) ; \ce{NH2} (gauche) et \ce{H} (droite) vers l'avant (pleins). Les quatre substituants sont différents : le carbone est asymétrique ($C^*$).}
\end{popfigure}

Permute deux substituants quelconques (par exemple \ce{NH2} et \ce{H}) : tu obtiens l'autre énantiomère, l'image miroir de la première molécule. Aucune rotation dans l'espace ne permet de superposer les deux formes, exactement comme un gant gauche ne peut pas devenir un gant droit.

\courssub{La représentation de Fischer}

Dessiner des tétraèdres en perspective est fastidieux. Le chimiste allemand Emil Fischer (prix Nobel 1902) a proposé une convention de représentation plane, devenue la norme pour les acides aminés et les sucres.

\begin{definition}[Représentation de Fischer d'un acide $\alpha$-aminé]
Dans la \cle{représentation de Fischer}, le carbone asymétrique est au centre d'une croix :
\begin{itemize}
\item la \textbf{chaîne carbonée est verticale} : le groupe \ce{-COOH} est placé \textbf{en haut} et la chaîne latérale \ce{-R} \textbf{en bas} ;
\item le groupe \ce{-NH2} et l'hydrogène \ce{-H} sont placés sur \textbf{l'horizontale}.
\end{itemize}
\textbf{Règles de lecture} : les liaisons \emph{horizontales} sortent du plan de la feuille (elles pointent \textbf{vers l'avant}, vers l'observateur) ; les liaisons \emph{verticales} rentrent dans le plan (elles pointent \textbf{vers l'arrière}).
\end{definition}

\begin{attention}[Manipuler une projection de Fischer]
Une projection de Fischer n'est pas un dessin ordinaire : elle code une géométrie 3D précise.
\begin{itemize}
\item On peut la faire \textbf{tourner de \ang{180} dans le plan} de la feuille : la configuration est conservée.
\item Une rotation de \ang{90}, ou l'échange de deux substituants, \textbf{change la configuration} (on passe à l'énantiomère).
\end{itemize}
Ne « pivote » donc jamais une Fischer de 90° en croyant obtenir la même molécule !
\end{attention}

\courssub{Configurations D et L}

Fischer a introduit une convention de nomenclature simple, toujours utilisée pour les acides aminés.

\begin{definition}[Configurations D et L des acides $\alpha$-aminés]
Dans la représentation de Fischer d'un acide $\alpha$-aminé (\ce{COOH} en haut, \ce{R} en bas) :
\begin{itemize}
\item si le groupe \ce{-NH2} est \textbf{à gauche}, l'acide aminé est de \cle{configuration L} ;
\item si le groupe \ce{-NH2} est \textbf{à droite}, il est de \cle{configuration D}.
\end{itemize}
Les deux formes D et L d'un même acide aminé sont des \textbf{énantiomères}.
\end{definition}

\begin{propriete}[Les acides aminés naturels sont de configuration L]
Les protéines de tous les êtres vivants sont constituées (presque) exclusivement d'acides $\alpha$-aminés de \textbf{configuration L}. Les formes D sont très rares dans la nature.
\end{propriete}

\begin{popfigure}
\begin{center}
\begin{tikzpicture}
\node at (-3,2.2) {\textbf{L-alanine}};
\node at (-3,0) {\chemfig{C(-[:90]COOH)(-[:-90]CH_3)(-[:180]NH_2)(-[:0]H)}};
\node at (-3,-1.6) {\small \ce{NH2} à gauche};
\node at (3,2.2) {\textbf{D-alanine}};
\node at (3,0) {\chemfig{C(-[:90]COOH)(-[:-90]CH_3)(-[:0]NH_2)(-[:180]H)}};
\node at (3,-1.6) {\small \ce{NH2} à droite};
\draw[popmuted, dashed, thick] (0,2.4) -- (0,-2.0);
\node[popmuted, font=\small\itshape] at (0,-2.4) {miroir};
\end{tikzpicture}
\end{center}
\legende{Projections de Fischer des deux énantiomères de l'alanine : \ce{COOH} en haut, \ce{CH3} en bas. À gauche, la L-alanine (\ce{NH2} à gauche), forme naturelle ; à droite, la D-alanine, son image miroir non superposable.}
\end{popfigure}

\begin{exemplebox}[La L-sérine en représentation de Fischer]
La sérine a pour chaîne latérale $\ce{R} = \ce{-CH2OH}$. Pour obtenir la projection de Fischer de la \textbf{L-sérine} :
\begin{enumerate}
\item on place \ce{COOH} en haut et \ce{CH2OH} en bas (verticale) ;
\item la configuration L impose \ce{NH2} \textbf{à gauche}, donc \ce{H} à droite.
\end{enumerate}
\begin{center}
\chemfig{C(-[:90]COOH)(-[:-90]CH_2OH)(-[:180]NH_2)(-[:0]H)}
\end{center}
Vérification : le carbone central porte \ce{NH2}, \ce{H}, \ce{COOH} et \ce{CH2OH} : quatre groupes différents, il est bien asymétrique. La D-sérine serait l'image miroir, avec \ce{NH2} à droite.
\end{exemplebox}

\begin{savaistu}[Pourquoi la vie a-t-elle choisi la configuration L ?]
C'est l'un des grands mystères de la biochimie : les protéines de \emph{tous} les êtres vivants — de la bactérie au baobab, du tilapia du fleuve Sanaga à l'être humain — sont construites exclusivement avec des acides aminés \textbf{L}, alors que la chimie de laboratoire produit spontanément des mélanges 50/50 de D et de L. Les formes D existent mais sont marginales : on les trouve par exemple dans les \textbf{parois de certaines bactéries} et dans certains antibiotiques (la pénicilline agit justement en bloquant l'assemblage de ces parois contenant de la D-alanine). Une conséquence pratique : un acide aminé D de synthèse ne peut pas être incorporé dans tes protéines — ton organisme est « gaucher » à l'échelle moléculaire !
\end{savaistu}

\courssub{Activité optique}

Les deux énantiomères d'un acide $\alpha$-aminé ont les mêmes propriétés physiques (température de fusion, solubilité...) et les mêmes propriétés chimiques vis-à-vis de réactifs achiraux. Comment les distinguer alors ? Grâce à la \textbf{lumière polarisée}.

\begin{propriete}[Activité optique]
Les acides $\alpha$-aminés (sauf la glycine) sont \textbf{chiraux} : leurs solutions dévient le plan de polarisation de la lumière polarisée. On dit qu'ils sont \cle{optiquement actifs}.
\begin{itemize}
\item L'énantiomère qui dévie le plan vers la droite est dit \emph{dextrogyre} (noté $(+)$) ;
\item celui qui dévie vers la gauche est dit \emph{lévogyre} (noté $(-)$).
\end{itemize}
Un mélange équimolaire des deux énantiomères, appelé \cle{mélange racémique}, est optiquement inactif (les effets se compensent).
\end{propriete}

\begin{attention}[D/L n'est pas $(+)$/ $(-)$ !]
Ne confonds jamais les deux notations. \textbf{D} et \textbf{L} décrivent la \emph{configuration} (position du \ce{NH2} dans la Fischer), tandis que $(+)$ et $(-)$ décrivent le \emph{sens de déviation mesuré expérimentalement}. Il n'y a aucun lien systématique : la L-alanine est $(+)$, mais la L-sérine est $(-)$. Au Bac, on te demande la configuration D ou L, jamais le pouvoir rotatoire.
\end{attention}

\begin{exoresolu}[Configuration d'un acide aminé]
On donne la projection de Fischer suivante d'un acide $\alpha$-aminé : \ce{COOH} en haut, \ce{CH2-CH(CH3)2} en bas, \ce{NH2} à gauche et \ce{H} à droite.
\begin{enumerate}
\item Identifier cet acide aminé et donner son abréviation.
\item Justifier que la molécule est chirale.
\item Préciser sa configuration (D ou L). S'agit-il de la forme naturelle ?
\end{enumerate}
\tcblower
\textbf{1.} La chaîne latérale $\ce{R} = \ce{-CH2-CH(CH3)2}$ (groupe isobutyle) est celle de la \textbf{leucine}, d'abréviation \textbf{Leu}.
\par\smallskip
\textbf{2.} Le carbone $\alpha$ est tétraédrique et porte quatre groupes différents : \ce{-NH2}, \ce{-H}, \ce{-COOH} et \ce{-CH2-CH(CH3)2}. C'est un carbone asymétrique ($C^*$) : la molécule est chirale et existe sous deux énantiomères.
\par\smallskip
\textbf{3.} Dans la projection de Fischer, le groupe \ce{NH2} est \textbf{à gauche} : il s'agit de la \textbf{configuration L}. C'est donc la forme naturelle, celle qui entre dans la composition des protéines.
\end{exoresolu}

\begin{exercice}[À toi de jouer]
\begin{enumerate}
\item La valine a pour chaîne latérale $\ce{R} = \ce{-CH(CH3)2}$. Dessine la représentation de Fischer de la L-valine, puis de la D-valine.
\item Parmi les molécules suivantes, lesquelles possèdent un carbone asymétrique : la glycine, la phénylalanine ($\ce{R} = \ce{-CH2-C6H5}$), le 2-propanol \ce{CH3-CHOH-CH3} ?
\end{enumerate}
\end{exercice}

\begin{corrige}[Exercice]
\textbf{1.} L-valine : \ce{COOH} en haut, \ce{CH(CH3)2} en bas, \ce{NH2} à gauche, \ce{H} à droite. D-valine : image miroir, \ce{NH2} à droite et \ce{H} à gauche.
\par\smallskip
\textbf{2.} Glycine : non ($\ce{R} = \ce{H}$, deux H sur le C $\alpha$). Phénylalanine : oui, le C $\alpha$ porte \ce{NH2}, \ce{H}, \ce{COOH} et \ce{CH2-C6H5}, tous différents. 2-propanol : non, le carbone central porte deux groupes \ce{CH3} identiques.
\end{corrige}

\begin{aretenir}[L'essentiel de la section]
\begin{itemize}
\item Un \cle{carbone asymétrique} est un carbone tétraédrique lié à quatre groupes différents ; il rend la molécule \textbf{chirale} (deux énantiomères non superposables).
\item Tout acide $\alpha$-aminé possède un C $\alpha$ asymétrique, \textbf{sauf la glycine} ($\ce{R} = \ce{H}$).
\item En \cle{représentation de Fischer} : \ce{COOH} en haut, \ce{R} en bas (verticale, vers l'arrière) ; \ce{NH2} et \ce{H} sur l'horizontale (vers l'avant).
\item \ce{NH2} à gauche $\Rightarrow$ configuration \textbf{L} (forme naturelle des protéines) ; \ce{NH2} à droite $\Rightarrow$ configuration \textbf{D}.
\item Les acides $\alpha$-aminés chiraux sont \textbf{optiquement actifs} : ils dévient la lumière polarisée.
\end{itemize}
\end{aretenir}

Maintenant que nous connaissons la structure spatiale des acides $\alpha$-aminés, une question surgit : comment se comportent ces molécules en solution aqueuse, elles qui portent à la fois un groupe acide (\ce{-COOH}) et un groupe basique (\ce{-NH2}) ? C'est l'objet de la section suivante, consacrée à leurs remarquables propriétés acido-basiques et à l'amphion.

\coursec{Propriétés acido-basiques : l'amphion}

Dans la section précédente, nous avons vu que les acides $\alpha$-aminés naturels sont des molécules chirales de configuration L (sauf la glycine). Nous les avons représentés sous leur forme « classique » \ce{H2N-CHR-COOH}, comme s'il s'agissait de molécules ordinaires. Pourtant, un indice expérimental aurait dû nous mettre la puce à l'oreille : les acides $\alpha$-aminés sont des \textbf{solides cristallins}, solubles dans l'eau, avec des températures de fusion élevées (souvent supérieures à \SI{200}{\celsius})... exactement comme des \emph{sels}, et non comme des molécules organiques classiques ! Ce comportement étrange s'explique par une propriété remarquable : en solution aqueuse, un acide $\alpha$-aminé n'existe pratiquement pas sous la forme \ce{H2N-CHR-COOH}, mais sous la forme d'un \cle{ion dipolaire} appelé \cle{amphion}. C'est toute cette richesse acido-basique que nous allons maintenant explorer.

\courssub{Une molécule à double visage : acide et base à la fois}

\courssub{Le transfert interne de proton}

Reprenons la formule générale d'un acide $\alpha$-aminé et observons-la avec l'œil du chimiste acido-basicien :

\begin{itemize}
\item le groupe \textbf{carboxyle} \ce{-COOH} possède un hydrogène acide : c'est un \textbf{acide} (au sens de Brønsted), capable de \emph{céder} un proton \ce{H+} ;
\item le groupe \textbf{amine} \ce{-NH2} porte un doublet non liant sur l'azote : c'est une \textbf{base}, capable de \emph{capter} un proton \ce{H+}.
\end{itemize}

Or ces deux groupes cohabitent \emph{au sein de la même molécule}. L'idée intuitive est alors simple : pourquoi le proton du groupe acide irait-il chercher une base à l'extérieur, alors qu'une base se trouve juste à côté, sur le même squelette carboné ? C'est exactement ce qui se produit : il y a \textbf{transfert interne (intramoléculaire) de proton} du groupe \ce{-COOH} vers le groupe \ce{-NH2}.

\begin{popfigure}
\begin{center}
\chemfig{H_2N-[:30]CH(-[2]R)-[:-30]C(=[2]O)-[:-30]O-H}
\hspace{1,5cm}
\ce{<=>}
\hspace{1,5cm}
\chemfig{H_3N^{+}-[:30]CH(-[2]R)-[:-30]C(=[2]O)-[:-30]O^{-}}
\end{center}
\legende{Équilibre entre la forme moléculaire (non chargée) et la forme ionique (amphion). Le proton du groupe carboxyle migre vers le groupe amine.}
\end{popfigure}

L'équation-bilan de cet équilibre s'écrit :

\[ \ce{H2N-CHR-COOH <=> ^{+}H3N-CHR-COO^{-}} \]

\begin{definition}[Amphion ou zwitterion]
L'\cle{amphion} (ou \cle{zwitterion}, de l'allemand \emph{Zwitter} : « hybride, hermaphrodite ») est la forme ionique dipolaire d'un acide $\alpha$-aminé, de formule générale \ce{^{+}H3N-CHR-COO^{-}}. Il porte \textbf{à la fois une charge positive} (sur le groupe ammonium \ce{-NH3+}) \textbf{et une charge négative} (sur le groupe carboxylate \ce{-COO-}), tout en étant \textbf{globalement neutre} électriquement.
\end{definition}

\begin{attention}[L'amphion n'est pas « sans charge »]
Erreur classique au Baccalauréat : écrire que l'amphion est « une molécule neutre, donc sans charge ». C'est faux ! L'amphion porte \textbf{deux charges} : une charge $+$ et une charge $-$, situées à des endroits différents de la molécule. Sa \emph{charge totale} est nulle ($+1-1=0$), mais il est bien \emph{chargé deux fois}. C'est précisément cette double charge qui explique la solubilité dans l'eau, le caractère solide et le point de fusion élevé des acides aminés : ce sont des espèces \emph{ioniques}. Ne confonds jamais « globalement neutre » avec « non chargé ».
\end{attention}

\begin{aretenir}[L'essentiel sur l'équilibre moléculaire/ionique]
En solution aqueuse, l'équilibre \ce{H2N-CHR-COOH <=> ^{+}H3N-CHR-COO^{-}} est \textbf{très fortement déplacé vers l'amphion} : la forme moléculaire non chargée est pratiquement inexistante. Un acide $\alpha$-aminé en solution, c'est d'abord et avant tout un amphion.
\end{aretenir}

\courssub{L'amphion est un ampholyte}

Puisque l'amphion possède encore un groupe \ce{-NH3+} (acide, porteur de protons cédables) et un groupe \ce{-COO-} (base, capable de capter un proton), il peut réagir \emph{de deux façons opposées} selon le milieu dans lequel on le place.

\begin{definition}[Ampholyte]
Un \cle{ampholyte} (ou espèce \cle{amphotère}) est une espèce chimique qui peut se comporter soit comme un \textbf{acide} (céder un proton), soit comme une \textbf{base} (capter un proton), selon le partenaire auquel elle est confrontée. L'amphion \ce{^{+}H3N-CHR-COO^{-}} est un ampholyte ; l'eau \ce{H2O} et l'ion hydrogénocarbonate \ce{HCO3-} en sont d'autres exemples célèbres.
\end{definition}

\begin{methode}[Écrire les équations acido-basiques de l'amphion]
Pour écrire correctement les équations mettant en jeu l'amphion, procède toujours ainsi :
\begin{enumerate}
\item \textbf{Pars de l'amphion} \ce{^{+}H3N-CHR-COO^{-}} comme espèce de départ.
\item \textbf{Demande-toi quel rôle il joue.} Face à un acide fort (\ce{H3O+}), l'amphion joue le rôle de \emph{base} : il capte le proton. Face à une base forte (\ce{HO-}), il joue le rôle d'\emph{acide} : il cède un proton.
\item \textbf{Identifie le site qui réagit.} Quand il capte un proton, c'est le groupe \ce{-COO-} qui devient \ce{-COOH}. Quand il cède un proton, c'est le groupe \ce{-NH3+} qui devient \ce{-NH2}.
\item \textbf{Vérifie la conservation} des atomes et des charges électriques des deux côtés de la flèche.
\end{enumerate}
\end{methode}

\courssub{Caractère basique de l'amphion : le couple cation/amphion}

Plaçons l'amphion en milieu \textbf{acide} (par exemple en ajoutant une solution d'acide chlorhydrique, donc des ions \ce{H3O+}). Le groupe carboxylate \ce{-COO-} capte un proton : l'amphion joue le rôle de \textbf{base}.

\[ \ce{^{+}H3N-CHR-COO^{-} + H3O+ -> ^{+}H3N-CHR-COOH + H2O} \]

Commentons cette équation :
\begin{itemize}
\item l'amphion (charge totale $0$) capte \ce{H+} et devient le \textbf{cation} \ce{^{+}H3N-CHR-COOH} (charge totale $+1$) ;
\item la charge est conservée : $0 + 1 = +1 + 0$ ;
\item il s'agit d'une réaction \textbf{quasi totale} (réaction entre une base et l'acide fort \ce{H3O+}).
\end{itemize}

On en déduit le premier couple acide/base :

\begin{propriete}[Premier couple : cation/amphion]
\[ \ce{^{+}H3N-CHR-COOH} \;/\; \ce{^{+}H3N-CHR-COO^{-}} \qquad pK_{a1} \approx 2{,}3 \]
L'acide de ce couple est le \textbf{cation} (forme protonée des deux groupes) ; la base est l'\textbf{amphion}. Ce couple correspond à l'acidité du groupe \ce{-COOH} (rendue plus forte que dans un acide carboxylique ordinaire par la proximité du groupe \ce{-NH3+} attracteur d'électrons).
\end{propriete}

\courssub{Caractère acide de l'amphion : le couple amphion/anion}

Plaçons maintenant l'amphion en milieu \textbf{basique} (par exemple en ajoutant de la soude, donc des ions \ce{HO-}). Le groupe ammonium \ce{-NH3+} cède un proton : l'amphion joue le rôle d'\textbf{acide}.

\[ \ce{^{+}H3N-CHR-COO^{-} + HO- -> H2N-CHR-COO^{-} + H2O} \]

Commentons cette équation :
\begin{itemize}
\item l'amphion (charge totale $0$) cède \ce{H+} et devient l'\textbf{anion} \ce{H2N-CHR-COO^{-}} (charge totale $-1$) ;
\item la charge est conservée : $0 + (-1) = -1 + 0$ ;
\item la réaction est \textbf{quasi totale} (réaction entre un acide et la base forte \ce{HO-}).
\end{itemize}

On en déduit le second couple acide/base :

\begin{propriete}[Second couple : amphion/anion]
\[ \ce{^{+}H3N-CHR-COO^{-}} \;/\; \ce{H2N-CHR-COO^{-}} \qquad pK_{a2} \approx 9{,}7 \]
L'acide de ce couple est l'\textbf{amphion} ; la base est l'\textbf{anion}. Ce couple correspond à l'acidité du groupe ammonium \ce{-NH3+}.
\end{propriete}

\begin{aretenir}[Les deux couples de l'acide $\alpha$-aminé]
Un acide $\alpha$-aminé possède \textbf{deux couples acide/base} qui se succèdent :
\begin{center}
\ce{^{+}H3N-CHR-COOH <=>[acide][base] ^{+}H3N-CHR-COO^{-} <=>[acide][base] H2N-CHR-COO^{-}} \\[4pt]
\textbf{cation} \hspace{2,2cm} \textbf{amphion} \hspace{2,2cm} \textbf{anion}
\end{center}
L'amphion est la base du couple 1 ($pK_{a1} \approx 2{,}3$) et l'acide du couple 2 ($pK_{a2} \approx 9{,}7$) : c'est bien un ampholyte.
\end{aretenir}

\courssub{Diagramme de prédominance}

Chaque couple obéit à la relation de Henderson-Hasselbalch : $pH = pK_a + \log \dfrac{[\text{base}]}{[\text{acide}]}$. Une espèce \emph{prédomine} lorsque sa concentration est supérieure à celle de son partenaire conjugué. En superposant les deux couples sur un axe de pH, on obtient le diagramme de prédominance complet de l'acide $\alpha$-aminé.

\begin{popfigure}
\begin{center}
\begin{tikzpicture}[x=1.05cm,font=\small]
% Axe
\draw[-{Stealth},thick,popdark] (0,0) -- (14.5,0) node[below left=2pt and -30pt]{pH};
% Graduations
\foreach \x in {0,2,4,6,8,10,12,14}{\draw[popdark] (\x,0) -- (\x,-0.12) node[below]{\x};}
% pKa
\draw[popdark,thick] (2.3,0) -- (2.3,0.55) node[above,black]{$pK_{a1}\approx 2{,}3$};
\draw[popdark,thick] (9.7,0) -- (9.7,0.55) node[above,black]{$pK_{a2}\approx 9{,}7$};
% Domaines
\fill[popblue!25] (0,1.3) rectangle (2.3,2.6);
\fill[popgreen!25] (2.3,1.3) rectangle (9.7,2.6);
\fill[poporange!30] (9.7,1.3) rectangle (14.3,2.6);
\draw[popdark] (0,1.3) rectangle (14.3,2.6);
\draw[popdark] (2.3,1.3) -- (2.3,2.6);
\draw[popdark] (9.7,1.3) -- (9.7,2.6);
% Textes
\node[align=center] at (1.15,1.95) {\textbf{cation}\\ \ce{^{+}H3N-CHR-COOH}};
\node[align=center] at (6.1,1.95) {\textbf{amphion}\\ \ce{^{+}H3N-CHR-COO^{-}}};
\node[align=center] at (12.1,1.95) {\textbf{anion}\\ \ce{H2N-CHR-COO^{-}}};
% Charges
\node[popblue,below=2pt] at (1.15,1.3) {charge $+1$};
\node[popgreen!60!black,below=2pt] at (6.1,1.3) {charge $0$};
\node[poporange!80!black,below=2pt] at (12.1,1.3) {charge $-1$};
% pH 7
\draw[-{Stealth},poppurple,very thick] (7,-0.9) -- (7,-0.1);
\node[poppurple,below] at (7,-0.9) {pH $= 7$ : l'amphion prédomine};
\end{tikzpicture}
\end{center}
\legende{Diagramme de prédominance d'un acide $\alpha$-aminé. En dessous de $pK_{a1}$, le cation prédomine ; entre $pK_{a1}$ et $pK_{a2}$, c'est l'amphion ; au-dessus de $pK_{a2}$, c'est l'anion.}
\end{popfigure}

\begin{propriete}[Domaines de prédominance]
Pour un acide $\alpha$-aminé de $pK_{a1} \approx 2{,}3$ et $pK_{a2} \approx 9{,}7$ :
\begin{itemize}
\item si $pH < pK_{a1}$ : le \textbf{cation} \ce{^{+}H3N-CHR-COOH} prédomine (milieu acide) ;
\item si $pK_{a1} < pH < pK_{a2}$ : l'\textbf{amphion} \ce{^{+}H3N-CHR-COO^{-}} prédomine ;
\item si $pH > pK_{a2}$ : l'\textbf{anion} \ce{H2N-CHR-COO^{-}} prédomine (milieu basique).
\end{itemize}
En particulier, en solution neutre ($pH \approx 7$), l'acide $\alpha$-aminé se trouve \textbf{quasi exclusivement sous forme d'amphion}.
\end{propriete}

\begin{exemplebox}[L'alanine à pH = 7]
Pour l'alanine \ce{CH3-CH(NH2)-COOH}, on donne $pK_{a1} = 2{,}3$ et $pK_{a2} = 9{,}7$. À $pH = 7$, on a $pK_{a1} < 7 < pK_{a2}$ : la forme prédominante est donc l'amphion \ce{^{+}H3N-CH(CH3)-COO^{-}}. On peut même quantifier : d'après Henderson-Hasselbalch, $\dfrac{[\text{anion}]}{[\text{amphion}]} = 10^{7-9{,}7} = 10^{-2{,}7} \approx 2 \times 10^{-3}$, soit environ \SI{0,2}{\percent} d'anion seulement. L'amphion règne en maître !
\end{exemplebox}

\begin{exoresolu}[Forme prédominante de la glycine dans le sang]
Le sang a un pH voisin de $7{,}4$. La glycine \ce{H2N-CH2-COOH} possède les constantes $pK_{a1} = 2{,}4$ (couple cation/amphion) et $pK_{a2} = 9{,}8$ (couple amphion/anion).
\begin{enumerate}
\item Sous quelle forme la glycine se trouve-t-elle majoritairement dans le sang ?
\item Calculer le rapport $\dfrac{[\text{amphion}]}{[\text{cation}]}$ dans le sang.
\end{enumerate}
\tcblower
\textbf{1. Forme prédominante.} On compare le pH du milieu aux deux $pK_a$ :
\[ pK_{a1} = 2{,}4 \;<\; pH = 7{,}4 \;<\; pK_{a2} = 9{,}8 \]
Le pH du sang se situe \emph{entre} les deux $pK_a$ : c'est donc l'\textbf{amphion} \ce{^{+}H3N-CH2-COO^{-}} qui prédomine.

\textbf{2. Rapport des concentrations.} Le couple cation/amphion obéit à :
\[ pH = pK_{a1} + \log \frac{[\text{amphion}]}{[\text{cation}]} \]
d'où :
\[ \frac{[\text{amphion}]}{[\text{cation}]} = 10^{pH - pK_{a1}} = 10^{7{,}4 - 2{,}4} = 10^{5} \]
Il y a donc \textbf{cent mille fois plus} d'amphion que de cation dans le sang : le cation y est tout à fait négligeable, ce qui confirme la réponse à la question 1.
\end{exoresolu}

\begin{attention}[Pièges fréquents]
\begin{itemize}
\item \textbf{Ne pas inverser les rôles des couples} : c'est le couple de $pK_{a1}$ (le plus petit) qui correspond au groupe \ce{-COOH}/\ce{-COO-}, et celui de $pK_{a2}$ au groupe \ce{-NH3+}/\ce{-NH2}. Logique : le \ce{-COOH} est un acide plus fort que le \ce{-NH3+}.
\item \textbf{Ne jamais écrire} l'équation \ce{H2N-CHR-COOH + H2O <=> ...} comme si la forme moléculaire était l'espèce de référence en solution : l'espèce de référence, c'est l'amphion.
\item \textbf{Vérifier les charges} dans chaque équation-bilan : c'est le moyen le plus rapide de détecter une erreur. Dans \ce{amphion + H3O+ -> cation + H2O}, on a $0+1$ à gauche et $+1+0$ à droite.
\end{itemize}
\end{attention}

\begin{savaistu}[Le point isoélectrique et l'électrophorèse]
Il existe un pH très particulier, appelé \cle{point isoélectrique} et noté $pH_i$, pour lequel la concentration du cation est exactement égale à celle de l'anion : l'acide aminé est alors, en moyenne, électriquement neutre et \textbf{ne migre pas dans un champ électrique}. Pour un acide aminé à chaîne latérale neutre, on montre que :
\[ pH_i = \frac{pK_{a1} + pK_{a2}}{2} \]
Pour l'alanine : $pH_i = \dfrac{2{,}3 + 9{,}7}{2} = 6{,}0$. Cette propriété est à la base de l'\textbf{électrophorèse}, technique de séparation des acides aminés et des protéines : placées dans un champ électrique à un pH donné, les molécules migrent vers l'anode ou la cathode selon le signe de leur charge nette, et chacune s'immobilise lorsque le pH local atteint son $pH_i$. C'est ainsi que les laboratoires — y compris ceux des hôpitaux de Yaoundé et de Douala — analysent les protéines du sang, par exemple pour dépister la drépanocytose (anomalie de l'hémoglobine très répandue en Afrique centrale).
\end{savaistu}

\begin{exercice}[À toi de jouer]
La valine possède $pK_{a1} = 2{,}3$ et $pK_{a2} = 9{,}6$.
\begin{enumerate}
\item Écrire les deux couples acide/base de la valine en précisant pour chacun l'acide et la base.
\item Écrire l'équation-bilan de la réaction de l'amphion de la valine avec les ions \ce{H3O+}, puis avec les ions \ce{HO-}.
\item Quelle est la forme prédominante de la valine dans le jus de bissap, boisson acide de $pH \approx 3$ ? Justifier.
\item Calculer le point isoélectrique de la valine.
\end{enumerate}
\end{exercice}

\begin{corrige}[Exercice : la valine]
\textbf{1.} Couple 1 : \ce{^{+}H3N-CH(CH(CH3)2)-COOH}/\ce{^{+}H3N-CH(CH(CH3)2)-COO^{-}} (acide : cation ; base : amphion), $pK_{a1} = 2{,}3$. Couple 2 : \ce{^{+}H3N-CH(CH(CH3)2)-COO^{-}}/\ce{H2N-CH(CH(CH3)2)-COO^{-}} (acide : amphion ; base : anion), $pK_{a2} = 9{,}6$.

\textbf{2.} Avec \ce{H3O+} : \ce{^{+}H3N-CHR-COO^{-} + H3O+ -> ^{+}H3N-CHR-COOH + H2O} (l'amphion est une base). Avec \ce{HO-} : \ce{^{+}H3N-CHR-COO^{-} + HO- -> H2N-CHR-COO^{-} + H2O} (l'amphion est un acide).

\textbf{3.} À $pH = 3$, on a $pH > pK_{a1} = 2{,}3$ mais $pH < pK_{a2} = 9{,}6$ : l'\textbf{amphion} prédomine (le cation est minoritaire mais non négligeable : $[\text{amphion}]/[\text{cation}] = 10^{0{,}7} \approx 5$).

\textbf{4.} $pH_i = \dfrac{2{,}3 + 9{,}6}{2} = 5{,}95 \approx 6{,}0$.
\end{corrige}

\begin{aretenir}[Bilan de la section]
\begin{itemize}
\item En solution, un acide $\alpha$-aminé existe sous forme d'\textbf{amphion} \ce{^{+}H3N-CHR-COO^{-}}, ion dipolaire globalement neutre, issu d'un transfert interne de proton.
\item L'amphion est un \textbf{ampholyte} : base du couple cation/amphion ($pK_{a1} \approx 2{,}3$) et acide du couple amphion/anion ($pK_{a2} \approx 9{,}7$).
\item Le \textbf{diagramme de prédominance} impose : cation si $pH < pK_{a1}$, amphion si $pK_{a1} < pH < pK_{a2}$, anion si $pH > pK_{a2}$.
\item Au \textbf{point isoélectrique} $pH_i = \dfrac{pK_{a1}+pK_{a2}}{2}$, l'acide aminé ne migre pas en champ électrique : principe de l'électrophorèse.
\end{itemize}
\end{aretenir}

Maintenant que nous connaissons la structure, la stéréochimie et le comportement acido-basique des acides $\alpha$-aminés, une question naturelle se pose : comment ces petites molécules s'assemblent-elles entre elles pour former les immenses protéines qui font vivre nos cellules ? La réponse tient en une liaison chimique d'une importance capitale : la \textbf{liaison peptidique}, que nous étudions dans la section suivante.

\coursec{La liaison peptidique et les polypeptides}

Dans la section précédente, nous avons vu que chaque acide $\alpha$-aminé existe en solution sous forme d'\cle{amphion}, porteur à la fois d'un groupe ammonium \ce{-NH3+} (acide) et d'un groupe carboxylate \ce{-COO-} (basique). Cette double personnalité chimique a une conséquence extraordinaire : les acides $\alpha$-aminés peuvent \textbf{réagir entre eux}, la fonction acide de l'un attaquant la fonction amine de l'autre. C'est ainsi que la nature assemble, maillon après maillon, les immenses chaînes que sont les protéines : les enzymes qui digèrent ton repas de ndolé, l'hémoglobine qui transporte l'oxygène dans ton sang, la kératine de tes cheveux. Toutes sont construites sur le même principe, celui d'une liaison très particulière que nous allons maintenant étudier : la \cle{liaison peptidique}.

\courssub{La condensation entre deux acides $\alpha$-aminés}

\textbf{Intuition.} Imagine deux perles que l'on veut enfiler sur un collier : il faut les \emph{relier}. Pour deux acides $\alpha$-aminés, la « colle » naturelle consiste à faire réagir le groupe carboxyle \ce{-COOH} du premier avec le groupe amine \ce{-NH2} du second. Au cours de cette union, un atome d'oxygène et deux atomes d'hydrogène sont éliminés sous forme d'une molécule d'eau : c'est une \cle{condensation}.

\begin{definition}[Liaison peptidique]
La \cle{liaison peptidique} est la liaison \textbf{amide} \ce{-CO-NH-} formée par \textbf{condensation} entre le groupe carboxyle \ce{-COOH} d'un acide $\alpha$-aminé et le groupe amine \ce{-NH2} d'un autre acide $\alpha$-aminé, avec \textbf{élimination d'une molécule d'eau}. La molécule obtenue à partir de deux acides $\alpha$-aminés est un \cle{dipeptide}.
\end{definition}

Prenons l'exemple fondateur : la condensation entre la \textbf{glycine} (Gly, $\mathrm{R = H}$) et l'\textbf{alanine} (Ala, $\mathrm{R = CH_3}$). Le groupe \ce{-COOH} de la glycine réagit avec le groupe \ce{-NH2} de l'alanine :

\begin{popfigure}
\begin{center}
\begin{tikzpicture}[scale=1.0, every node/.style={font=\small, transform shape}]
% Réactifs
\node (gly) at (0,0) {\chemfig{H_2N-CH_2-C(=[1]O)(-[7]OH)}};
\node[popmuted, below=0.2cm] at (gly.south) {Glycine (Gly)};
\node at (3.5,0) {$+$};
\node (ala) at (6.5,0) {\chemfig{H_2N-CH(-[2]CH_3)-C(=[1]O)(-[7]OH)}};
\node[popmuted, below=0.2cm] at (ala.south) {Alanine (Ala)};
\node at (9.5,0) {$\xrightarrow[\text{condensation}]{\ -\ce{H2O}\ }$};
% Produit
\node (dipep) at (14,0) {\chemfig{H_2N-CH_2-C(=[1]O)-[7]NH-CH(-[2]CH_3)-C(=[1]O)(-[7]OH)}};
\node[popmuted, below=0.2cm] at (dipep.south) {Dipeptide Gly-Ala};
\node at (17.5,0) {$+$ \ \ce{H2O}};
% Encadrement liaison peptidique sur le produit
\draw[poporange, very thick, rounded corners=3pt] ([xshift=-1.8cm,yshift=-0.6cm]dipep.center) rectangle ([xshift=-0.1cm,yshift=0.9cm]dipep.center);
\node[poporange, font=\footnotesize\bfseries, anchor=south] at ([xshift=-0.95cm,yshift=0.9cm]dipep.center) {liaison peptidique};
\end{tikzpicture}
\legende{Condensation de la glycine et de l'alanine : formation du dipeptide Gly-Ala avec élimination d'eau. La liaison peptidique \ce{-CO-NH-} est encadrée en orange.}
\end{center}
\end{popfigure}

Écrivons l'équation-bilan sous forme semi-développée, comme attendu au Baccalauréat :

\[ \ce{H2N-CH2-COOH + H2N-CH(CH3)-COOH -> H2N-CH2-CO-NH-CH(CH3)-COOH + H2O} \]

Commentons cette équation-bilan avec la rigueur exigée :

\begin{itemize}
\item Il s'agit d'une \textbf{condensation} : deux molécules s'unissent avec élimination d'une petite molécule, ici \ce{H2O}.
\item L'atome d'oxygène de l'eau provient du groupe \ce{-COOH} (sous forme de \ce{-OH}), et les deux atomes d'hydrogène proviennent de ce même \ce{-OH} et du groupe \ce{-NH2} (qui perd un H pour devenir \ce{-NH-}).
\item La réaction est \textbf{lente} à température ordinaire ; elle est réalisée en pratique avec activation de la fonction acide (nous y reviendrons dans la section 5) ou catalysée par des enzymes (ribosomes) dans le vivant.
\item La réaction est \textbf{réversible} : la réaction inverse est l'\cle{hydrolyse}.
\end{itemize}

\begin{attention}[N'oublie pas la molécule d'eau !]
L'erreur la plus fréquente au Baccalauréat est d'écrire la formation du dipeptide \textbf{sans faire apparaître \ce{H2O}} dans les produits. L'équation est alors déséquilibrée (il manque O et H à droite) et le correcteur sanctionne. Vérifie toujours le bilan des atomes : à gauche, Gly + Ala comptent au total $\mathrm{C_5H_{10}N_2O_4}$ ; à droite, le dipeptide $\mathrm{C_5H_{10}N_2O_3}$ plus \ce{H2O} redonnent bien $\mathrm{C_5H_{10}N_2O_4}$.
\end{attention}

\begin{exemplebox}[Vérification du bilan atomique]
Comptons les atomes pour la condensation Gly + Ala $\rightarrow$ Gly-Ala + \ce{H2O} :
\begin{center}
\begin{tabularx}{\linewidth}{|l|X|X|X|X|}
\hline
\rowcolor{popblueL} & \textbf{C} & \textbf{H} & \textbf{N} & \textbf{O} \\
\hline
Réactifs (Gly + Ala) & $2+3=5$ & $5+7=12$ & $1+1=2$ & $2+2=4$ \\
\hline
Produits (Gly-Ala + \ce{H2O}) & $5$ & $10+2=12$ & $2$ & $3+1=4$ \\
\hline
\end{tabularx}
\end{center}
Le bilan est équilibré : l'équation est correcte.
\end{exemplebox}

\courssub{Géométrie de la liaison peptidique : un groupe plan}

\textbf{Intuition.} Une liaison simple \ce{C-N} « ordinaire » (comme dans une amine) permet une rotation libre autour d'elle : la molécule peut se tortiller. Dans la liaison peptidique, il n'en est rien : le groupe \ce{-CO-NH-} est \textbf{rigide et plan}. Pourquoi ?

L'atome d'azote possède un doublet non liant qui se \emph{délocalise} vers le groupe carbonyle voisin \ce{C=O}. On peut représenter cette délocalisation par deux formes limites (mésomères) :

\begin{center}
\chemfig{-[::-30]C(=[::60]O)-[::-60]N(-[::60]H)-[::-30]}
\quad
\ce{<->}
\quad
\chemfig{-[::-30]C(-[::60]\charge{90=\ominus}{O})=_[::-60]\charge{90=\oplus}{N}(-[::60]H)-[::-30]}
\quad
(\text{formes limites de résonance de l'amide})
\end{center}

La liaison \ce{C-N} acquiert ainsi un \textbf{caractère partiel de double liaison}. Or, on ne peut pas tourner autour d'une double liaison sans la rompre : la rotation autour de \ce{C-N} est \textbf{bloquée}.

\begin{propriete}[Planéité du groupe amide]
Dans une liaison peptidique, les quatre atomes \ce{C}, \ce{O}, \ce{N} et \ce{H} du groupe \ce{-CO-NH-} — ainsi que les deux atomes de carbone $\alpha$ voisins — sont situés \textbf{dans un même plan}. La rotation autour de la liaison \ce{C-N} est bloquée en raison de son caractère partiel de double liaison. La longueur de la liaison \ce{C-N} peptidique (environ \SI{0,132}{nm}) est intermédiaire entre une simple liaison \ce{C-N} (\SI{0,149}{nm}) et une double liaison \ce{C=N} (\SI{0,127}{nm}), ce qui confirme expérimentalement ce caractère partiel.
\end{propriete}

\begin{popfigure}
\begin{center}
\begin{tikzpicture}[scale=1.3, font=\small]
% Plan grisé
\fill[popblue!15] (-3.5,-1.5) -- (3.5,-1.5) -- (4.2,1.8) -- (-2.8,1.8) -- cycle;
\node[popmuted, font=\footnotesize\itshape] at (3.2,1.45) {plan du groupe amide};
% Atomes (nœuds nommés pour liaisons)
\node[circle, draw=popdark, fill=popcream, minimum size=0.7cm] (Cat) at (-0.9,0.15) {\ce{C}};
\node[circle, draw=popdark, fill=popcream, minimum size=0.7cm] (Oat) at (-2.1,1.1) {\ce{O}};
\node[circle, draw=popdark, fill=popcream, minimum size=0.7cm] (Nat) at (0.9,0.15) {\ce{N}};
\node[circle, draw=popdark, fill=popcream, minimum size=0.7cm] (Hat) at (2.1,1.1) {\ce{H}};
\node[circle, draw=popmuted, fill=white, minimum size=0.65cm] (Ca1) at (-2.2,-0.9) {\ce{C_{\alpha 1}}};
\node[circle, draw=popmuted, fill=white, minimum size=0.65cm] (Ca2) at (2.2,-0.9) {\ce{C_{\alpha 2}}};
% Liaisons
\draw[very thick, double distance=1.5pt] (Cat) -- (Oat); % Double liaison C=O
\draw[very thick, poporange] (Cat) -- (Nat) node[midway, below, font=\footnotesize\itshape, popdark]{caractère partiel double liaison};
\draw[very thick] (Nat) -- (Hat);
\draw[very thick] (Cat) -- (Ca1);
\draw[very thick] (Nat) -- (Ca2);
% Flèche rotation bloquée
\draw[poporange, very thick, -{Stealth[length=3mm]}, bend left=40] (0.2,0.8) to (1.6,0.8);
\node[poporange, font=\footnotesize\bfseries] at (0.9,1.4) {rotation bloquée};
\end{tikzpicture}
\legende{Géométrie plane du groupe peptidique : les atomes \ce{C}, \ce{O}, \ce{N}, \ce{H} et les deux carbones $\alpha$ adjacents sont coplanaires. La délocalisation électronique confère à la liaison \ce{C-N} un caractère partiel de double liaison qui bloque la rotation.}
\end{center}
\end{popfigure}

\begin{savaistu}[Pourquoi cette planéité est-elle capitale ?]
C'est la rigidité des liaisons peptidiques qui permet à une chaîne protéique de se replier de manière précise et reproductible en hélices ($\alpha$-hélice) et en feuillets ($\beta$-feuillet), comme un ruban que l'on plierait selon des charnières fixes. Ce repliement détermine la \textbf{fonction} de la protéine : une enzyme ne reconnaît sa « cible » que si sa forme dans l'espace est exacte. Une protéine « dépliée » (on dit \emph{dénaturée}, par exemple par la chaleur lors de la cuisson d'un œuf ou du poisson) perd son activité biologique.
\end{savaistu}

\courssub{Polypeptides et protéines}

La condensation peut se poursuivre : le dipeptide possède encore un groupe \ce{-NH2} libre à une extrémité et un groupe \ce{-COOH} libre à l'autre. Il peut donc condenser avec un troisième acide $\alpha$-aminé, puis un quatrième, et ainsi de suite.

\begin{definition}[Polypeptide et protéine]
Un \cle{polypeptide} est une chaîne formée par l'enchaînement de \textbf{nombreux} acides $\alpha$-aminés (résidus) reliés par des liaisons peptidiques. Une \cle{protéine} est un polypeptide de grande masse molaire (généralement plus d'une centaine de résidus) possédant une structure spatiale et une fonction biologique déterminées. Chaque acide $\alpha$-aminé engagé dans la chaîne est appelé un \cle{résidu}.
\end{definition}

Toute chaîne peptidique possède deux extrémités distinctes :

\begin{itemize}
\item l'\cle{extrémité N-terminale} : le bout de chaîne qui porte le groupe amine \ce{-NH2} (ou \ce{-NH3+}) \textbf{libre} ;
\item l'\cle{extrémité C-terminale} : le bout de chaîne qui porte le groupe carboxyle \ce{-COOH} (ou \ce{-COO-}) \textbf{libre}.
\end{itemize}

\begin{popfigure}
\begin{center}
% Utilisation de chemfig avec ancres nommées pour annotation robuste
\chemfig{
    @{Nterm}H_2N-CH(-[2]CH_3)-C(=[1]O)-[7]NH-CH(-[2]CH_2OH)-C(=[1]O)-[7]NH-CH_2-C(=[1]O)(-[7]OH)@{Cterm}
}
\chemmove{
    % Extrémité N-terminale
    \draw[popgreen, very thick, rounded corners=3pt] ([xshift=-4pt,yshift=-4pt]Nterm.north west) rectangle ([xshift=4pt,yshift=4pt]Nterm.north east);
    \node[popgreen, font=\footnotesize\bfseries, align=center, anchor=north] at ([yshift=-6pt]Nterm.south west) {extrémité\\N-terminale};
    % Extrémité C-terminale
    \draw[poppink, very thick, rounded corners=3pt] ([xshift=-4pt,yshift=-4pt]Cterm.south west) rectangle ([xshift=4pt,yshift=4pt]Cterm.north east);
    \node[poppink, font=\footnotesize\bfseries, align=center, anchor=north] at ([yshift=-6pt]Cterm.south) {extrémité\\C-terminale};
    % Liaisons peptidiques (approximation visuelle sur les atomes C=O et NH)
    % On repère les liaisons amide : 1ère entre Ala et Ser, 2ème entre Ser et Gly
    % Coordonnées relatives complexes en chemfig pur, on utilise des boîtes relatives aux nœuds si possible
    % Alternative simple : annotations textuelles
    \node[poporange, font=\footnotesize\bfseries, anchor=south] at ($(Nterm)!0.35!(Cterm) + (0,1.2)$) {liaisons peptidiques};
}
\legende{Le tripeptide Ala-Ser-Gly : trois résidus (Ala, Ser, Gly) reliés par deux liaisons peptidiques. À gauche, l'extrémité N-terminale (\ce{-NH2} libre) ; à droite, l'extrémité C-terminale (\ce{-COOH} libre).}
\end{center}
\end{popfigure}

\begin{savaistu}[L'insuline, un polypeptide qui sauve des vies]
L'\textbf{insuline}, hormone qui régule la glycémie (taux de glucose dans le sang), est un polypeptide de \textbf{51 acides aminés} répartis sur deux chaînes reliées par des ponts disulfure (\ce{-S-S-}). Chez les personnes diabétiques, sa production est insuffisante : elles doivent s'injecter de l'insuline, aujourd'hui produite industriellement par des bactéries génétiquement modifiées (recombinantes). C'est l'une des premières protéines dont la séquence complète fut déterminée (Frederick Sanger, prix Nobel de chimie 1958).
\end{savaistu}

\courssub{Nomenclature des peptides : l'ordre des résidus compte}

Pour nommer un peptide sans écrire toute sa formule, on utilise une convention simple et universelle.

\begin{methode}[Nommer un peptide]
\begin{enumerate}
\item Repérer l'\textbf{extrémité N-terminale} (le résidu dont le groupe \ce{-NH2} est libre).
\item Énumérer les résidus \textbf{de l'extrémité N-terminale vers l'extrémité C-terminale}, en utilisant leurs abréviations à trois lettres (Gly, Ala, Val, Leu, Ser, Phe\ldots) séparées par des tirets.
\item Le dernier résidu cité est celui qui porte le \ce{-COOH} libre (C-terminale).
\end{enumerate}
Exemple : \textbf{Ala-Val-Gly} désigne le tripeptide dont l'alanine porte le \ce{-NH2} libre, la valine est au milieu et la glycine porte le \ce{-COOH} libre.
\end{methode}

\begin{attention}[Gly-Ala et Ala-Gly sont deux molécules différentes]
L'ordre des résidus est \textbf{essentiel}. Dans \textbf{Gly-Ala}, c'est le \ce{-COOH} de la glycine qui a réagi avec le \ce{-NH2} de l'alanine : la glycine est N-terminale. Dans \textbf{Ala-Gly}, c'est l'inverse : l'alanine est N-terminale. Ce sont deux dipeptides \textbf{distincts}, de propriétés physiques et chimiques différentes — ce ne sont pas des isomères interchangeables ! De même, Ala-Val-Gly, Val-Ala-Gly et Gly-Ala-Val sont trois tripeptides différents. En général, avec $n$ acides $\alpha$-aminés différents, on peut former $n!$ peptides différents (par exemple $3! = 6$ tripeptides avec trois acides aminés distincts).
\end{attention}

\courssub{Hydrolyse des peptides et des protéines}

La condensation étant réversible, la réaction inverse existe : c'est l'\cle{hydrolyse}. Une molécule d'eau « casse » la liaison peptidique et régénère les acides $\alpha$-aminés :

\[ \ce{H2N-CH2-CO-NH-CH(CH3)-COOH + H2O -> H2N-CH2-COOH + H2N-CH(CH3)-COOH} \]

\begin{propriete}[Hydrolyse]
L'hydrolyse d'une liaison peptidique est la réaction \textbf{inverse de la condensation} : l'eau régénère un groupe \ce{-COOH} d'un côté et un groupe \ce{-NH2} de l'autre. Elle est réalisée :
\begin{itemize}
\item en laboratoire, à chaud, en milieu acide (\ce{HCl} concentré, vers \SI{110}{\celsius}) ou basique (\ce{NaOH}) pour l'analyse des protéines (dosage des acides aminés) ;
\item dans l'organisme, lors de la \textbf{digestion}, grâce à des enzymes (pepsine dans l'estomac, trypsine et chymotrypsine du pancréas) qui découpent les protéines des aliments (viande, poisson fumé, niébé, arachides\ldots) en acides $\alpha$-aminés absorbables par l'intestin.
\end{itemize}
\end{propriete}

C'est tout le sens de la situation de notre famille : l'organisme humain ne sait pas fabriquer tous les acides $\alpha$-aminés (certains sont dits \emph{essentiels}, comme la valine, la leucine ou la phénylalanine). Il doit donc \textbf{manger des protéines}, les hydrolyser lors de la digestion, puis réassembler les acides aminés récupérés en ses propres protéines, grâce à la condensation.

\begin{exoresolu}[Formation et hydrolyse d'un dipeptide]
\textbf{Énoncé.} On fait réagir la valine \ce{H2N-CH(CH(CH3)2)-COOH} avec la glycine \ce{H2N-CH2-COOH} dans les conditions où seul le dipeptide Val-Gly se forme.

1. Écrire l'équation-bilan de la condensation.
2. Entourer la liaison peptidique et indiquer les extrémités N- et C-terminales.
3. Écrire l'équation-bilan de l'hydrolyse totale de ce dipeptide.
4. Calculer la masse molaire du dipeptide Val-Gly. On donne : $M(\ce{C}) = \SI{12}{g.mol^{-1}}$, $M(\ce{H}) = \SI{1}{g.mol^{-1}}$, $M(\ce{N}) = \SI{14}{g.mol^{-1}}$, $M(\ce{O}) = \SI{16}{g.mol^{-1}}$.

\tcblower
\textbf{Solution détaillée.}

\textbf{1. Équation de condensation.} Le \ce{-COOH} de la valine (N-terminale) réagit avec le \ce{-NH2} de la glycine, avec élimination d'eau :
\[ \ce{H2N-CH(CH(CH3)2)-COOH + H2N-CH2-COOH -> H2N-CH(CH(CH3)2)-CO-NH-CH2-COOH + H2O} \]

\textbf{2. Liaison peptidique et extrémités.} La liaison peptidique est le groupe \ce{-CO-NH-} situé entre les deux carbones $\alpha$. L'extrémité \textbf{N-terminale} est le \ce{-NH2} libre porté par la valine (à gauche) ; l'extrémité \textbf{C-terminale} est le \ce{-COOH} libre de la glycine (à droite). Le peptide se nomme donc \textbf{Val-Gly} (et non Gly-Val).

\textbf{3. Hydrolyse totale.} C'est la réaction inverse :
\[ \ce{Val-Gly + H2O -> H2N-CH(CH(CH3)2)-COOH + H2N-CH2-COOH} \]
L'eau restaure le \ce{-COOH} de la valine et le \ce{-NH2} de la glycine.

\textbf{4. Masse molaire.} Formule brute de la valine : $\mathrm{C_5H_{11}NO_2}$ ($M = \SI{117}{g.mol^{-1}}$) ; de la glycine : $\mathrm{C_2H_5NO_2}$ ($M = \SI{75}{g.mol^{-1}}$). Le dipeptide a pour formule la somme des deux \textbf{moins} \ce{H2O} :
\[ \mathrm{C_5H_{11}NO_2 + C_2H_5NO_2 - H_2O = C_7H_{14}N_2O_3} \]
\begin{align*}
M(\text{Val-Gly}) &= 7 \times 12 + 14 \times 1 + 2 \times 14 + 3 \times 16 \\
&= 84 + 14 + 28 + 48 = \SI{174}{g.mol^{-1}}.
\end{align*}
Remarque utile : la masse molaire d'un peptide s'obtient toujours en retranchant $M(\ce{H2O}) = \SI{18}{g.mol^{-1}}$ \textbf{par liaison peptidique formée} de la somme des masses des acides aminés : $M = 117 + 75 - 18 = \SI{174}{g.mol^{-1}}$.
\end{exoresolu}

\begin{exercice}[À toi de jouer]
1. Écrire l'équation-bilan de formation du dipeptide \textbf{Ala-Gly} et montrer qu'il diffère de Gly-Ala.
2. Combien de tripeptides différents peut-on former avec les trois acides aminés Gly, Ala et Ser, chacun utilisé une seule fois ? Les citer tous.
3. L'hydrolyse totale d'un tripeptide fournit 2 molécules de glycine et 1 molécule d'alanine. Proposer deux séquences possibles pour ce tripeptide.
\end{exercice}

\begin{corrige}[Exercice]
1. \ce{H2N-CH(CH3)-COOH + H2N-CH2-COOH -> H2N-CH(CH3)-CO-NH-CH2-COOH + H2O}. Dans Ala-Gly, l'alanine porte le \ce{-NH2} libre (N-terminale) ; dans Gly-Ala, c'est la glycine. Les liaisons peptidiques relient des atomes de carbone $\alpha$ différents, les deux dipeptides sont des isomères de constitution distincts.

2. Avec trois résidus différents utilisés une fois chacun : $3! = 6$ tripeptides : Gly-Ala-Ser, Gly-Ser-Ala, Ala-Gly-Ser, Ala-Ser-Gly, Ser-Gly-Ala, Ser-Ala-Gly.

3. Le tripeptide contient les résidus Gly, Gly, Ala. Séquences possibles : Gly-Gly-Ala, Gly-Ala-Gly, Ala-Gly-Gly (trois possibilités, deux suffisent).
\end{corrige}

\begin{aretenir}[L'essentiel de la section]
\begin{itemize}
\item La \textbf{liaison peptidique} est une liaison amide \ce{-CO-NH-} formée par condensation entre \ce{-COOH} et \ce{-NH2}, avec élimination de \ce{H2O}.
\item Le groupe peptidique est \textbf{plan} : la liaison \ce{C-N} a un caractère partiel de double liaison (résonance) qui bloque la rotation.
\item Un \textbf{polypeptide} est une chaîne de résidus ; il possède une extrémité \textbf{N-terminale} (\ce{-NH2} libre) et une extrémité \textbf{C-terminale} (\ce{-COOH} libre).
\item On nomme un peptide en citant les résidus \textbf{depuis l'extrémité N-terminale} : Gly-Ala $\neq$ Ala-Gly.
\item L'\textbf{hydrolyse} (digestion, analyse) est la réaction inverse : elle régénère les acides $\alpha$-aminés.
\end{itemize}
\end{aretenir}

\coursec{Synthèse sélective d'un dipeptide}

Dans la section précédente, nous avons découvert la \cle{liaison peptidique} et vu qu'elle se forme par condensation entre la fonction \ce{-COOH} d'un acide $\alpha$-aminé et la fonction \ce{-NH2} d'un autre, avec élimination d'une molécule d'eau. Une question se pose alors naturellement : si l'on mélange simplement deux acides $\alpha$-aminés différents, obtient-on le dipeptide souhaité ? La réponse est non, et c'est tout l'objet de cette section : comprendre pourquoi, puis apprendre la \cle{stratégie de synthèse sélective} qui permet d'obtenir \emph{un seul} dipeptide, dans l'ordre voulu. C'est une compétence classique et très rentable au Baccalauréat.

\courssub{Le problème posé : une condensation non sélective}

\textbf{Intuition.} Imagine que tu mélanges dans un bécher des billes rouges (glycine, notée Gly) et des billes bleues (alanine, notée Ala), chaque bille possédant un « crochet » (la fonction \ce{-NH2}) et un « œillet » (la fonction \ce{-COOH}). Si tu agites le mélange, les crochets et les œillets s'accrochent \emph{au hasard} : tu obtiens toutes les paires possibles, pas seulement celle que tu visais.

\textbf{Explication rigoureuse.} Chaque acide $\alpha$-aminé possède \emph{deux} fonctions réactives : une fonction amine (nucléophile) et une fonction acide carboxylique (électrophile). Lorsqu'on mélange la glycine \ce{H2N-CH2-COOH} et l'alanine \ce{H2N-CH(CH3)-COOH} dans les conditions de condensation, quatre couples de réactifs sont possibles :

\begin{itemize}
\item le \ce{-COOH} d'une glycine réagit avec le \ce{-NH2} d'une autre glycine : dipeptide \textbf{Gly--Gly} ;
\item le \ce{-COOH} d'une alanine réagit avec le \ce{-NH2} d'une autre alanine : dipeptide \textbf{Ala--Ala} ;
\item le \ce{-COOH} de la glycine réagit avec le \ce{-NH2} de l'alanine : dipeptide \textbf{Gly--Ala} ;
\item le \ce{-COOH} de l'alanine réagit avec le \ce{-NH2} de la glycine : dipeptide \textbf{Ala--Gly}.
\end{itemize}

\begin{attention}[Quatre dipeptides sans protection]
Sans précaution particulière, la condensation d'un mélange de \textbf{deux} acides $\alpha$-aminés différents A et B conduit à un mélange de \textbf{quatre dipeptides} : A--A, B--B, A--B et B--A. De manière générale, avec $n$ acides aminés différents, on obtient $n^2$ dipeptides. Ce résultat est \textbf{très demandé au Bac}. De plus, la réaction ne s'arrête pas aux dipeptides : des tripeptides, tétrapeptides\dots{} se forment aussi, ce qui aggrave encore le mélange.
\end{attention}

\begin{popfigure}
\centering
\begin{tikzpicture}[scale=0.82, transform shape]
% Gly-Gly
\node[font=\small\bfseries, popblue] at (0,2.3) {Gly--Gly};
\node at (0,1.2) {\chemfig{H_2N-CH_2-C(=[1]O)-[7]NH-CH_2-COOH}};
% Ala-Ala
\node[font=\small\bfseries, poporange] at (8.2,2.3) {Ala--Ala};
\node at (8.2,1.2) {\chemfig{H_2N-CH(-[2]CH_3)-C(=[1]O)-[7]NH-CH(-[7]CH_3)-COOH}};
% Gly-Ala
\node[font=\small\bfseries, popgreen] at (0,-1.3) {Gly--Ala};
\node at (0,-2.4) {\chemfig{H_2N-CH_2-C(=[1]O)-[7]NH-CH(-[7]CH_3)-COOH}};
% Ala-Gly
\node[font=\small\bfseries, poppurple] at (8.2,-1.3) {Ala--Gly};
\node at (8.2,-2.4) {\chemfig{H_2N-CH(-[2]CH_3)-C(=[1]O)-[7]NH-CH_2-COOH}};
\end{tikzpicture}
\legende{Les quatre dipeptides obtenus par condensation directe d'un mélange glycine + alanine, sans protection des fonctions. La liaison peptidique est le groupe \ce{-CO-NH-} central.}
\end{popfigure}

\courssub{La stratégie générale : quatre étapes}

Pour n'obtenir \emph{qu'un seul} dipeptide, dans un ordre précis (par exemple Gly--Ala et non Ala--Gly), le chimiste procède comme un serrurier : il \emph{verrouille} les portes par lesquelles la réaction ne doit pas passer, et il \emph{huile} celle par laquelle elle doit passer.

\begin{definition}[Synthèse sélective d'un dipeptide]
La synthèse sélective d'un dipeptide A--B s'effectue en quatre étapes :
\begin{enumerate}
\item \textbf{Blocage} (protection) : on protège la fonction \ce{-NH2} de l'acide aminé A qui doit réagir par son \ce{-COOH} \textbf{et} le \ce{-COOH} de B ;
\item \textbf{Activation} : on rend la fonction \ce{-COOH} de A plus réactive ;
\item \textbf{Condensation} : on fait réagir le \ce{-COOH} activé de A avec le \ce{-NH2} libre de B : formation du dipeptide \emph{protégé} ;
\item \textbf{Déblocage} : on régénère les fonctions protégées sans rompre la liaison peptidique.
\end{enumerate}
\end{definition}

\begin{methode}[Réussir une synthèse peptidique sélective]
\begin{enumerate}
\item Identifier le dipeptide cible et son \emph{ordre} (ex. Gly--Ala : la glycine fournit le \ce{-COOH}, l'alanine fournit le \ce{-NH2}).
\item \textbf{Bloquer} la fonction qui ne doit \emph{pas} réagir : le \ce{-NH2} de la glycine (et le \ce{-COOH} de l'alanine).
\item \textbf{Activer} la fonction qui doit réagir : le \ce{-COOH} de la glycine.
\item \textbf{Condenser} : mélanger les deux partenaires ; une seule liaison peptidique peut se former.
\item \textbf{Débloquer} dans des conditions ménagées pour retrouver les fonctions libres.
\end{enumerate}
\end{methode}

\courssub{Étape 1 : le blocage (protection) des fonctions}

\begin{definition}[Blocage d'une fonction]
\cle{Bloquer} (ou protéger) une fonction chimique, c'est la transformer temporairement en un groupe \emph{inerte} dans les conditions de la condensation, de façon à empêcher sa participation à la réaction. Le groupe protecteur doit pouvoir être retiré ensuite, dans des conditions douces.
\end{definition}

\textbf{Blocage de la fonction amine.} On transforme \ce{-NH2} en un groupe \textbf{amide} (ou carbamate) non nucléophile. En synthèse moderne, on utilise couramment le groupe protecteur \emph{tert-butyloxycarbonyle} (Boc), introduit par l'anhydride \ce{(Boc)2O} en milieu basique. À titre d'illustration simple, on peut aussi utiliser un chlorure d'acyle comme le chlorure d'éthanoyle \ce{CH3-COCl} :
\[ \ce{H2N-CH2-COOH + CH3-COCl -> CH3-CO-NH-CH2-COOH + HCl} \]
L'azote, désormais engagé dans un groupe amide, n'a plus de doublet disponible : il ne peut plus attaquer un \ce{-COOH}.

\textbf{Blocage de la fonction acide carboxylique.} On la transforme en \textbf{ester}, par exemple par estérification de Fischer avec le méthanol en présence d'acide sulfurique concentré (catalyseur), à chaud :
\[ \ce{H2N-CH(CH3)-COOH + CH3OH <=>[H+] H2N-CH(CH3)-COOCH3 + H2O} \]
L'ester est beaucoup moins électrophile qu'un acide activé : il ne participe pas à la condensation.

\begin{attention}[Ne pas confondre blocage et activation]
Un \textbf{ester} sert ici de \emph{protection} du \ce{-COOH} de l'alanine (peu réactif), alors qu'un \textbf{chlorure d'acyle} sert d'\emph{activation} du \ce{-COOH} de la glycine (très réactif). Même famille de dérivés d'acides carboxyliques, mais rôles opposés : tout dépend de la réactivité du dérivé choisi.
\end{attention}

\courssub{Étape 2 : l'activation de la fonction acide carboxylique}

La fonction \ce{-COOH} réagit mal avec les amines : la condensation directe est lente, réversible et peu rentable. Il faut donc \cle{activer} la fonction, c'est-à-dire la transformer en un dérivé \emph{plus électrophile}.

\begin{definition}[Activation d'un acide carboxylique]
Activer une fonction acide carboxylique, c'est la convertir en un dérivé plus réactif : un \textbf{chlorure d'acyle} \ce{R-COCl} (ou un anhydride d'acide, ou un ester activé). Le chlorure d'acyle est obtenu par action du chlorure de thionyle \ce{SOCl2} :
\[ \ce{R-COOH + SOCl2 -> R-COCl + SO2 ^ + HCl ^} \]
Réaction totale, rapide, à température modérée ; les sous-produits \ce{SO2} et \ce{HCl} sont gazeux, ce qui simplifie la purification.
\end{definition}

\courssub{Étape 3 : la condensation}

On mélange alors les deux partenaires : l'acide aminé A dont le \ce{-NH2} est bloqué et le \ce{-COOH} activé, et l'acide aminé B dont le \ce{-NH2} est libre et le \ce{-COOH} bloqué. \emph{Une seule} combinaison réactive est possible : le \ce{-NH2} libre de B attaque le \ce{-COCl} de A. Il se forme le dipeptide \textbf{protégé} :
\[ \ce{Boc-NH-CH2-COCl + H2N-CH(CH3)-COOCH3 -> Boc-NH-CH2-CO-NH-CH(CH3)-COOCH3 + HCl} \]
La liaison peptidique \ce{-CO-NH-} est créée, dans l'ordre voulu (Gly avant Ala), et aucune autre liaison ne peut se former puisque toutes les autres fonctions sont verrouillées.

\courssub{Étape 4 : le déblocage}

\begin{definition}[Déblocage]
Le \cle{déblocage} (ou déprotection) consiste à régénérer les fonctions \ce{-NH2} et \ce{-COOH} initiales, par hydrolyse ménagée ou par action d'un réactif spécifique (par exemple un acide fort comme l'acide trifluoroacétique pour le groupe Boc, une saponification douce pour l'ester), \textbf{sans rompre la liaison peptidique} fraîchement formée.
\end{definition}

\begin{attention}[Un déblocage ménagé]
La liaison peptidique est elle-même une liaison \emph{amide}, hydrolysable en milieu acide ou basique concentré et à chaud. Si le déblocage est effectué dans des conditions trop brutales (acide concentré, ébullition prolongée), on hydrolyse aussi la liaison \ce{-CO-NH-} du dipeptide : tout le travail de synthèse est perdu ! On choisit donc des groupes protecteurs « clivables » dans des conditions douces, sélectives (ex. : Boc clivé par TFA à froid, ester méthylique saponifié par \ce{LiOH} à \SI{0}{\celsius}).
\end{attention}

\begin{exemplebox}[Synthèse complète de Gly--Ala]
Pour obtenir \emph{uniquement} le dipeptide Gly--Ala :
\begin{enumerate}
\item on \textbf{bloque} le \ce{-NH2} de la glycine (groupement Boc) et le \ce{-COOH} de l'alanine (ester méthylique) ;
\item on \textbf{active} le \ce{-COOH} de la glycine protégée en chlorure d'acyle (\ce{SOCl2}) ;
\item on \textbf{condense} : \ce{Boc-NH-CH2-COCl + H2N-CH(CH3)-COOCH3} $\longrightarrow$ dipeptide protégé ;
\item on \textbf{débloque} : hydrolyse douce de l'ester (saponification) et clivage acide ménagé du Boc $\longrightarrow$ \ce{H2N-CH2-CO-NH-CH(CH3)-COOH}, le seul Gly--Ala.
\end{enumerate}
Pour obtenir Ala--Gly, il suffit d'inverser les rôles : bloquer le \ce{-NH2} de l'alanine, activer son \ce{-COOH}, et utiliser la glycine à \ce{-NH2} libre.
\end{exemplebox}

\begin{popfigure}
\centering
\begin{tikzpicture}[node distance=0.55cm,
  etape/.style={draw=popblue, thick, fill=popblueL, rounded corners, text width=3.1cm, minimum height=2.6cm, align=center, font=\small},
  fleche/.style={-{Stealth[length=3mm]}, very thick, poporange}]
\node[etape] (e1) {\textbf{1. Blocage}\\[2pt] \ce{-NH2} de Gly protégé (Boc)\\ \ce{-COOH} de Ala estérifié};
\node[etape, right=of e1, fill=poporangeL, draw=poporange] (e2) {\textbf{2. Activation}\\[2pt] \ce{-COOH} de Gly $\to$ \ce{-COCl} (\ce{SOCl2})};
\node[etape, right=of e2, fill=popgreenL, draw=popgreen] (e3) {\textbf{3. Condensation}\\[2pt] \ce{-COCl} + \ce{H2N}-Ala $\to$ dipeptide protégé};
\node[etape, right=of e3, fill=poppurpleL, draw=poppurple] (e4) {\textbf{4. Déblocage}\\[2pt] conditions douces $\to$ \textbf{Gly--Ala} pur};
\draw[fleche] (e1) -- (e2);
\draw[fleche] (e2) -- (e3);
\draw[fleche] (e3) -- (e4);
\end{tikzpicture}
\legende{Les quatre étapes de la synthèse sélective du dipeptide Gly--Ala : bloquer, activer, condenser, débloquer.}
\end{popfigure}

\begin{exoresolu}[Combien de dipeptides ?]
On réalise la condensation, \textbf{sans aucune protection}, d'un mélange équimolaire de trois acides $\alpha$-aminés : glycine (Gly), alanine (Ala) et valine (Val). Combien de dipeptides différents peut-on obtenir ? Les citer.
\tcblower
\textbf{Solution.} Chaque dipeptide est formé d'un résidu N-terminal (celui qui fournit \ce{-COOH}) et d'un résidu C-terminal (celui qui fournit \ce{-NH2}). Il y a 3 choix pour le premier résidu et 3 choix pour le second, soit $3 \times 3 = 3^2 = \textbf{9 dipeptides}$ :
\begin{center}
Gly--Gly, Gly--Ala, Gly--Val, Ala--Gly, Ala--Ala, Ala--Val, Val--Gly, Val--Ala, Val--Val.
\end{center}
\emph{Remarque :} Gly--Ala et Ala--Gly sont deux composés \textbf{différents} (l'ordre des résidus compte). Pour n'obtenir qu'un seul de ces neuf dipeptides, il faudrait bloquer les fonctions qui ne doivent pas réagir et activer celle qui doit réagir.
\end{exoresolu}

\begin{savaistu}[L'aspartame, un dipeptide dans ta boisson]
L'\textbf{aspartame}, édulcorant intense (environ 200 fois plus sucré que le saccharose) présent dans de nombreuses boissons « light » vendues à Douala ou Yaoundé, est un dérivé du dipeptide \textbf{Asp--Phe} (acide aspartique -- phénylalanine), dont la fonction \ce{-COOH} terminale est estérifiée par un groupe méthyle. Sa fabrication industrielle illustre parfaitement la synthèse peptidique sélective : il a fallu protéger la seconde fonction acide de l'acide aspartique pour n'obtenir que l'enchaînement désiré ! Les personnes atteintes de phénylcétonurie ne peuvent pas métaboliser la phénylalanine : d'où la mention obligatoire « contient une source de phénylalanine » sur les étiquettes.
\end{savaistu}

\begin{aretenir}[L'essentiel de la section]
\begin{itemize}
\item Sans protection, deux acides $\alpha$-aminés A et B donnent \textbf{quatre dipeptides} : A--A, B--B, A--B, B--A (et $n^2$ dipeptides pour $n$ acides aminés).
\item La synthèse sélective suit quatre étapes : \textbf{blocage} de la fonction qui ne doit pas réagir, \textbf{activation} de celle qui doit réagir (chlorure d'acyle), \textbf{condensation}, puis \textbf{déblocage} ménagé.
\item L'ordre des résidus compte : Gly--Ala $\neq$ Ala--Gly ; c'est le choix des protections qui fixe l'ordre.
\item Le déblocage doit être suffisamment doux pour ne pas hydrolyser la liaison peptidique.
\end{itemize}
\end{aretenir}

\newpage

\coursec{Travaux pratiques}

\courssub{Travaux pratiques : Détection et caractère amphotère des acides $\alpha$-aminés}

\textbf{Objectifs}\par
À l'issue de cette séance de travaux pratiques, tu dois être capable de :

\begin{itemize}
\item mettre en évidence la présence d'\cle{acides $\alpha$-aminés} dans des extraits alimentaires courants (blanc d'œuf, lait de soja) à l'aide du \cle{test à la ninhydrine} (ou du test du biuret sur les protéines) ;
\item vérifier expérimentalement le \cle{caractère amphotère} de la glycine, c'est-à-dire la double capacité de l'\cle{amphion} à capter un proton (rôle de base) et à en céder un (rôle d'acide) ;
\item suivre une évolution de pH au papier pH et interpréter la \emph{résistance aux variations de pH} (effet tampon) ;
\item écrire les deux équations-bilans traduisant le caractère acide et le caractère basique de l'amphion.
\end{itemize}

\begin{savaistu}[Pourquoi ce TP dans un lycée camerounais ?]
Le soja, le niébé, l'arachide et le poisson fumé sont des sources majeures de protéines dans l'alimentation camerounaise. Or les protéines sont des assemblages d'acides $\alpha$-aminés. Ce TP te permet donc de relier directement le cours à ton assiette : identifier les acides aminés, c'est comprendre la valeur nutritionnelle des aliments.
\end{savaistu}

\textbf{Matériel et produits}\par

\begin{center}
\begin{tabularx}{\linewidth}{|l|X|}
\hline
\rowcolor{popblueL} \textbf{Matériel} & \textbf{Produits} \\
\hline
6 tubes à essais + portoir & Solution de ninhydrine à \SI{0,2}{\%} dans l'éthanol (à défaut : réactif de Fehling ou solution de sulfate de cuivre + soude pour le test du biuret) \\
\hline
2 béchers de \SI{100}{mL} & Blanc d'œuf frais dilué 10 fois dans l'eau \\
\hline
1 pipette jaugée de \SI{10}{mL} + poire & Lait de soja (jus de graines de soja trempées et mixées, filtré) \\
\hline
1 agitateur en verre & Glycine en poudre (pharmacie) ou gélatine alimentaire hydrolysée \\
\hline
Papier pH (ou pH-mètre si disponible) & Acide chlorhydrique dilué à \SI{0,1}{mol.L^{-1}} \\
\hline
2 compte-gouttes ou pipettes Pasteur & Soude diluée à \SI{0,1}{mol.L^{-1}} \\
\hline
Bain-marie (casserole + plaque ou réchaud) & Eau distillée (ou eau de pluie bouillie et refroidie) \\
\hline
Balance au centième (ou cuillère doseuse) & Lunettes, gants, blouse \\
\hline
\end{tabularx}
\end{center}

\textbf{Alternatives locales.} Si la ninhydrine est introuvable, réalise le \cle{test du biuret} : quelques gouttes de solution de sulfate de cuivre \ce{CuSO4} puis de soude concentrée sur l'échantillon ; une coloration \emph{violette} indique les liaisons peptidiques des protéines. Si la glycine manque, fais bouillir de la gélatine alimentaire dans l'eau acidifiée pendant 20 min : l'hydrolyse partielle libère des acides aminés. Le jus de soja se prépare en trempant les graines une nuit, en les mixant avec de l'eau puis en filtrant au tamis propre.

\textbf{Sécurité}\par

\begin{attention}[Consignes de sécurité]
\begin{itemize}
\item \textbf{Ninhydrine} : produit \emph{irritant} (pictogramme « point d'exclamation ») et inflammable car en solution alcoolique (pictogramme « flamme »). Manipuler loin de toute flamme, sous le ventilateur ou près d'une fenêtre ouverte, avec gants.
\item \textbf{Soude \ce{NaOH}} : \emph{corrosive} (pictogramme « corrosion »). En cas de contact avec la peau, rincer abondamment à l'eau claire pendant plusieurs minutes et prévenir l'enseignant.
\item \textbf{Acide chlorhydrique \ce{HCl}} : \emph{corrosif} également. Verser toujours l'acide \emph{dans} l'eau lors des dilutions, jamais l'inverse.
\item Port obligatoire de la \textbf{blouse}, des \textbf{lunettes} et des \textbf{gants}. Ne jamais goûter les échantillons, même alimentaires, une fois introduits au laboratoire.
\item Le bain-marie est chaud : manipuler les tubes avec une pince en bois.
\end{itemize}
\end{attention}

\textbf{Schéma du dispositif}\par

\begin{popfigure}
\begin{center}
\begin{tikzpicture}[scale=0.95, every node/.style={font=\small}]
% Tube à essais 1 : test ninhydrine
\draw[popdark, thick] (0,0) -- (0,2.6) arc[start angle=180, end angle=360, radius=0.4] -- (0.8,0);
\draw[popdark, thick] (0,2.6) -- (0,3.0) (0.8,2.6) -- (0.8,3.0);
\fill[poppurple!45] (0.05,0.35) arc[start angle=180, end angle=360, radius=0.35] -- (0.75,1.6) -- (0.05,1.6) -- cycle;
\draw[poppurple, thick] (0.05,1.6) -- (0.75,1.6);
\node[below, align=center] at (0.4,-0.25) {Tube 1 : blanc d'œuf\\ + ninhydrine (violet)};
% Tube 2 : témoin
\draw[popdark, thick] (2.2,0) -- (2.2,2.6) arc[start angle=180, end angle=360, radius=0.4] -- (3.0,0);
\draw[popdark, thick] (2.2,2.6) -- (2.2,3.0) (3.0,2.6) -- (3.0,3.0);
\fill[popblue!20] (2.25,0.35) arc[start angle=180, end angle=360, radius=0.35] -- (2.95,1.6) -- (2.25,1.6) -- cycle;
\draw[popblue, thick] (2.25,1.6) -- (2.95,1.6);
\node[below, align=center] at (2.6,-0.25) {Tube témoin : eau\\ + ninhydrine (incolore)};
% Bécher avec solution de glycine
\draw[popdark, thick] (5.0,0) -- (5.0,2.8) (7.4,0) -- (7.4,2.8);
\draw[popdark, thick] (5.0,0) -- (7.4,0);
\fill[popgreen!25] (5.05,0.05) rectangle (7.35,1.7);
\draw[popgreen, thick] (5.05,1.7) -- (7.35,1.7);
\node[below, align=center] at (6.2,-0.25) {Bécher : solution de glycine\\ suivi du pH};
% agitateur
\draw[popdark, thick] (6.9,3.2) -- (6.5,0.6);
\node[right, font=\scriptsize] at (6.9,3.2) {agitateur};
% bande papier pH
\draw[popgold!60, fill=popgold!30] (4.2,2.2) rectangle (4.9,2.5);
\node[font=\scriptsize, below] at (4.55,2.2) {papier pH};
\draw[-{Stealth}, popdark] (4.9,2.35) -- (5.4,1.9);
% compte-gouttes HCl
\draw[popdark] (5.6,3.6) ellipse (0.18 and 0.3);
\draw[popdark, thick] (5.6,3.3) -- (5.6,2.2);
\draw[popdark, thick] (5.5,2.2) -- (5.6,1.9) -- (5.7,2.2);
\node[left, font=\scriptsize, align=right] at (5.3,3.6) {\ce{HCl} \SI{0,1}{mol.L^{-1}}\\ goutte à goutte};
% compte-gouttes NaOH
\draw[popdark] (7.0,3.6) ellipse (0.18 and 0.3);
\draw[popdark, thick] (7.0,3.3) -- (7.0,2.2);
\draw[popdark, thick] (6.9,2.2) -- (7.0,1.9) -- (7.1,2.2);
\node[right, font=\scriptsize, align=left] at (7.3,3.6) {\ce{NaOH} \SI{0,1}{mol.L^{-1}}\\ goutte à goutte};
\end{tikzpicture}
\end{center}
\legende{Dispositif du TP : à gauche, les tubes du test à la ninhydrine (échantillon et témoin) ; à droite, le suivi pH-métrique qualitatif de la solution de glycine lors de l'ajout d'acide puis de base.}
\end{popfigure}

\textbf{Protocole}\par

\textbf{Partie A — Mise en évidence des acides $\alpha$-aminés (test à la ninhydrine)}

\begin{enumerate}
\item Numéroter trois tubes à essais : tube 1 (blanc d'œuf dilué), tube 2 (jus de soja), tube 3 (eau distillée, témoin).
\item Introduire dans chaque tube \SI{2}{mL} de l'échantillon correspondant.
\item Ajouter dans chaque tube 5 gouttes de solution de ninhydrine. Agiter.
\item Plonger les trois tubes dans le bain-marie à environ \SI{80}{\celsius} pendant \SI{5}{min}.
\item Sortir les tubes avec la pince en bois, laisser refroidir et comparer les colorations à l'œil nu.
\end{enumerate}

\textbf{Partie B — Caractère amphotère de la glycine}

\begin{enumerate}
\item Peser environ \SI{0,75}{g} de glycine (masse molaire \SI{75,07}{g.mol^{-1}}, soit une quantité de matière $n \approx \SI{0,010}{mol}$) et la dissoudre dans \SI{50}{mL} d'eau distillée dans le bécher. La concentration de la solution est d'environ \SI{0,20}{mol.L^{-1}}. Agiter jusqu'à dissolution complète.
\item Déposer une goutte de solution sur le papier pH (ou plonger l'électrode du pH-mètre) : noter le pH initial (attendu vers 6,0, point isoélectrique de la glycine).
\item Ajouter \ce{HCl} à \SI{0,1}{mol.L^{-1}} goutte à goutte en agitant. Relever le pH après 5, 10, 15 et 20 gouttes.
\item Rincer le bécher, préparer une nouvelle solution identique de glycine, puis ajouter cette fois la soude à \SI{0,1}{mol.L^{-1}} goutte à goutte. Relever le pH après 5, 10, 15 et 20 gouttes.
\item Comparer avec un essai témoin : mêmes ajouts dans \SI{50}{mL} d'eau distillée pure.
\end{enumerate}

\textbf{Observations et résultats attendus}\par

\textbf{Partie A.}

\begin{center}
\begin{tabularx}{\linewidth}{|l|X|X|}
\hline
\rowcolor{popblueL} \textbf{Tube} & \textbf{Coloration observée} & \textbf{Conclusion} \\
\hline
1 : blanc d'œuf + ninhydrine & Violette intense (bleu-violet de Ruhemann) & Présence d'acides $\alpha$-aminés / protéines \\
\hline
2 : jus de soja + ninhydrine & Violette (parfois rosée) & Présence d'acides $\alpha$-aminés \\
\hline
3 : eau + ninhydrine & Incolore à jaune pâle & Témoin négatif : pas d'acides aminés \\
\hline
\end{tabularx}
\end{center}

\textbf{Partie B.} Valeurs réalistes attendues (papier pH, précision $\pm 0,5$) pour une solution de glycine \SI{0,20}{mol.L^{-1}} :

\begin{center}
\begin{tabularx}{\linewidth}{|l|X|X|X|X|}
\hline
\rowcolor{popblueL} \textbf{Gouttes ajoutées} & 0 & 5 & 10 & 20 \\
\hline
Glycine + \ce{HCl} (pH) & 6,0 & 5,5 & 5,0 & 4,0 \\
\hline
Eau pure + \ce{HCl} (pH) & 7,0 & 4,0 & 3,0 & 2,0 \\
\hline
Glycine + \ce{NaOH} (pH) & 6,0 & 6,5 & 7,0 & 8,0 \\
\hline
Eau pure + \ce{NaOH} (pH) & 7,0 & 10,0 & 11,0 & 12,0 \\
\hline
\end{tabularx}
\end{center}

On constate que le pH de la solution de glycine varie \emph{beaucoup moins} que celui de l'eau pure, dans les deux sens : la solution de glycine \textbf{résiste aux variations de pH}.

\textbf{Exploitation}\par

\begin{exoresolu}[Questions d'exploitation]
\textbf{1.} Pourquoi le tube 3 (eau) est-il indispensable ? \par
\textbf{2.} La glycine est en solution sous forme d'amphion \ce{^{+}H3N-CH2-COO-}. Écrire l'équation-bilan de sa réaction avec les ions \ce{H3O+} apportés par l'acide chlorhydrique. Quel rôle (acide ou base) joue l'amphion ? Préciser le couple acide/base concerné. \par
\textbf{3.} Écrire l'équation-bilan de la réaction de l'amphion avec les ions \ce{HO-} apportés par la soude. Quel rôle joue alors l'amphion ? Préciser le couple acide/base concerné. \par
\textbf{4.} Expliquer pourquoi le pH varie peu dans les deux cas et conclure sur le caractère amphotère de la glycine.
\tcblower
\textbf{1.} Le tube témoin permet de vérifier que la coloration violette ne provient ni du solvant ni du réactif seul : seule la présence d'acides $\alpha$-aminés (ou de protéines) provoque le virage au violet (formation du colorant de Ruhemann). Sans témoin, aucune conclusion valable n'est possible. \par
\textbf{2.} L'ion carboxylate \ce{-COO-} de l'amphion capte un proton (réaction de neutralisation) :
\[ \ce{^{+}H3N-CH2-COO- + H3O+ -> ^{+}H3N-CH2-COOH + H2O} \]
L'amphion \emph{accepte} un proton : il joue le rôle de \textbf{base} (couple \ce{^{+}H3N-CH2-COOH}/\ce{^{+}H3N-CH2-COO-}, pK$_{\text{a}1} \approx 2,34$). \par
\textbf{3.} Le groupe ammonium \ce{-NH3+} de l'amphion cède un proton :
\[ \ce{^{+}H3N-CH2-COO- + HO- -> H2N-CH2-COO- + H2O} \]
L'amphion \emph{cède} un proton : il joue le rôle d'\textbf{acide} (couple \ce{^{+}H3N-CH2-COO-}/\ce{H2N-CH2-COO-}, pK$_{\text{a}2} \approx 9,60$). \par
\textbf{4.} Les ions \ce{H3O+} ajoutés sont consommés par la réaction de la question 2 (l'amphion agit comme base), et les ions \ce{HO-} ajoutés sont consommés par la réaction de la question 3 (l'amphion agit comme acide) : leur concentration dans la solution varie donc peu, d'où la faible variation de pH (effet tampon). La glycine réagit à la fois comme un acide et comme une base : elle est \textbf{amphotère}, ce qui confirme le comportement de l'amphion en solution aqueuse.
\end{exoresolu}

\textbf{Conclusion}\par

Ce TP a permis de vérifier deux résultats essentiels du cours. D'une part, le test à la ninhydrine (coloration violette) révèle la présence d'acides $\alpha$-aminés libres ou engagés dans des protéines dans des aliments courants comme le blanc d'œuf et le jus de soja. D'autre part, la solution de glycine résiste aux variations de pH aussi bien lors de l'ajout d'acide que lors de l'ajout de base : l'\cle{amphion} \ce{^{+}H3N-CH2-COO-} se comporte à la fois comme une base (il capte \ce{H+} sur son groupe \ce{-COO-}, couple \ce{^{+}H3N-CH2-COOH}/\ce{^{+}H3N-CH2-COO-}) et comme un acide (il cède \ce{H+} depuis son groupe \ce{-NH3+}, couple \ce{^{+}H3N-CH2-COO-}/\ce{H2N-CH2-COO-}). Cette double réactivité, caractéristique des acides $\alpha$-aminés, explique le rôle \emph{tampon} des protéines dans les milieux biologiques, par exemple dans le sang (hémoglobine) ou dans les cellules.

\newpage

\coursec{Méthodes et exercices résolus}

\coursec{Les acides $\alpha$-aminés et la synthèse des peptides}

\courssub{Méthode 1 : Reconnaître et écrire la formule d'un acide $\alpha$-aminé}

\begin{methode}[Identifier un acide $\alpha$-aminé]
\begin{enumerate}
\item Repérer la fonction \cle{acide carboxylique} \ce{-COOH} : c'est elle qui fixe la numérotation de la chaîne (le carbone du \ce{-COOH} porte le numéro 1).
\item Repérer la fonction \cle{amine} \ce{-NH2}.
\item Vérifier que l'amine est portée par le carbone en \cle{$\alpha$}, c'est-à-dire le carbone \textbf{n° 2}, voisin du \ce{-COOH}.
\item Écrire la formule semi-développée sous la forme générale \chemfig{R-CH(-[2]NH_2)-COOH} où R est la chaîne latérale (R = H pour la glycine).
\item Nommer : on nomme l'acide carboxylique correspondant, puis on précise la position de l'amine par le préfixe « amino- » (ex. : acide 2-aminopropanoïque).
\end{enumerate}
\textbf{Réflexe Bac :} si l'amine n'est pas sur le carbone $\alpha$, ce n'est \emph{pas} un acide $\alpha$-aminé (ex. : \ce{H2N-CH2-CH2-COOH} est un acide $\beta$-aminé).
\end{methode}

\begin{exoresolu}[Acides aminés ou pas ?]
On considère les composés suivants :
\textbf{A} : \ce{CH3-CH(NH2)-COOH} ; \textbf{B} : \ce{H2N-CH2-CH2-COOH} ; \textbf{C} : \ce{CH3-CH2-CH(NH2)-COOH}.
\begin{enumerate}
\item Pour chaque composé, dire s'il s'agit d'un acide $\alpha$-aminé. Justifier.
\item Donner le nom systématique des acides $\alpha$-aminés identifiés.
\end{enumerate}
\tcblower
\textbf{1. Analyse de chaque composé.}

\textbf{Composé A :} \ce{CH3-CH(NH2)-COOH}. Le carbone n° 1 est celui du \ce{-COOH} ; le carbone n° 2 (le \ce{CH}) porte le groupe \ce{-NH2}. L'amine est en position $\alpha$ : \textbf{A est un acide $\alpha$-aminé} (c'est l'alanine, R = \ce{CH3}).

\textbf{Composé B :} \ce{H2N-CH2-CH2-COOH}. Le carbone $\alpha$ (n° 2) est le \ce{CH2} voisin du \ce{-COOH} : il ne porte \emph{pas} l'amine. L'amine est sur le carbone n° 3 (position $\beta$). \textbf{B n'est pas un acide $\alpha$-aminé} ; c'est un acide $\beta$-aminé (acide 3-aminopropanoïque, ou $\beta$-alanine).

\textbf{Composé C :} \ce{CH3-CH2-CH(NH2)-COOH}. Le carbone n° 2 porte le \ce{-NH2} : \textbf{C est un acide $\alpha$-aminé} (R = \ce{CH3-CH2}-, c'est l'acide 2-aminobutanoïque, aussi appelé $\alpha$-aminobutyrique).

\textbf{2. Noms systématiques.}
\begin{itemize}
\item A : acide \textbf{2-aminopropanoïque} (chaîne à 3 carbones : propane ; acide : propanoïque ; amine en 2).
\item C : acide \textbf{2-aminobutanoïque} (chaîne à 4 carbones).
\end{itemize}
\textbf{Conclusion :} A et C sont des acides $\alpha$-aminés ; B ne l'est pas car son amine est en position $\beta$.
\end{exoresolu}

\courssub{Méthode 2 : Carbone asymétrique et représentation de Fischer}

\begin{methode}[Repérer un carbone asymétrique et construire la projection de Fischer]
\begin{enumerate}
\item Un \cle{carbone asymétrique} est un carbone \textbf{tétraédrique} ($sp^3$) portant \textbf{quatre groupes différents}.
\item Dans un acide $\alpha$-aminé \chemfig{R-CH(-[2]NH_2)-COOH}, le carbone $\alpha$ porte : \ce{-COOH}, \ce{-NH2}, \ce{-H} et \ce{-R}. Il est asymétrique \textbf{sauf si R = H} (glycine, seul acide $\alpha$-aminé naturel \textbf{achiral}).
\item \textbf{Projection de Fischer :} placer la chaîne carbonée \emph{verticale}, le groupe le plus oxydé (\ce{-COOH}) \emph{en haut}, R en bas ; \ce{-NH2} et \ce{-H} sur l'horizontale. Convention : les liaisons horizontales sortent du plan vers le lecteur, les liaisons verticales rentrent dans le plan.
\item \textbf{Configuration :} \ce{-NH2} à \textbf{gauche} $\Rightarrow$ série \textbf{L} ; \ce{-NH2} à \textbf{droite} $\Rightarrow$ série \textbf{D}. Les acides aminés naturels des protéines sont de configuration \textbf{L}.
\end{enumerate}
\textbf{Réflexe :} dans une projection de Fischer, on ne fait pivoter la molécule que de 180° dans le plan (jamais de 90°, qui changerait la configuration).
\end{methode}

\begin{exoresolu}[L'alanine est-elle chirale ?]
L'alanine a pour formule semi-développée \ce{CH3-CH(NH2)-COOH}.
\begin{enumerate}
\item Montrer que la molécule d'alanine possède un carbone asymétrique.
\item Donner les représentations de Fischer des deux énantiomères et préciser lequel correspond à la L-alanine, celle que l'on trouve dans les protéines.
\item La glycine \ce{H2N-CH2-COOH} est-elle chirale ? Justifier.
\end{enumerate}
\tcblower
\textbf{1. Carbone asymétrique.} Le carbone $\alpha$ (celui du \ce{CH}) est tétraédrique et porte quatre groupes tous différents : \ce{-COOH}, \ce{-NH2}, \ce{-H} et \ce{-CH3}. C'est donc un \textbf{carbone asymétrique}, noté C* : \ce{CH3-\overset{*}{C}H(NH2)-COOH}. La molécule d'alanine est chirale : elle existe sous deux énantiomères non superposables, images l'un de l'autre dans un miroir.

\textbf{2. Représentations de Fischer.} On place \ce{-COOH} en haut et \ce{-CH3} en bas.

\begin{popfigure}
\begin{tikzpicture}[scale=1.0]
% L-alanine
\draw[thick] (0,0.7) -- (0,-0.7);
\draw[thick] (-0.9,0) -- (0.9,0);
\node[above] at (0,0.7) {\ce{COOH}};
\node[below] at (0,-0.7) {\ce{CH3}};
\node[left] at (-0.9,0) {\ce{H2N}};
\node[right] at (0.9,0) {\ce{H}};
\fill (0,0) circle (1.5pt);
\node[below=1.0cm] at (0,-0.7) {\textbf{L-alanine}};
% D-alanine
\begin{scope}[xshift=5cm]
\draw[thick] (0,0.7) -- (0,-0.7);
\draw[thick] (-0.9,0) -- (0.9,0);
\node[above] at (0,0.7) {\ce{COOH}};
\node[below] at (0,-0.7) {\ce{CH3}};
\node[left] at (-0.9,0) {\ce{H}};
\node[right] at (0.9,0) {\ce{NH2}};
\fill (0,0) circle (1.5pt);
\node[below=1.0cm] at (0,-0.7) {\textbf{D-alanine}};
\end{scope}
\end{tikzpicture}
\legende{Projections de Fischer des deux énantiomères de l'alanine.}
\end{popfigure}

Dans la \textbf{L-alanine}, le groupe \ce{-NH2} est à \textbf{gauche} : c'est l'énantiomère présent dans les protéines naturelles. Dans la D-alanine, il est à droite.

\textbf{3. Cas de la glycine.} Dans \ce{H2N-CH2-COOH}, le carbone $\alpha$ porte \textbf{deux atomes d'hydrogène identiques} : il n'a pas quatre groupes différents. La glycine ne possède donc \textbf{pas de carbone asymétrique} : elle est \textbf{achirale} (non chirale). C'est le seul acide $\alpha$-aminé naturel dans ce cas.

\textbf{Conclusion :} l'alanine est chirale (un C*) et la forme naturelle est la L-alanine (\ce{NH2} à gauche en Fischer) ; la glycine est achirale.
\end{exoresolu}

\courssub{Méthode 3 : Propriétés acido-basiques de l'amphion}

\begin{methode}[Exploiter le caractère amphotère d'un acide $\alpha$-aminé]
\begin{enumerate}
\item En solution aqueuse, l'acide $\alpha$-aminé existe sous forme d'\cle{amphion} (zwitterion) : \ce{^{+}H3N-CHR-COO-}, issu du transfert interne de proton :
\[ \ce{H2N-CHR-COOH <=> ^{+}H3N-CHR-COO-} \]
\item L'amphion possède \textbf{deux couples acide/base} :
\begin{itemize}
\item couple \ce{-COOH}/\ce{-COO-} de l'amphion : \ce{^{+}H3N-CHR-COOH}/\ce{^{+}H3N-CHR-COO-} ($pK_{a1} \approx 2{,}3$) ;
\item couple \ce{-NH3+}/\ce{-NH2} : \ce{^{+}H3N-CHR-COO-}/\ce{H2N-CHR-COO-} ($pK_{a2} \approx 9{,}7$).
\end{itemize}
\item \textbf{Caractère acide} de l'amphion (il cède un proton à une base, ex. \ce{HO-}) :
\[ \ce{^{+}H3N-CHR-COO- + HO- -> H2N-CHR-COO- + H2O} \]
\item \textbf{Caractère basique} de l'amphion (il capte un proton d'un acide, ex. \ce{H3O+}) :
\[ \ce{^{+}H3N-CHR-COO- + H3O+ -> ^{+}H3N-CHR-COOH + H2O} \]
\item \textbf{Diagramme de prédominance :} cation \ce{^{+}H3N-CHR-COOH} pour $pH < pK_{a1}$ ; amphion pour $pK_{a1} < pH < pK_{a2}$ ; anion \ce{H2N-CHR-COO-} pour $pH > pK_{a2}$.
\end{enumerate}
\end{methode}

\begin{exoresolu}[L'amphion de la glycine, un amphotère]
On dissout de la glycine Gly ($pK_{a1} = 2{,}3$ ; $pK_{a2} = 9{,}6$) dans l'eau. Le pH de la solution obtenue est voisin de 6.
\begin{enumerate}
\item Écrire l'équation-bilan de l'équilibre entre la forme moléculaire et la forme ionique (amphion) de la glycine.
\item Sous quelle forme prédominante se trouve la glycine dans cette solution ? Justifier à l'aide d'un diagramme de prédominance.
\item On ajoute quelques gouttes d'une solution concentrée d'acide chlorhydrique. Écrire l'équation-bilan de la réaction qui se produit et dire quel rôle (acide ou base) joue l'amphion.
\item On ajoute ensuite une solution d'hydroxyde de sodium en excès. Écrire l'équation-bilan correspondante.
\end{enumerate}
\tcblower
\textbf{1. Équilibre forme moléculaire / forme ionique.} Il s'agit d'un transfert de proton \emph{intramoléculaire} du groupe \ce{-COOH} vers le groupe \ce{-NH2} :
\[ \ce{H2N-CH2-COOH <=> ^{+}H3N-CH2-COO-} \]
L'espèce \ce{^{+}H3N-CH2-COO-}, qui porte à la fois une charge positive et une charge négative, est l'\textbf{amphion} (ou zwitterion).

\textbf{2. Forme prédominante à pH = 6.} Traçons le diagramme de prédominance :

\begin{popfigure}
\begin{tikzpicture}[scale=1.0]
\draw[->,thick] (0,0) -- (11,0) node[below left]{pH};
\draw[thick] (2.3,0.15) -- (2.3,-0.15) node[below]{$pK_{a1}=2{,}3$};
\draw[thick] (9.6,0.15) -- (9.6,-0.15) node[below]{$pK_{a2}=9{,}6$};
\node[above,align=center] at (1.1,0.5) {\small cation\\ \small \ce{^{+}H3N-CH2-COOH}};
\node[above,align=center] at (5.9,0.5) {\small amphion\\ \small \ce{^{+}H3N-CH2-COO-}};
\node[above,align=center] at (10.3,0.5) {\small anion\\ \small \ce{H2N-CH2-COO-}};
\fill[popblue] (6,0) circle (2pt) node[below=2pt]{\small pH = 6};
\end{tikzpicture}
\legende{Diagramme de prédominance de la glycine.}
\end{popfigure}

Le point $pH = 6$ se situe entre $pK_{a1}$ et $pK_{a2}$ : la forme \textbf{prédominante est l'amphion} \ce{^{+}H3N-CH2-COO-}.

\textbf{3. Ajout d'acide chlorhydrique.} L'acide chlorhydrique libère des ions \ce{H3O+}. L'amphion capte un proton sur son groupe \ce{-COO-} :
\[ \ce{^{+}H3N-CH2-COO- + H3O+ -> ^{+}H3N-CH2-COOH + H2O} \]
L'amphion \textbf{capte} un proton : il joue le rôle de \textbf{base} (couple \ce{^{+}H3N-CH2-COOH}/\ce{^{+}H3N-CH2-COO-}).

\textbf{4. Ajout de soude en excès.} Les ions \ce{HO-} arrachent le proton du groupe \ce{-NH3+} :
\[ \ce{^{+}H3N-CH2-COO- + HO- -> H2N-CH2-COO- + H2O} \]
L'amphion \textbf{cède} un proton : il joue le rôle d'\textbf{acide} (couple \ce{^{+}H3N-CH2-COO-}/\ce{H2N-CH2-COO-}).

\textbf{Conclusion :} l'amphion est un \textbf{amphotère} (espèce amphotérique) : il réagit comme base avec les acides forts et comme acide avec les bases fortes.
\end{exoresolu}

\courssub{Méthode 4 : Écrire la formation d'un dipeptide et nommer un peptide}

\begin{methode}[Liaison peptidique et nomenclature des peptides]
\begin{enumerate}
\item La \cle{liaison peptidique} est une \textbf{liaison amide} \ce{-CO-NH-} formée entre le groupe \ce{-COOH} d'un acide aminé et le groupe \ce{-NH2} d'un autre, avec \textbf{élimination d'une molécule d'eau} (condensation).
\item Écrire l'équation en alignant \ce{HOOC-} et \ce{H2N-} en vis-à-vis, puis barrer \ce{-OH} du premier et un \ce{-H} du second.
\item \textbf{Nomenclature :} on lit le peptide depuis l'extrémité \textbf{N-terminale} (porte le \ce{-NH2} libre) vers l'extrémité \textbf{C-terminale} (porte le \ce{-COOH} libre). Chaque résidu, sauf le dernier, prend le suffixe « -yle » : Gly-Ala = glycylalanine.
\item \textbf{Attention à l'ordre :} Gly-Ala $\neq$ Ala-Gly : ce sont deux dipeptides \emph{différents}.
\item La liaison peptidique est \textbf{plane} (caractère partiel de double liaison \ce{C-N}) : les six atomes \ce{C, O, N, H} et les deux C$\alpha$ voisins sont coplanaires.
\end{enumerate}
\end{methode}

\begin{exoresolu}[Formation du dipeptide Gly-Ala]
On fait réagir la glycine \ce{H2N-CH2-COOH} (Gly) avec l'alanine \ce{CH3-CH(NH2)-COOH} (Ala) dans des conditions permettant la condensation.
\begin{enumerate}
\item Écrire l'équation-bilan de formation du dipeptide Gly-Ala.
\item Entourer la liaison peptidique sur la formule obtenue et nommer ses deux extrémités.
\item Écrire la formule semi-développée du dipeptide Ala-Gly. Est-il identique au précédent ?
\item On réalise la condensation d'un mélange équimolaire de glycine et d'alanine \emph{sans aucune protection}. Combien de dipeptides différents peut-on obtenir ? Les citer.
\end{enumerate}
\tcblower
\textbf{1. Équation-bilan de formation de Gly-Ala.} Le groupe \ce{-COOH} de la glycine réagit avec le groupe \ce{-NH2} de l'alanine :
\[ \ce{H2N-CH2-COOH + H2N-CH(CH3)-COOH -> H2N-CH2-CO-NH-CH(CH3)-COOH + H2O} \]

\textbf{2. Liaison peptidique et extrémités.} Le dipeptide obtenu s'écrit :
\begin{center}
\chemfig{H_2N-CH_2-C(=O)-NH-CH(-[2]CH_3)-COOH}
\end{center}
La liaison peptidique est le groupe \textbf{\ce{-CO-NH-}} (liaison amide). L'extrémité portant le \ce{-NH2} libre (côté glycine) est l'extrémité \textbf{N-terminale} ; celle portant le \ce{-COOH} libre (côté alanine) est l'extrémité \textbf{C-terminale}.

\textbf{3. Le dipeptide Ala-Gly.} Cette fois, c'est le \ce{-COOH} de l'alanine qui réagit avec le \ce{-NH2} de la glycine :
\[ \ce{H2N-CH(CH3)-CO-NH-CH2-COOH} \]
Ce composé est \textbf{différent} de Gly-Ala : le groupe \ce{-CH3} est porté par le résidu N-terminal au lieu du résidu C-terminal. L'ordre des résidus compte : \textbf{Gly-Ala et Ala-Gly sont deux dipeptides distincts}.

\textbf{4. Nombre de dipeptides sans protection.} Sans blocage des fonctions, chaque \ce{-COOH} peut réagir avec chaque \ce{-NH2} présent. Les condensations possibles sont :
\begin{itemize}
\item Gly + Gly $\rightarrow$ \textbf{Gly-Gly} ;
\item Ala + Ala $\rightarrow$ \textbf{Ala-Ala} ;
\item Gly(\ce{COOH}) + Ala(\ce{NH2}) $\rightarrow$ \textbf{Gly-Ala} ;
\item Ala(\ce{COOH}) + Gly(\ce{NH2}) $\rightarrow$ \textbf{Ala-Gly}.
\end{itemize}
On obtient donc un mélange de \textbf{4 dipeptides différents} (règle générale : $n^2$ dipeptides à partir de $n$ acides aminés, ici $2^2 = 4$).

\textbf{Conclusion :} la condensation non sélective de deux acides aminés conduit à un mélange de quatre dipeptides, d'où la nécessité d'une synthèse \emph{sélective} pour n'obtenir que Gly-Ala.
\end{exoresolu}

\courssub{Méthode 5 : Synthèse sélective d'un dipeptide}

\begin{methode}[Les quatre étapes de la synthèse sélective]
Pour synthétiser \textbf{uniquement} le dipeptide A-B (A en N-terminal, B en C-terminal) :
\begin{enumerate}
\item \textbf{Bloquer (protéger)} la fonction \ce{-NH2} de l'acide aminé A (celui dont on veut utiliser le \ce{-COOH}) pour l'empêcher de réagir.
\item \textbf{Activer} la fonction \ce{-COOH} de A (ou bloquer le \ce{-COOH} de B et activer le \ce{-NH2} de B) pour rendre la condensation possible et dirigée.
\item \textbf{Condensation} : faire réagir A protégé/activé avec B : seule la liaison \ce{COOH}(A)--\ce{NH2}(B) se forme.
\item \textbf{Débloquer (déprotéger)} les fonctions protégées pour retrouver le dipeptide libre A-B.
\end{enumerate}
\textbf{Réflexe Bac :} sans protection, $n$ acides aminés donnent $n^2$ dipeptides ; la stratégie « blocage -- activation -- condensation -- déblocage » permet d'en obtenir \emph{un seul}.
\end{methode}

\begin{exoresolu}[Synthèse sélective de Gly-Ala]
Un laboratoire pharmaceutique de Yaoundé veut préparer le dipeptide Gly-Ala de façon sélective à partir de glycine et d'alanine.
\begin{enumerate}
\item Expliquer pourquoi la simple mise en contact des deux acides aminés ne convient pas.
\item Décrire, en précisant le rôle de chaque étape, le protocole de synthèse sélective de Gly-Ala.
\item Indiquer sur quel acide aminé porte le blocage de la fonction amine. Justifier.
\end{enumerate}
\tcblower
\textbf{1. Inconvénient de la condensation directe.} La glycine et l'alanine possèdent chacune une fonction acide et une fonction amine, toutes réactives. La mise en contact directe conduit à un mélange de \textbf{quatre dipeptides} : Gly-Gly, Ala-Ala, Gly-Ala et Ala-Gly (voir exercice précédent). Le rendement en Gly-Ala serait faible et la séparation coûteuse : il faut donc \textbf{orienter} la réaction.

\textbf{2. Protocole de synthèse sélective.}
\begin{itemize}
\item \textbf{Étape 1 -- Blocage :} on protège la fonction \ce{-NH2} de la \textbf{glycine} (par exemple par un groupe protecteur comme le BOC ou le FMOC) : elle ne pourra plus s'auto-condenser ni réagir avec le \ce{-COOH} de l'alanine.
\item \textbf{Étape 2 -- Activation :} on active la fonction \ce{-COOH} de la glycine (transformation en dérivé plus réactif, ex. chlorure d'acyle, anhydride mixte ou utilisation d'un agent de couplage type DCC/EDC) pour faciliter l'attaque par l'amine de l'alanine. On peut aussi bloquer le \ce{-COOH} de l'alanine (ex. ester méthylique) pour éviter l'auto-condensation de celle-ci.
\item \textbf{Étape 3 -- Condensation :} on fait réagir la glycine protégée/activée avec l'alanine. Seule la liaison \ce{COOH}(Gly)--\ce{NH2}(Ala) se forme :
\[ \ce{Gly(NH2\ protege, COOH\ active) + Ala -> Gly-Ala(protege) + H2O} \]
\item \textbf{Étape 4 -- Déblocage :} on élimine les groupes protecteurs dans des conditions douces qui \textbf{n'hydrolysent pas la liaison peptidique} (ex. TFA pour BOC, pipéridine pour FMOC, saponification pour l'ester). On isole le dipeptide pur \textbf{Gly-Ala}.
\end{itemize}

\textbf{3. Choix de la fonction à bloquer.} Dans Gly-Ala, c'est le \ce{-COOH} de la glycine qui forme la liaison peptidique : la glycine doit donc réagir \emph{par son acide}. C'est donc sa fonction \textbf{amine} qu'il faut bloquer (sinon elle attaquerait un autre \ce{-COOH} et donnerait Gly-Gly ou Ala-Gly). Symétriquement, l'alanine doit réagir par son \ce{-NH2} : on peut bloquer son \ce{-COOH}.

\textbf{Conclusion :} la stratégie « blocage de \ce{NH2}(Gly) -- activation de \ce{COOH}(Gly) -- condensation -- déblocage » conduit exclusivement au dipeptide Gly-Ala.
\end{exoresolu}

\courssub{Méthode 6 : Exploiter un dosage ou une composition centésimale d'un acide aminé}

\begin{methode}[Déterminer une formule ou une quantité de matière]
\begin{enumerate}
\item \textbf{Données utiles :} masses molaires atomiques ($M(\ce{C}) = \SI{12}{g.mol^{-1}}$, $M(\ce{H}) = \SI{1}{g.mol^{-1}}$, $M(\ce{O}) = \SI{16}{g.mol^{-1}}$, $M(\ce{N}) = \SI{14}{g.mol^{-1}}$).
\item \textbf{Formule brute d'un acide $\alpha$-aminé saturé à une amine :} \ce{C_nH_{2n+1}NO2} (chaîne latérale saturée, acyclique).
\item Pour trouver $n$ : poser $M = 14n + 47$ (car $12n + (2n+1) + 32 + 14 = 14n + 47$) et résoudre.
\item \textbf{Quantité de matière :} $n = \dfrac{m}{M}$ ; \textbf{concentration :} $c = \dfrac{n}{V}$ ; toujours convertir les volumes en litres.
\item \textbf{Dosage de l'amphion :} l'amphion réagit \emph{mole à mole} avec \ce{HO-} ou avec \ce{H3O+} : à l'équivalence, $n(\text{amphion}) = n(\ce{HO-}) = c_B \times V_{BE}$.
\end{enumerate}
\end{methode}

\begin{exoresolu}[Dosage de la valine]
La valine (Val) est un acide $\alpha$-aminé de formule \ce{(CH3)2CH-CH(NH2)-COOH}. On dissout un comprimé contenant uniquement de la valine dans de l'eau distillée ; on obtient $V = \SI{100,0}{mL}$ de solution S. On dose $V_A = \SI{20,0}{mL}$ de S par une solution d'hydroxyde de sodium de concentration $c_B = \SI{5,00e-2}{mol.L^{-1}}$. L'équivalence est atteinte pour $V_{BE} = \SI{16,0}{mL}$.
\begin{enumerate}
\item Calculer la masse molaire de la valine.
\item Écrire l'équation-bilan de la réaction de dosage.
\item Déterminer la quantité de valine dans la prise d'essai, puis la masse de valine dans le comprimé.
\end{enumerate}
\tcblower
\textbf{1. Masse molaire de la valine.} Formule brute : \ce{C5H11NO2}.
\[ M = 5 \times 12 + 11 \times 1 + 2 \times 16 + 14 = 60 + 11 + 32 + 14 = \SI{117}{g.mol^{-1}} \]
(Vérification avec la formule générale $M = 14n + 47$ pour $n=5$ : $70 + 47 = 117$.)

\textbf{2. Équation-bilan du dosage.} En solution, la valine est sous forme d'amphion \ce{^{+}H3N-CH(iPr)-COO-}. Les ions \ce{HO-} arrachent le proton du groupe \ce{-NH3+} :
\[ \ce{^{+}H3N-CHR-COO- + HO- -> H2N-CHR-COO- + H2O} \quad (R = \ce{(CH3)2CH}-) \]
La réaction est totale, rapide et unique : c'est bien une réaction de dosage, \textbf{mole à mole}.

\textbf{3. Quantité de matière et masse.}

À l'équivalence : $n_A = n(\ce{HO-})_{\text{versé}} = c_B \times V_{BE}$.
\[ n_A = 5{,}00 \times 10^{-2} \times 16{,}0 \times 10^{-3} = 8{,}00 \times 10^{-4}\ \text{mol} \]
C'est la quantité de valine dans la prise d'essai de \SI{20,0}{mL}.

La solution S complète (\SI{100,0}{mL}) contient 5 fois plus :
\[ n_{\text{comprimé}} = n_A \times \frac{V}{V_A} = 8{,}00 \times 10^{-4} \times \frac{100{,}0}{20{,}0} = 4{,}00 \times 10^{-3}\ \text{mol} \]

Masse de valine dans le comprimé :
\[ m = n_{\text{comprimé}} \times M = 4{,}00 \times 10^{-3} \times 117 = \SI{0,468}{g} \]

\textbf{Conclusion :} le comprimé contient $m = \SI{0,468}{g}$ de valine, soit environ \SI{468}{mg} (3 chiffres significatifs, cohérents avec les données).
\end{exoresolu}

\begin{attention}[Les 5 erreurs qui coûtent des points au Bac]
\begin{enumerate}
\item \textbf{Oublier les charges de l'amphion.} L'amphion s'écrit \ce{^{+}H3N-CHR-COO-} : \ce{-NH3+} (et non \ce{-NH2}) et \ce{-COO-} (et non \ce{-COOH}). Une équation de dosage écrite avec la forme moléculaire \ce{H2N-CHR-COOH} est fausse en solution aqueuse.
\item \textbf{Confondre Gly-Ala et Ala-Gly.} L'ordre des résidus compte : on lit toujours du \textbf{N-terminal vers le C-terminal}. Échanger les deux acides aminés donne un dipeptide \emph{différent}.
\item \textbf{Dire que la glycine est chirale.} Son carbone $\alpha$ porte deux hydrogènes : elle est le \textbf{seul} acide $\alpha$-aminé naturel sans carbone asymétrique. Vérifier toujours que les quatre groupes sont différents avant de conclure à l'asymétrie.
\item \textbf{Inverser les séries D et L en Fischer.} \ce{COOH} en haut, R en bas : \ce{-NH2} à \textbf{gauche} = \textbf{L} (naturel), à droite = D. Et ne jamais tourner une projection de Fischer de 90°.
\item \textbf{Oublier les dipeptides « croisés » et « identiques ».} Sans protection, deux acides aminés A et B donnent \textbf{4} dipeptides : A-A, B-B, A-B \emph{et} B-A. Beaucoup de candidats n'en citent que deux. De même, dans la synthèse sélective, préciser \emph{quelle} fonction on bloque sur \emph{quel} acide aminé.
\end{enumerate}
\end{attention}

\newpage

\coursec{Exercices}

\courssub{Je vérifie mes connaissances}

\begin{exercice}[QCM : les acides $\alpha$-aminés]
Pour chaque question, choisir la (ou les) bonne(s) réponse(s) en justifiant brièvement.
\begin{enumerate}
\item La formule générale d'un acide $\alpha$-aminé est :
\begin{enumerate}
\item \ce{R-CH(NH2)-COOH} \quad b) \ce{R-CH(OH)-COOH} \quad c) \ce{R-CH(NH2)-CH2OH} \quad d) \ce{R-CO-NH2}
\end{enumerate}
\item Le groupe caractéristique de la liaison peptidique est :
\begin{enumerate}
\item le groupe ester \item le groupe amide \item le groupe carbonyle \item le groupe hydroxyle
\end{enumerate}
\item En solution aqueuse, un acide $\alpha$-aminé se trouve principalement sous la forme :
\begin{enumerate}
\item \ce{R-CH(NH2)-COOH} \item \ce{R-CH(NH3+)-COO-} \item \ce{R-CH(NH3+)-COOH} \item \ce{R-CH(NH2)-COO-}
\end{enumerate}
\item Les acides $\alpha$-aminés naturels ont, sauf exception, la configuration :
\begin{enumerate}
\item D \item L \item R \item Z
\end{enumerate}
\end{enumerate}
\end{exercice}

\begin{exercice}[Vrai ou faux]
Répondre par vrai ou faux en justifiant chaque réponse.
\begin{enumerate}
\item Tous les acides $\alpha$-aminés possèdent un carbone asymétrique.
\item L'amphion est une espèce globalement chargée qui porte deux charges de signes opposés.
\item La liaison peptidique est une liaison covalente entre le groupe carboxyle d'un acide aminé et le groupe amine d'un autre, avec élimination d'une molécule d'eau.
\item La condensation d'un mélange de deux acides $\alpha$-aminés différents, sans précaution particulière, ne fournit que deux dipeptides.
\end{enumerate}
\end{exercice}

\begin{exercice}[Texte à trous]
Compléter le texte suivant avec les mots ou groupes de mots : \emph{amphion}, \emph{asymétrique}, \emph{amide}, \emph{plane}, \emph{L}, \emph{zwitterion}, \emph{peptidique}, \emph{$\alpha$}.

Un acide $\alpha$-aminé possède un groupe amine et un groupe carboxyle portés par le même atome de carbone, dit carbone \ldots\ldots\ldots Ce carbone est, sauf pour la glycine, un carbone \ldots\ldots\ldots, ce qui confère à la molécule deux configurations notées D et \ldots\ldots\ldots ; les acides aminés naturels appartiennent à la série \ldots\ldots\ldots En solution aqueuse, il existe un transfert interne de proton donnant un ion dipolaire appelé \ldots\ldots\ldots ou \ldots\ldots\ldots La condensation de deux acides aminés crée une liaison \ldots\ldots\ldots, qui est le groupe \ldots\ldots\ldots ; cette liaison présente une géométrie \ldots\ldots\ldots
\end{exercice}

\begin{exercice}[Définitions et formules]
\begin{enumerate}
\item Définir : acide $\alpha$-aminé ; carbone asymétrique ; amphion ; liaison peptidique.
\item Écrire les formules semi-développées de la glycine (Gly), de l'alanine (Ala) et de la sérine (Ser).
\item Donner la signification des abréviations Val, Leu, Phe.
\end{enumerate}
\end{exercice}

\courssub{Je m'entraîne}

\begin{exercice}[Nomenclature et carbone asymétrique]
On donne la formule semi-développée de la valine : \chemfig{CH_3-CH(-[6]CH_3)-CH(-[2]NH_2)-C(=[1]O)-[7]OH}.
\begin{enumerate}
\item Écrire sa formule brute et vérifier qu'elle correspond à la formule générale d'un acide $\alpha$-aminé.
\item Nommer cette molécule en nomenclature systématique.
\item Repérer par un astérisque le (ou les) carbone(s) asymétrique(s).
\item Donner sa représentation de Fischer en précisant sa configuration (D ou L), sachant qu'il s'agit d'un acide aminé naturel.
\end{enumerate}
\end{exercice}

\begin{exercice}[Amphion et diagramme de prédominance]
On considère l'alanine \ce{CH3-CH(NH2)-COOH}, dont les deux couples acide/base ont pour $\mathrm{p}K_a$ : $\mathrm{p}K_{a1} = 2,3$ (couple \ce{-COOH}/\ce{-COO-}) et $\mathrm{p}K_{a2} = 9,7$ (couple \ce{-NH3+}/\ce{-NH2}).
\begin{enumerate}
\item Écrire l'équation-bilan de l'équilibre entre la forme moléculaire et la forme ionique (amphion) de l'alanine.
\item Écrire les deux équations-bilans mettant en évidence le caractère acide de l'amphion (avec \ce{HO-}) et son caractère basique (avec \ce{H3O+}).
\item Tracer le diagramme de prédominance des trois espèces (cation, amphion, anion) en fonction du pH.
\item Déterminer l'espèce majoritaire à $\mathrm{pH} = 2$, à $\mathrm{pH} = 6$ et à $\mathrm{pH} = 11$.
\end{enumerate}
\end{exercice}

\begin{exercice}[Formation d'un dipeptide]
On fait réagir l'alanine \ce{CH3-CH(NH2)-COOH} avec la sérine \ce{HO-CH2-CH(NH2)-COOH}.
\begin{enumerate}
\item Écrire l'équation-bilan de formation du dipeptide Ala-Ser et entourer la liaison peptidique.
\item Nommer tous les dipeptides obtenus par condensation d'un mélange équimolaire d'alanine et de sérine, sans protection des fonctions. Combien y en a-t-il ?
\item Écrire la formule semi-développée du dipeptide Ser-Ala et montrer qu'il est différent de Ala-Ser.
\end{enumerate}
\end{exercice}

\begin{exercice}[Hydrolyse d'un tripeptide]
L'hydrolyse totale d'un tripeptide fournit trois acides $\alpha$-aminés : la glycine, l'alanine et la valine.
\begin{enumerate}
\item Proposer deux enchaînements possibles de ce tripeptide et les nommer à l'aide des abréviations.
\item Expliquer pourquoi l'hydrolyse totale seule ne permet pas de déterminer la séquence du tripeptide.
\item Quelle information supplémentaire faudrait-il pour lever l'ambiguïté ?
\end{enumerate}
\end{exercice}

\courssub{Je me prépare au Bac}

\begin{exercice}[Type Bac : autour de la leucine]
\emph{Données :} $M(\ce{C}) = \SI{12}{g.mol^{-1}}$ ; $M(\ce{H}) = \SI{1}{g.mol^{-1}}$ ; $M(\ce{O}) = \SI{16}{g.mol^{-1}}$ ; $M(\ce{N}) = \SI{14}{g.mol^{-1}}$. Pour la leucine : $\mathrm{p}K_{a1} = 2,4$ et $\mathrm{p}K_{a2} = 9,6$.

Les protéines du niébé, légumineuse très consommée au Cameroun, renferment de la leucine, acide $\alpha$-aminé de formule semi-développée \chemfig{CH_3-CH(-[6]CH_3)-CH_2-CH(-[2]NH_2)-C(=[1]O)-[7]OH}.
\begin{enumerate}
\item Déterminer la formule brute de la leucine et calculer sa masse molaire.
\item Nommer la leucine en nomenclature systématique.
\item Identifier le carbone asymétrique et donner la représentation de Fischer de la leucine naturelle en précisant sa configuration.
\item Écrire l'équation-bilan de formation de l'amphion de la leucine en solution aqueuse.
\item Tracer le diagramme de prédominance des espèces de la leucine et indiquer l'espèce majoritaire dans le sang ($\mathrm{pH} = 7,4$).
\item On dissout $m = \SI{1,31}{g}$ de leucine dans de l'eau distillée pour obtenir $V = \SI{100}{mL}$ de solution. Calculer la concentration molaire de la solution.
\item Écrire l'équation-bilan de la réaction de la leucine avec la glycine \ce{NH2-CH2-COOH} conduisant au dipeptide Leu-Gly.
\end{enumerate}
\end{exercice}

\begin{exercice}[Type Bac : synthèse sélective d'un dipeptide]
\emph{Données :} masses molaires : $M(\text{Leu}) = \SI{131}{g.mol^{-1}}$ ; $M(\text{Gly}) = \SI{75}{g.mol^{-1}}$ ; $M(\text{Leu-Gly}) = \SI{188}{g.mol^{-1}}$.

Un laboratoire pharmaceutique de Douala souhaite préparer sélectivement le dipeptide Leu-Gly à partir de leucine \ce{(CH3)2CH-CH2-CH(NH2)-COOH} et de glycine \ce{NH2-CH2-COOH}.
\begin{enumerate}
\item Montrer que la condensation directe d'un mélange de leucine et de glycine conduit à quatre dipeptides différents ; les nommer.
\item Pour obtenir sélectivement Leu-Gly, décrire les quatre étapes de la synthèse (blocage, activation, condensation, déblocage) en précisant pour chacune la fonction concernée et le rôle de l'opération. On s'aidera de schémas réactionnels.
\item Écrire l'équation-bilan de l'étape de condensation.
\item On réalise la synthèse à partir de $n_1 = \SI{0,050}{mol}$ de leucine et d'un excès de glycine. Le rendement global de la synthèse est de $60\,\%$. Calculer la masse de dipeptide Leu-Gly effectivement obtenue.
\item Pourquoi dit-on que la liaison peptidique a une géométrie plane ? Quelle conséquence cela a-t-il sur la structure des protéines ?
\end{enumerate}
\end{exercice}

\begin{exercice}[Type Bac : dosage et propriétés de la glycine]
\emph{Données :} pour la glycine : $\mathrm{p}K_{a1} = 2,3$ ; $\mathrm{p}K_{a2} = 9,8$ ; $M(\text{Gly}) = \SI{75}{g.mol^{-1}}$ ; $K_e = 10^{-14}$ à \SI{25}{\celsius}.

La glycine, acide $\alpha$-aminé le plus simple, intervient dans la composition de nombreux compléments alimentaires vendus dans les pharmacies de Yaoundé.
\begin{enumerate}
\item Écrire la formule semi-développée de la glycine. La glycine possède-t-elle un carbone asymétrique ? Justifier.
\item Écrire les deux équations-bilans montrant le caractère amphotère de l'amphion de la glycine.
\item Tracer le diagramme de prédominance des espèces de la glycine. À quel pH l'amphion est-il l'espèce quasi exclusive ? (On admettra que ce pH, appelé point isoélectrique, vaut $\mathrm{pH}_i = \dfrac{\mathrm{p}K_{a1} + \mathrm{p}K_{a2}}{2}$ ; calculer sa valeur.)
\item On prépare une solution S de glycine de concentration $C = \SI{0,10}{mol.L^{-1}}$. On dose $V_a = \SI{20,0}{mL}$ de S, préalablement acidifiée (forme cationique \ce{NH3+-CH2-COOH} majoritaire), par une solution d'hydroxyde de sodium de concentration $C_b = \SI{0,10}{mol.L^{-1}}$.
\begin{enumerate}
\item Écrire les deux équations-bilans successives du dosage.
\item Déterminer les volumes équivalents $V_{e1}$ et $V_{e2}$ correspondant aux deux sauts de pH.
\item Indiquer l'espèce majoritaire en solution pour $V_b = V_{e1}$.
\end{enumerate}
\item La glycine peut réagir avec l'alanine. Écrire l'équation-bilan de formation du dipeptide Gly-Ala et calculer la masse de ce dipeptide obtenue à partir de $\SI{7,5}{g}$ de glycine en supposant la réaction totale et sélective.
\end{enumerate}
\end{exercice}

\newpage

\coursec{Corrigés des exercices}

\coursec{Corrigés détaillés : Acides $\alpha$-aminés et synthèse des peptides}

\begin{corrige}[Exercice 1 — QCM : les acides $\alpha$-aminés]
\begin{enumerate}
\item \textbf{Réponse a)} : \ce{R-CH(NH2)-COOH}. Un acide $\alpha$-aminé porte le groupe amine \ce{-NH2} et le groupe carboxyle \ce{-COOH} sur le \emph{même} carbone, le carbone $\alpha$. La réponse b) correspond à un acide $\alpha$-hydroxylé, c) à un aminoalcool, d) à un amide.
\item \textbf{Réponse b)} : le groupe \cle{amide} \ce{-CO-NH-}. C'est le groupe caractéristique de la liaison peptidique, formé par condensation entre \ce{-COOH} et \ce{-NH2} avec élimination d'eau.
\item \textbf{Réponse b)} : \ce{R-CH(NH3+)-COO-}. En solution aqueuse, un transfert interne de proton a lieu entre le groupe carboxyle (acide) et le groupe amine (basique) : l'espèce majoritaire au voisinage de la neutralité est l'\cle{amphion} (ou zwitterion).
\item \textbf{Réponse b)} : configuration \cle{L}. Les acides $\alpha$-aminés naturels (constituants des protéines) appartiennent à la série L, la glycine exceptée car elle n'est pas chirale.
\end{enumerate}
\end{corrige}

\begin{corrige}[Exercice 2 — Vrai ou faux]
\begin{enumerate}
\item \textbf{Faux.} La glycine \ce{NH2-CH2-COOH} possède deux atomes d'hydrogène sur le carbone $\alpha$ : celui-ci ne porte pas quatre groupes différents et n'est donc pas asymétrique. Tous les autres acides $\alpha$-aminés naturels possèdent un carbone asymétrique.
\item \textbf{Vrai.} L'amphion \ce{$R$-CH(NH3+)-COO-} porte une charge positive (sur \ce{-NH3+}) et une charge négative (sur \ce{-COO-}) ; il est globalement chargé mais électriquement neutre (charge nette nulle) : c'est un ion dipolaire.
\item \textbf{Vrai.} C'est une condensation :
\[ \ce{$R$-CH(NH2)-C(=O)OH + H2N-CH($R'$)-C(=O)OH -> $R$-CH(NH2)-C(=O)-NH-CH($R'$)-C(=O)OH + H2O} \]
La liaison \ce{C-N} créée au sein du groupe \ce{-C(=O)-NH-} est covalente.
\item \textbf{Faux.} Sans protection, chaque acide aminé peut réagir avec lui-même et avec l'autre : on obtient \textbf{quatre} dipeptides. Pour deux acides aminés A et B : A-A, B-B, A-B et B-A (les dipeptides A-B et B-A sont distincts car la liaison peptidique est orientée).
\end{enumerate}
\end{corrige}

\begin{corrige}[Exercice 3 — Texte à trous]
Un acide $\alpha$-aminé possède un groupe amine et un groupe carboxyle portés par le même atome de carbone, dit carbone \textbf{$\alpha$}. Ce carbone est, sauf pour la glycine, un carbone \textbf{asymétrique}, ce qui confère à la molécule deux configurations notées D et \textbf{L} ; les acides aminés naturels appartiennent à la série \textbf{L}. En solution aqueuse, il existe un transfert interne de proton donnant un ion dipolaire appelé \textbf{amphion} ou \textbf{zwitterion}. La condensation de deux acides aminés crée une liaison \textbf{peptidique}, qui est le groupe \textbf{amide} ; cette liaison présente une géométrie \textbf{plane}.
\end{corrige}

\begin{corrige}[Exercice 4 — Définitions et formules]
\begin{enumerate}
\item \textbf{Définitions.}
\begin{itemize}
\item \cle{Acide $\alpha$-aminé} : composé organique possédant un groupe carboxyle \ce{-COOH} et un groupe amine \ce{-NH2} portés par le même atome de carbone (le carbone $\alpha$) ; formule générale \ce{R-CH(NH2)-COOH}.
\item \cle{Carbone asymétrique} : atome de carbone tétragonal (hybridation $sp^3$) lié à quatre atomes ou groupes d'atomes tous différents.
\item \cle{Amphion} (ou zwitterion) : ion dipolaire \ce{R-CH(NH3+)-COO-} résultant du transfert interne d'un proton du groupe carboxyle vers le groupe amine.
\item \cle{Liaison peptidique} : liaison covalente \ce{-CO-NH-} (groupe amide) formée par condensation entre le groupe carboxyle d'un acide aminé et le groupe amine d'un autre, avec élimination d'eau.
\end{itemize}
\item \textbf{Formules semi-développées.}
\begin{itemize}
\item Glycine (Gly) : \chemfig{NH_2-CH_2-C(=[1]O)-[7]OH}
\item Alanine (Ala) : \chemfig{CH_3-CH(-[2]NH_2)-C(=[1]O)-[7]OH}
\item Sérine (Ser) : \chemfig{HO-CH_2-CH(-[2]NH_2)-C(=[1]O)-[7]OH}
\end{itemize}
\item \textbf{Abréviations usuelles (trois lettres) :} Val = valine ; Leu = leucine ; Phe = phénylalanine.
\end{enumerate}
\end{corrige}

\begin{corrige}[Exercice 5 — Nomenclature et carbone asymétrique (valine)]
\begin{enumerate}
\item \textbf{Formule brute.} En comptant les atomes de \chemfig{CH_3-CH(-[6]CH_3)-CH(-[2]NH_2)-C(=[1]O)-[7]OH} : 5 atomes de carbone, 11 atomes d'hydrogène, 1 atome d'azote, 2 atomes d'oxygène, soit \ce{C5H11NO2}. Elle s'écrit \ce{(CH3)2CH-CH(NH2)-COOH}, de la forme \ce{R-CH(NH2)-COOH} avec \ce{R = -CH(CH3)2} : c'est bien un acide $\alpha$-aminé.
\item \textbf{Nomenclature systématique.} La chaîne principale contenant \ce{-COOH} comporte 4 atomes de carbone (acide butanoïque), avec un groupe amino en position 2 et un groupe méthyle en position 3 : \textbf{acide 2-amino-3-méthylbutanoïque}.
\item \textbf{Carbone asymétrique.} Le carbone $\alpha$ (C2) porte quatre groupes différents : \ce{-H}, \ce{-NH2}, \ce{-COOH} et \ce{-CH(CH3)2} :
\[ \chemfig{CH_3-CH(-[6]CH_3)-C^{*}(-[2]NH_2)(-[4]H)-C(=[1]O)-[7]OH} \]
Les autres carbones portent au moins deux groupes identiques (\ce{CH3} ou \ce{H}) : un seul carbone asymétrique.
\item \textbf{Représentation de Fischer.} La chaîne carbonée est verticale, \ce{-COOH} en haut ; pour un acide aminé naturel (série L), le groupe \ce{-NH2} est à \textbf{gauche} :
\begin{center}
\chemfig{COOH-[:-90](-[180]NH_2)(-[0]H)-[:-90]CH(-[180]CH_3)(-[0]CH_3)}
\end{center}
Configuration \textbf{L} : c'est la L-valine.
\end{enumerate}
\end{corrige}

\begin{corrige}[Exercice 6 — Amphion et diagramme de prédominance (alanine)]
\begin{enumerate}
\item \textbf{Équilibre forme moléculaire / forme ionique} (transfert interne de proton) :
\[ \ce{CH3-CH(NH2)-COOH <=> CH3-CH(NH3+)-COO-} \]
\item \textbf{Caractère acide de l'amphion} (il cède le proton de \ce{-NH3+} à \ce{HO-}) :
\[ \ce{CH3-CH(NH3+)-COO- + HO- -> CH3-CH(NH2)-COO- + H2O} \]
\textbf{Caractère basique de l'amphion} (il capte un proton sur \ce{-COO-}) :
\[ \ce{CH3-CH(NH3+)-COO- + H3O+ -> CH3-CH(NH3+)-COOH + H2O} \]
L'amphion est donc \cle{amphotère}.
\item \textbf{Diagramme de prédominance.} Le couple $\mathrm{p}K_{a1} = 2,3$ concerne \ce{-COOH}/\ce{-COO-} (transition cation $\to$ amphion) et $\mathrm{p}K_{a2} = 9,7$ concerne \ce{-NH3+}/\ce{-NH2} (transition amphion $\to$ anion) :
\begin{popfigure}
\begin{tikzpicture}[x=0.72cm]
\draw[->, thick, popdark] (0,0) -- (14.5,0) node[below left=2pt and 0pt, font=\small]{pH};
\foreach \x/\lab in {2.3/{2,3}, 9.7/{9,7}}{
  \draw[thick, popdark] (\x,0.08) -- (\x,-0.08) node[below, font=\small]{\lab};
}
\draw[dashed, popmuted] (2.3,0) -- (2.3,0.9);
\draw[dashed, popmuted] (9.7,0) -- (9.7,0.9);
\node[font=\small, text=popblue] at (1.15,0.55) {cation};
\node[font=\small, text=popblue] at (1.15,0.2) {\ce{CH3-CH(NH3+)-COOH}};
\node[font=\small, text=poppurple] at (6.0,0.55) {amphion};
\node[font=\small, text=poppurple] at (6.0,0.2) {\ce{CH3-CH(NH3+)-COO-}};
\node[font=\small, text=popgreen] at (12.1,0.55) {anion};
\node[font=\small, text=popgreen] at (12.1,0.2) {\ce{CH3-CH(NH2)-COO-}};
\end{tikzpicture}
\legende{Diagramme de prédominance des espèces de l'alanine.}
\end{popfigure}
\item \textbf{Espèce majoritaire.}
\begin{itemize}
\item $\mathrm{pH} = 2 < 2,3$ : le \textbf{cation} \ce{CH3-CH(NH3+)-COOH} ;
\item $\mathrm{pH} = 6$ (entre 2,3 et 9,7) : l'\textbf{amphion} \ce{CH3-CH(NH3+)-COO-} ;
\item $\mathrm{pH} = 11 > 9,7$ : l'\textbf{anion} \ce{CH3-CH(NH2)-COO-}.
\end{itemize}
\end{enumerate}
\end{corrige}

\begin{corrige}[Exercice 7 — Formation d'un dipeptide (Ala-Ser)]
\begin{enumerate}
\item \textbf{Équation-bilan.} Le groupe carboxyle de l'alanine réagit avec le groupe amine de la sérine :
\[ \ce{CH3-CH(NH2)-COOH + H2N-CH(CH2OH)-COOH -> CH3-CH(NH2)-\boxed{CO-NH}-CH(CH2OH)-COOH + H2O} \]
Formellement, la liaison peptidique est le groupe \ce{-C(=O)-NH-} :
\begin{center}
\chemfig{CH_3-CH(-[2]NH_2)-C(=[1]O)-[7]NH-CH(-[2]CH_2OH)-C(=[1]O)-[7]OH}
\end{center}
Le groupe \ce{-CO-NH-} encadré est la liaison peptidique (groupe amide).
\item \textbf{Dipeptides obtenus sans protection.} Chaque acide aminé peut fournir soit sa fonction \ce{-COOH}, soit sa fonction \ce{-NH2} : on obtient \textbf{quatre} dipeptides :
\begin{itemize}
\item \textbf{Ala-Ala} (deux alanines) ;
\item \textbf{Ser-Ser} (deux sérines) ;
\item \textbf{Ala-Ser} (COOH de l'alanine + \ce{NH2} de la sérine) ;
\item \textbf{Ser-Ala} (COOH de la sérine + \ce{NH2} de l'alanine).
\end{itemize}
\item \textbf{Ser-Ala} : \chemfig{HO-CH_2-CH(-[2]NH_2)-C(=[1]O)-[7]NH-CH(-[2]CH_3)-C(=[1]O)-[7]OH}, soit \ce{HO-CH2-CH(NH2)-CO-NH-CH(CH3)-COOH}.

Dans Ser-Ala, l'extrémité N-terminale porte le résidu sérine (groupe \ce{-CH2OH} du côté de l'amine libre) ; dans Ala-Ser, c'est l'alanine qui est N-terminale. Les groupes latéraux \ce{-CH3} et \ce{-CH2OH} ne sont pas placés aux mêmes positions : ce sont deux molécules \textbf{différentes} (la liaison peptidique est orientée).
\end{enumerate}
\end{corrige}

\begin{corrige}[Exercice 8 — Hydrolyse d'un tripeptide]
\begin{enumerate}
\item \textbf{Enchaînements possibles} (deux exemples parmi les six permutations) :
\begin{itemize}
\item Gly-Ala-Val : \ce{H2N-CH2-CO-NH-CH(CH3)-CO-NH-CH(CH(CH3)2)-COOH} ;
\item Val-Gly-Ala : \ce{H2N-CH(CH(CH3)2)-CO-NH-CH2-CO-NH-CH(CH3)-COOH}.
\end{itemize}
\item \textbf{Limite de l'hydrolyse totale.} L'hydrolyse totale rompt toutes les liaisons peptidiques : elle fournit la \emph{composition} du peptide (nature et proportions des acides aminés) mais détruit l'information sur leur \emph{ordre} d'enchaînement (la séquence). Avec trois acides aminés différents, il existe $3! = 6$ séquences possibles.
\item \textbf{Information supplémentaire.} Il faut déterminer la séquence, par exemple en identifiant l'acide aminé \cle{N-terminal} (extrémité portant la fonction amine libre) puis l'acide aminé \cle{C-terminal}, ou en réalisant des hydrolyses partielles et en analysant les dipeptides formés.
\end{enumerate}
\end{corrige}

\begin{corrige}[Exercice 9 — Type Bac : autour de la leucine]
\emph{Barème indicatif (sur 10 pts) : 1) 1,5 pt ; 2) 1 pt ; 3) 2 pts ; 4) 1 pt ; 5) 1,5 pt ; 6) 1,5 pt ; 7) 1,5 pt.}
\begin{enumerate}
\item \textbf{Formule brute et masse molaire.} La formule \chemfig{CH_3-CH(-[6]CH_3)-CH_2-CH(-[2]NH_2)-C(=[1]O)-[7]OH} donne \ce{C6H13NO2}.
\[ M = 6\times 12 + 13\times 1 + 14 + 2\times 16 = 72 + 13 + 14 + 32 = \SI{131}{g.mol^{-1}}. \]
\item \textbf{Nomenclature systématique} : acide \textbf{2-amino-4-méthylpentanoïque} (chaîne de 5 C incluant le carboxyle, amino en C2, méthyle en C4).
\item \textbf{Carbone asymétrique et Fischer.} Le carbone $\alpha$ (C2) porte \ce{-H}, \ce{-NH2}, \ce{-COOH} et \ce{-CH2-CH(CH3)2} : quatre groupes différents, il est asymétrique. La leucine naturelle est de configuration \textbf{L} : \ce{-NH2} à gauche dans la représentation de Fischer, \ce{-COOH} en haut :
\begin{center}
\chemfig{COOH-[:-90](-[180]NH_2)(-[0]H)-[:-90]CH_2-CH(-[180]CH_3)(-[0]CH_3)}
\end{center}
\item \textbf{Formation de l'amphion} :
\[ \ce{(CH3)2CH-CH2-CH(NH2)-COOH <=> (CH3)2CH-CH2-CH(NH3+)-COO-} \]
\item \textbf{Diagramme de prédominance} ($\mathrm{p}K_{a1} = 2,4$ ; $\mathrm{p}K_{a2} = 9,6$) :
\begin{popfigure}
\begin{tikzpicture}[x=0.72cm]
\draw[->, thick, popdark] (0,0) -- (14.5,0) node[below left=2pt and 0pt, font=\small]{pH};
\foreach \x/\lab in {2.4/{2,4}, 7.4/{7,4}, 9.6/{9,6}}{
  \draw[thick, popdark] (\x,0.08) -- (\x,-0.08) node[below, font=\small]{\lab};
}
\draw[dashed, popmuted] (2.4,0) -- (2.4,0.9);
\draw[dashed, popmuted] (9.6,0) -- (9.6,0.9);
\draw[dashed, poporange] (7.4,0) -- (7.4,0.9);
\node[font=\small, text=popblue] at (1.2,0.5) {cation};
\node[font=\small, text=poppurple] at (6.0,0.5) {amphion \ce{Leu^{\pm}}};
\node[font=\small, text=popgreen] at (12.0,0.5) {anion};
\end{tikzpicture}
\legende{Diagramme de prédominance de la leucine ; le trait orange repère le pH du sang.}
\end{popfigure}
Dans le sang, $\mathrm{pH} = 7,4$ est compris entre 2,4 et 9,6 : l'espèce majoritaire est l'\textbf{amphion} \ce{(CH3)2CH-CH2-CH(NH3+)-COO-}.
\item \textbf{Concentration molaire.}
\[ n = \frac{m}{M} = \frac{1,31}{131} = \num{1.00e-2}\,\text{mol} \;;\qquad C = \frac{n}{V} = \frac{\num{1.00e-2}}{0,100} = \SI{0.100}{mol.L^{-1}}. \]
\item \textbf{Formation de Leu-Gly} :
\[ \ce{(CH3)2CH-CH2-CH(NH2)-COOH + H2N-CH2-COOH -> (CH3)2CH-CH2-CH(NH2)-CO-NH-CH2-COOH + H2O} \]
\end{enumerate}
\textbf{Points de vigilance.} Ne pas oublier les deux \ce{CH3} du groupe isopropyle dans le comptage des atomes ; en Fischer, la chaîne carbonée est verticale et le \ce{-NH2} à gauche pour la série L ; vérifier que $\mathrm{pH} = 7,4$ se situe bien entre les deux $\mathrm{p}K_a$ avant de conclure.
\end{corrige}

\begin{corrige}[Exercice 10 — Type Bac : synthèse sélective de Leu-Gly]
\emph{Barème indicatif (sur 10 pts) : 1) 2 pts ; 2) 3 pts ; 3) 1,5 pt ; 4) 2 pts ; 5) 1,5 pt.}
\begin{enumerate}
\item \textbf{Quatre dipeptides possibles.} Sans protection, la fonction \ce{-COOH} de chaque acide aminé peut réagir avec la fonction \ce{-NH2} de l'autre ou de lui-même :
\begin{center}
\begin{tabularx}{\linewidth}{|l|X|}\hline
\rowcolor{popblueL} \textbf{Dipeptide} & \textbf{Fonctions ayant réagi} \\\hline
Leu-Leu & \ce{-COOH} de Leu + \ce{-NH2} de Leu \\\hline
Gly-Gly & \ce{-COOH} de Gly + \ce{-NH2} de Gly \\\hline
Leu-Gly & \ce{-COOH} de Leu + \ce{-NH2} de Gly \\\hline
Gly-Leu & \ce{-COOH} de Gly + \ce{-NH2} de Leu \\\hline
\end{tabularx}
\end{center}
On obtient donc un mélange de \textbf{quatre} dipeptides : la condensation directe n'est pas sélective.
\item \textbf{Les quatre étapes de la synthèse sélective de Leu-Gly.} On veut faire réagir le \ce{-COOH} de la leucine avec le \ce{-NH2} de la glycine, et uniquement ceux-là.
\begin{enumerate}
\item \textbf{Blocage (protection)} de la fonction \ce{-NH2} de la \emph{leucine} : on la transforme en une fonction non nucléophile (carbamate, par exemple) pour empêcher l'auto-condensation Leu-Leu :
\[ \ce{Leu-NH2 ->[protection] Leu-NH-P} \]
\item \textbf{Activation} de la fonction \ce{-COOH} de la \emph{leucine} : on la transforme en dérivé plus réactif (chlorure d'acyle, ester activé) pour faciliter l'attaque par l'amine :
\[ \ce{Leu-COOH ->[activation] Leu-COCl} \]
\item \textbf{Condensation} : on fait réagir la leucine (amine bloquée, acide activé) avec la glycine dont la fonction \ce{-NH2} est libre ; seule la liaison Leu-Gly peut se former.
\item \textbf{Déblocage (déprotection)} : on régénère la fonction \ce{-NH2} de la leucine sans rompre la liaison peptidique, ce qui livre le dipeptide Leu-Gly pur.
\end{enumerate}
\item \textbf{Équation-bilan de la condensation} (avec A = groupe activant, P = groupe protecteur) :
\[ \ce{P-HN-Leu-CO-A + H2N-CH2-COOH -> P-HN-Leu-CO-NH-CH2-COOH + HA} \]
puis, après déblocage : \ce{H2N-Leu-CO-NH-CH2-COOH}, soit globalement :
\[ \ce{(CH3)2CH-CH2-CH(NH2)-COOH + H2N-CH2-COOH -> (CH3)2CH-CH2-CH(NH2)-CO-NH-CH2-COOH + H2O} \]
\item \textbf{Masse de Leu-Gly obtenue.} Tableau d'avancement (glycine en excès, leucine limitante) :
\begin{center}
\begin{tabularx}{\linewidth}{|l|X|X|X|X|}\hline
\rowcolor{popblueL} État & Leu & Gly & Leu-Gly & \ce{H2O} \\\hline
Initial (mol) & 0,050 & excès & 0 & --- \\\hline
Final théorique (mol) & 0 & excès & 0,050 & 0,050 \\\hline
\end{tabularx}
\end{center}
Quantité théorique : $n_{\text{th}} = \SI{0.050}{mol}$. Avec un rendement $\eta = 60\,\%$ :
\[ n_{\text{réel}} = \eta \times n_{\text{th}} = 0,60 \times 0,050 = \SI{0.030}{mol} \]
\[ m = n_{\text{réel}} \times M(\text{Leu-Gly}) = 0,030 \times 188 = \SI{5.64}{g}. \]
\item \textbf{Géométrie plane de la liaison peptidique.} La liaison \ce{C-N} du groupe amide possède un \emph{caractère partiel de double liaison} (délocalisation du doublet de l'azote vers le carbonyle) : la rotation autour de cette liaison est bloquée et les six atomes \ce{C}, \ce{O}, \ce{N}, \ce{H} et les deux carbones $\alpha$ voisins sont \textbf{coplanaires}. Conséquence : la chaîne polypeptidique ne peut tourner librement qu'autour des liaisons des carbones $\alpha$, ce qui impose aux protéines des structures ordonnées (hélice $\alpha$, feuillet $\beta$) et donc leur forme spatiale, condition de leur activité biologique.
\end{enumerate}
\textbf{Points de vigilance.} Bien protéger la fonction \ce{-NH2} de l'acide aminé dont on veut engager le \ce{-COOH} (ici la leucine) ; ne pas confondre blocage et activation ; dans le calcul de rendement, appliquer le pourcentage à la quantité \emph{théorique} et non à la masse des réactifs.
\end{corrige}

\begin{corrige}[Exercice 11 — Type Bac : dosage et propriétés de la glycine]
\emph{Barème indicatif (sur 10 pts) : 1) 1,5 pt ; 2) 1,5 pt ; 3) 2 pts ; 4) 3 pts ; 5) 2 pts.}
\begin{enumerate}
\item \textbf{Formule de la glycine} : \chemfig{NH_2-CH_2-C(=[1]O)-[7]OH}, soit \ce{NH2-CH2-COOH}. Le carbone $\alpha$ porte \textbf{deux} atomes d'hydrogène identiques : il n'est lié qu'à trois groupes différents, donc il n'est \textbf{pas asymétrique}. La glycine est le seul acide $\alpha$-aminé naturel \emph{achiral} (pas de configurations D/L).
\item \textbf{Caractère amphotère de l'amphion} \ce{NH3+-CH2-COO-} :
\begin{itemize}
\item caractère acide : $\ce{NH3+-CH2-COO- + HO- -> NH2-CH2-COO- + H2O}$ ;
\item caractère basique : $\ce{NH3+-CH2-COO- + H3O+ -> NH3+-CH2-COOH + H2O}$.
\end{itemize}
\item \textbf{Diagramme de prédominance} ($\mathrm{p}K_{a1} = 2,3$ ; $\mathrm{p}K_{a2} = 9,8$) :
\begin{popfigure}
\begin{tikzpicture}[x=0.72cm]
\draw[->, thick, popdark] (0,0) -- (14.5,0) node[below left=2pt and 0pt, font=\small]{pH};
\foreach \x/\lab in {2.3/{2,3}, 6.05/{6,05}, 9.8/{9,8}}{
  \draw[thick, popdark] (\x,0.08) -- (\x,-0.08) node[below, font=\small]{\lab};
}
\draw[dashed, popmuted] (2.3,0) -- (2.3,0.9);
\draw[dashed, popmuted] (9.8,0) -- (9.8,0.9);
\draw[dashed, poporange] (6.05,0) -- (6.05,0.9);
\node[font=\small, text=popblue] at (1.15,0.5) {cation \ce{NH3+-CH2-COOH}};
\node[font=\small, text=poppurple] at (6.05,0.5) {amphion \ce{NH3+-CH2-COO-}};
\node[font=\small, text=popgreen] at (12.15,0.5) {anion \ce{NH2-CH2-COO-}};
\end{tikzpicture}
\legende{Diagramme de prédominance de la glycine ; le trait orange repère le point isoélectrique.}
\end{popfigure}
Point isoélectrique :
\[ \mathrm{pH}_i = \frac{\mathrm{p}K_{a1} + \mathrm{p}K_{a2}}{2} = \frac{2,3 + 9,8}{2} = \frac{12,1}{2} = 6,05. \]
À $\mathrm{pH} = \mathrm{pH}_i = 6,05$, l'amphion est l'espèce quasi exclusive (sa charge nette moyenne est nulle).
\item \textbf{Dosage de la forme cationique par \ce{NaOH}.}
\begin{enumerate}
\item \textbf{Équations successives.} Le cation est un diacide : on neutralise d'abord le groupe \ce{-COOH} (première acidité), puis le groupe \ce{-NH3+} (deuxième acidité) :
\[ \text{(1)}\quad \ce{NH3+-CH2-COOH + HO- -> NH3+-CH2-COO- + H2O} \]
\[ \text{(2)}\quad \ce{NH3+-CH2-COO- + HO- -> NH2-CH2-COO- + H2O} \]
\item \textbf{Volumes équivalents.} Quantité de glycine dosée :
\[ n_a = C \times V_a = 0,10 \times 20,0\times 10^{-3} = \num{2.0e-3}\,\text{mol}. \]
Première équivalence : $C_b\,V_{e1} = n_a$, d'où
\[ V_{e1} = \frac{n_a}{C_b} = \frac{\num{2.0e-3}}{0,10} = \SI{2.0e-2}{L} = \SI{20.0}{mL}. \]
La deuxième équivalence nécessite le même volume supplémentaire :
\[ V_{e2} = 2\,V_{e1} = \SI{40.0}{mL}. \]
\item \textbf{Espèce majoritaire à $V_b = V_{e1}$.} La première acidité vient d'être totalement neutralisée et la deuxième n'a pas encore réagi : l'espèce majoritaire est l'\textbf{amphion} \ce{NH3+-CH2-COO-} (le pH vaut alors environ $\mathrm{pH}_i = 6,05$).
\end{enumerate}
\item \textbf{Formation de Gly-Ala et masse obtenue.}
\[ \ce{NH2-CH2-COOH + CH3-CH(NH2)-COOH -> NH2-CH2-CO-NH-CH(CH3)-COOH + H2O} \]
Quantité de glycine : $n = \dfrac{m}{M} = \dfrac{7,5}{75} = \SI{0.10}{mol}$. Réaction supposée totale et sélective : $n(\text{Gly-Ala}) = \SI{0.10}{mol}$.

Masse molaire du dipeptide : $M(\text{Gly-Ala}) = M(\text{Gly}) + M(\text{Ala}) - M(\ce{H2O}) = 75 + 89 - 18 = \SI{146}{g.mol^{-1}}$.
\[ m(\text{Gly-Ala}) = 0,10 \times 146 = \SI{14.6}{g}. \]
\end{enumerate}
\textbf{Points de vigilance.} Ne pas oublier que la forme cationique est un \emph{diacide} : il y a deux équivalences et $V_{e2} = 2V_{e1}$ ; pour la masse molaire d'un dipeptide, penser à \emph{retrancher} la molécule d'eau éliminée lors de la condensation ; conserver les virgules décimales et les unités à chaque étape.
\end{corrige}

\newpage

\coursec{Fiche bilan}

\begin{popfigure}
\begin{tikzpicture}[
  central/.style={circle, draw=popdark, fill=popblue!25, text width=3.2cm, align=center, font=\bfseries\small, inner sep=3pt},
  branche/.style={rectangle, rounded corners=4pt, draw=popdark, text width=3.1cm, align=center, font=\scriptsize, inner sep=3pt},
  lien/.style={thick, popline}
]
\node[central] (C) {Acides $\alpha$-aminés et peptides};
\node[branche, fill=poporangeL] (N1) at (-5.4,2.6) {\textbf{Structure}\\ \ce{H2N-CHR-COOH}\\ 20 acides naturels (Gly, Ala, Val\ldots)};
\node[branche, fill=poppinkL] (N2) at (-5.4,-2.6) {\textbf{Stéréochimie}\\ C$^*$ asymétrique\\ Fischer : séries \textsc{d} et \textsc{l} (naturels = \textsc{l})};
\node[branche, fill=popgreenL] (N3) at (0,3.4) {\textbf{Amphion}\\ \ce{^{+}H3N-CHR-COO^{-}}\\ 2 couples, 2 p$K_a$\\ diagramme de prédominance};
\node[branche, fill=popgoldL] (N4) at (5.4,2.6) {\textbf{Liaison peptidique}\\ \ce{-CO-NH-} plane\\ condensation $-$ \ce{H2O}};
\node[branche, fill=poppurpleL] (N5) at (5.4,-2.6) {\textbf{Polypeptides}\\ Gly-Ala\ldots\\ extrémités N et C\\ protéines ($n > 50$)};
\node[branche, fill=popblueL] (N6) at (0,-3.4) {\textbf{Synthèse sélective}\\ bloquer -- activer -- condenser -- débloquer};
\draw[lien] (C) -- (N1);
\draw[lien] (C) -- (N2);
\draw[lien] (C) -- (N3);
\draw[lien] (C) -- (N4);
\draw[lien] (C) -- (N5);
\draw[lien] (C) -- (N6);
\end{tikzpicture}
\legende{Carte mentale de la leçon : les six idées-forces à maîtriser.}
\end{popfigure}

\begin{bilan}[L'essentiel de la leçon]
\textbf{1. Définitions clés}
\begin{itemize}
  \item \cle{Acide $\alpha$-aminé} : molécule portant sur le \emph{même} carbone ($\alpha$) une fonction amine \ce{-NH2} et une fonction acide carboxylique \ce{-COOH}. Formule générale : \ce{H2N-CHR-COOH}.
  \item \cle{Carbone asymétrique} (C$^*$) : carbone tétraédrique lié à 4 groupes différents $\Rightarrow$ molécule chirale, 2 énantiomères. Tous les acides $\alpha$-aminés naturels sont chiraux \textbf{sauf la glycine} (R = H).
  \item \cle{Amphion} (zwitterion) : forme ionique \ce{^{+}H3N-CHR-COO^{-}}, résultant d'un transfert interne de \ce{H+}. C'est un \cle{ampholyte}.
  \item \cle{Liaison peptidique} : groupe amide \ce{-CO-NH-} formé par condensation entre \ce{-COOH} et \ce{-NH2} ; elle est \textbf{plane} (caractère partiel de double liaison).
  \item \cle{Peptide} : nommé en citant les résidus depuis l'extrémité N-terminale : Gly-Ala $\neq$ Ala-Gly.
\end{itemize}

\textbf{2. Équations-types à connaître par c\oe ur}
\begin{itemize}
  \item Équilibre forme moléculaire / amphion :
  \[ \mathrm{H_2N-CH(R)-COOH} \ce{<=>} \mathrm{H_3N^+-CH(R)-COO^-} \]
  \item Caractère basique de l'amphion (milieu acide) :
  \[ \mathrm{H_3N^+-CH(R)-COO^-} + \ce{H3O+} \ce{->} \mathrm{H_3N^+-CH(R)-COOH} + \ce{H2O} \]
  \item Caractère acide de l'amphion (milieu basique) :
  \[ \mathrm{H_3N^+-CH(R)-COO^-} + \ce{OH-} \ce{->} \mathrm{H_2N-CH(R)-COO^-} + \ce{H2O} \]
  \item Formation d'un dipeptide (condensation) :
  \[ \mathrm{H_2N-CH(R)-COOH} + \mathrm{H_2N-CH(R')-COOH} \ce{->} \mathrm{H_2N-CH(R)-CO-NH-CH(R')-COOH} + \ce{H2O} \]
\end{itemize}

\textbf{3. Les deux couples et le diagramme de prédominance}
\begin{itemize}
  \item Couples : \ce{^{+}H3N-CHR-COOH}/\ce{^{+}H3N-CHR-COO^{-}} (p$K_{a1} \approx 2$) et \ce{^{+}H3N-CHR-COO^{-}}/\ce{H2N-CHR-COO^{-}} (p$K_{a2} \approx 9$).
  \item pH $<$ p$K_{a1}$ : \textbf{cation} ; p$K_{a1}$ $<$ pH $<$ p$K_{a2}$ : \textbf{amphion} ; pH $>$ p$K_{a2}$ : \textbf{anion}.
\end{itemize}

\textbf{4. Quelques acides $\alpha$-aminés usuels}

\begin{center}
\begin{tabularx}{\linewidth}{|l|X|l|}
\hline
\rowcolor{popblueL} \textbf{Nom} & \textbf{Chaîne latérale R} & \textbf{Abréviation} \\
\hline
Glycine & \ce{-H} & Gly \\
\hline
Alanine & \ce{-CH3} & Ala \\
\hline
Valine & \ce{-CH(CH3)2} & Val \\
\hline
Leucine & \ce{-CH2-CH(CH3)2} & Leu \\
\hline
Sérine & \ce{-CH2-OH} & Ser \\
\hline
Phénylalanine & \ce{-CH2-C6H5} & Phe \\
\hline
\end{tabularx}
\end{center}

\textbf{5. Méthodes express}
\begin{itemize}
  \item \cle{Fischer} : chaîne carbonée verticale, \ce{COOH} en haut ; \ce{NH2} à droite = série \textsc{d}, à gauche = série \textsc{l} (naturels).
  \item \cle{Nombre de dipeptides} sans protection : 2 acides aminés A et B donnent \textbf{4} dipeptides (A-A, B-B, A-B, B-A) ; avec $n$ acides différents : $n^2$ dipeptides.
  \item \cle{Synthèse sélective} d'un dipeptide A-B : (1) \textbf{bloquer} \ce{-NH2} de A et \ce{-COOH} de B ; (2) \textbf{activer} \ce{-COOH} de A ; (3) \textbf{condenser} A + B ; (4) \textbf{débloquer} les fonctions protégées.
\end{itemize}
\end{bilan}

\begin{aretenir}[Checklist avant le Bac]
\begin{itemize}
  \item $\square$ Je sais écrire la formule générale \ce{H2N-CHR-COOH} et repérer le carbone $\alpha$.
  \item $\square$ Je sais identifier un carbone asymétrique (4 substituants différents) et citer l'exception : la glycine.
  \item $\square$ Je sais construire une représentation de Fischer et distinguer les configurations \textsc{d} et \textsc{l}.
  \item $\square$ Je sais que les acides aminés naturels (protéines) sont de configuration \textsc{l}.
  \item $\square$ Je sais écrire l'équilibre forme moléculaire $\rightleftharpoons$ amphion.
  \item $\square$ Je sais écrire les deux équations prouvant le caractère acide \emph{et} basique de l'amphion.
  \item $\square$ Je sais placer cation, amphion et anion sur un axe de pH à l'aide des deux p$K_a$.
  \item $\square$ Je sais écrire l'équation de formation d'un dipeptide avec élimination de \ce{H2O}.
  \item $\square$ Je sais repérer la liaison peptidique \ce{-CO-NH-} et rappeler qu'elle est plane.
  \item $\square$ Je sais nommer un peptide depuis l'extrémité N-terminale (Gly-Ala $\neq$ Ala-Gly).
  \item $\square$ Je sais dénombrer les dipeptides possibles sans protection ($n^2$).
  \item $\square$ Je sais décrire les 4 étapes de la synthèse sélective : blocage, activation, condensation, déblocage.
\end{itemize}
\end{aretenir}

\findecours

\end{document}
